% Nombres complexes : forme algébrique — Mathématiques Tle D — Cours signé OnBuch+
% Programme officiel MINESEC. Compiler avec : tectonic M02-nombres-complexes-forme-algebrique.tex
\newcommand{\DOCMATIERE}{Mathématiques}
\newcommand{\DOCNIVEAU}{Tle D}
\newcommand{\DOCMODULE}{Relations et opérations fondamentales dans l'ensemble des nombres réels}
\newcommand{\DOCLECON}{02}
\newcommand{\DOCTITRE}{Nombres complexes : forme algébrique}
% OnBuch+ — préambule « cours signés » ENRICHI (compilé avec tectonic / XeLaTeX)
% Système visuel « Pop » d'OnBuch+ : fond crème (jamais blanc), bordures encre
% épaisses, ombres « dures » décalées, titres Archivo Black en capitales.
% Version MATHÉMATIQUES : graphes (pgfplots/tikz), boîtes pédagogiques
% (définition, propriété, méthode, attention, le savais-tu, exercice résolu,
% exercice, TP, bilan), page de garde et signature OnBuch+ ancrée en bas.
%
% Macros attendues dans main.tex AVANT \input{preamble} :
%   \DOCMATIERE  \DOCNIVEAU  \DOCMODULE  \DOCLECON (numéro)  \DOCTITRE
\documentclass[11pt]{article}

\usepackage{fontspec}
\usepackage[french]{babel}
\usepackage[a4paper,margin=2cm,top=3.4cm,bottom=2.8cm,headheight=1.4cm,footskip=1.3cm]{geometry}
\usepackage{amsmath,amssymb,amsfonts}
\usepackage{mathtools} % \xmapsto, \coloneqq... souvent utilisés par les modèles
\usepackage{cancel} % simplifications barrées
% \abs{z} (module, valeur absolue) et \norm{u} (norme), utilisés par les modèles
\providecommand{\abs}{}\let\abs\relax\DeclarePairedDelimiter{\abs}{\lvert}{\rvert}
\providecommand{\norm}{}\let\norm\relax\DeclarePairedDelimiter{\norm}{\lVert}{\rVert}
\usepackage[table]{xcolor} % [table] : fournit aussi \rowcolors (alternance de lignes)
\usepackage{pagecolor}
\usepackage{tikz}
\usetikzlibrary{arrows.meta,positioning,calc,shapes.geometric,shapes.misc,shapes.arrows,decorations.pathmorphing,decorations.markings,patterns,shadows,mindmap,trees,fit,backgrounds,matrix,intersections}
\usepackage{pgfplots}
\usepackage{pgf-pie} % diagrammes circulaires \pie (statistiques)
\pgfplotsset{compat=1.18}
\usepgfplotslibrary{fillbetween,groupplots} % aires entre courbes (calcul intégral)
\usepackage{enumitem}
\usepackage{fancyhdr}
% [most] charge toutes les bibliothèques tcolorbox (listings, minted, raster,
% documentation...) et gonfle inutilement la mémoire TeX sur un document long
% (~300 boîtes) : on ne charge que ce qui est réellement utilisé.
\usepackage[skins,breakable]{tcolorbox}
\usepackage{graphicx}
\usepackage{colortbl}
\usepackage{booktabs}
\usepackage{tabularx}
\usepackage{diagbox} % case d'en-tête diagonale des tableaux à double entrée
% Colonnes extensibles centrée / gauche / droite (C, L, R), souvent utilisées par les modèles
\newcolumntype{C}{>{\centering\arraybackslash}X}
\newcolumntype{L}{>{\raggedright\arraybackslash}X}
\newcolumntype{R}{>{\raggedleft\arraybackslash}X}
\usepackage{array}
\usepackage{multicol}
\usepackage{siunitx}
% parse-numbers=false : accepte aussi \SI{2,5 \times 10^{-3}}{...}
\sisetup{locale=FR,output-decimal-marker={,},group-minimum-digits=4,parse-numbers=false}
\DeclareSIUnit\ohm{\ensuremath{\symup{\Omega}}} % U+2126 absent de la police
\usepackage{qrcode}
% \degree : commande fréquemment utilisée par les modèles pour un angle ou une
% température en dehors de \SI{}{} (non fournie par les paquets chargés).
\providecommand{\degree}{^{\circ}}
% \qed (fin de démonstration), utilisé par les modèles sans amsthm
\providecommand{\qed}{\hfill$\square$}
% \barre : double barre des valeurs interdites dans un tableau de variations (tkz-tab)
\providecommand{\barre}{\ensuremath{\|}}
% environnement proof (amsthm non chargé) utilisé par les modèles
\ifdefined\proof\else\newenvironment{proof}[1][Démonstration]{\par\noindent\textit{#1.}\ }{\qed\par}\fi
\usepackage{lastpage}
\usepackage{hyperref}
\hypersetup{hidelinks}

% ============================================================
% Palette « Pop » OnBuch+ — mêmes hex que la classe P (plus_ui.dart)
% ============================================================
\definecolor{poporange}{HTML}{F59321}
\definecolor{popdark}{HTML}{E07A0C}
\definecolor{popcream}{HTML}{FFF6E8}
\definecolor{popink}{HTML}{1C1714}
\definecolor{poppurple}{HTML}{7A5AE0}
\definecolor{popgreen}{HTML}{2E9E6A}
\definecolor{poppink}{HTML}{E0457B}
\definecolor{popblue}{HTML}{2F7DE1}
\definecolor{popgold}{HTML}{C98A12}
\definecolor{popmuted}{HTML}{8A7F76}
\definecolor{popline}{HTML}{EADFD2}
\definecolor{popgray}{HTML}{8A7F76}
\colorlet{popgrey}{popgray}
% Alias tolérants (le modèle nomme parfois les couleurs différemment)
\colorlet{popwhite}{white}
\colorlet{popblack}{popink}
\colorlet{popred}{poppink}
\colorlet{popyellow}{popgold}
% Teintes claires pour fonds de tableaux / zones
\colorlet{popblueL}{popblue!10!white}
\colorlet{poporangeL}{poporange!14!white}
\colorlet{popgreenL}{popgreen!12!white}
\colorlet{poppurpleL}{poppurple!12!white}
\colorlet{poppinkL}{poppink!12!white}
\colorlet{popgoldL}{popgold!14!white}

\pagecolor{popcream}

% ============================================================
% Typographie
% ============================================================
\defaultfontfeatures{Ligatures=TeX}
\setmainfont{PlusJakartaSans-Medium.ttf}[Path=../../../fonts/,BoldFont=PlusJakartaSans-Bold.ttf]
% Maths en sans-serif (Fira Math) avec chiffres et lettres droites de Plus
% Jakarta, pour que les formules se fondent dans le texte.
\usepackage[math-style=ISO,bold-style=ISO]{unicode-math}
\setmathfont{FiraMath-Regular.otf}[Path=../../../fonts/]
\setmathfont{PlusJakartaSans-Medium.ttf}[Path=../../../fonts/,range={up/num,up/{Latin,latin},\mathup}]
\usepackage{icomma} % virgule décimale sans espace en mode math
% Caractères mathématiques Unicode absents des polices de texte (Plus Jakarta,
% Archivo Black) : rendus par leur équivalent mathématique, en texte comme en
% formule — sinon ils s'affichent en carré vide (ex. « Arithmétique dans ℤ »).
\usepackage{newunicodechar}
\newunicodechar{ℝ}{\ensuremath{\mathbb{R}}}
\newunicodechar{ℤ}{\ensuremath{\mathbb{Z}}}
\newunicodechar{ℂ}{\ensuremath{\mathbb{C}}}
\newunicodechar{ℕ}{\ensuremath{\mathbb{N}}}
\newunicodechar{ℚ}{\ensuremath{\mathbb{Q}}}
\newunicodechar{𝔻}{\ensuremath{\mathbb{D}}}
\newunicodechar{α}{\ensuremath{\alpha}}
\newunicodechar{β}{\ensuremath{\beta}}
\newunicodechar{γ}{\ensuremath{\gamma}}
\newunicodechar{δ}{\ensuremath{\delta}}
\newunicodechar{ε}{\ensuremath{\varepsilon}}
\newunicodechar{θ}{\ensuremath{\theta}}
\newunicodechar{λ}{\ensuremath{\lambda}}
\newunicodechar{μ}{\ensuremath{\mu}}
\newunicodechar{σ}{\ensuremath{\sigma}}
\newunicodechar{τ}{\ensuremath{\tau}}
\newunicodechar{φ}{\ensuremath{\varphi}}
\newunicodechar{ψ}{\ensuremath{\psi}}
\newunicodechar{ω}{\ensuremath{\omega}}
\newunicodechar{Σ}{\ensuremath{\Sigma}}
% majuscules produites par \MakeUppercase dans les titres (α-aminés -> Α)
\newunicodechar{Α}{\ensuremath{\alpha}}
\newunicodechar{Β}{\ensuremath{\beta}}
\newunicodechar{Γ}{\ensuremath{\Gamma}}
\newunicodechar{Δ}{\ensuremath{\Delta}}
\newunicodechar{Ε}{\ensuremath{\varepsilon}}
\newunicodechar{Θ}{\ensuremath{\Theta}}
\newunicodechar{Λ}{\ensuremath{\Lambda}}
\newunicodechar{Μ}{\ensuremath{\mu}}
\newunicodechar{Τ}{\ensuremath{\tau}}
\newunicodechar{Φ}{\ensuremath{\Phi}}
\newunicodechar{Ψ}{\ensuremath{\Psi}}
\newunicodechar{Ω}{\ensuremath{\Omega}}
\newunicodechar{↦}{\ensuremath{\mapsto}}
\newunicodechar{∈}{\ensuremath{\in}}
\newunicodechar{∉}{\ensuremath{\notin}}
\newunicodechar{∧}{\ensuremath{\wedge}}
\newunicodechar{⇔}{\ensuremath{\Leftrightarrow}}
\newunicodechar{⟺}{\ensuremath{\Longleftrightarrow}}
\newunicodechar{⇒}{\ensuremath{\Rightarrow}}
\newunicodechar{⟹}{\ensuremath{\Longrightarrow}}
\newunicodechar{∘}{\ensuremath{\circ}}
\newunicodechar{∪}{\ensuremath{\cup}}
\newunicodechar{∩}{\ensuremath{\cap}}
\newunicodechar{≡}{\ensuremath{\equiv}}
\newunicodechar{⊂}{\ensuremath{\subset}}
\newunicodechar{⊥}{\ensuremath{\perp}}
\newunicodechar{∃}{\ensuremath{\exists}}
\newunicodechar{∀}{\ensuremath{\forall}}
\newunicodechar{∛}{\ensuremath{\sqrt[3]{\,}}}
\newunicodechar{∜}{\ensuremath{\sqrt[4]{\,}}}
\newunicodechar{⃗}{\ensuremath{{}^{\to}}}
\newunicodechar{ⁿ}{\ifmmode^{n}\else\textsuperscript{\ensuremath{n}}\fi}
\newunicodechar{ˣ}{\ifmmode^{x}\else\textsuperscript{\ensuremath{x}}\fi}
\newunicodechar{ᵇ}{\ifmmode^{b}\else\textsuperscript{\ensuremath{b}}\fi}
\newunicodechar{ʳ}{\ifmmode^{r}\else\textsuperscript{\ensuremath{r}}\fi}
\newunicodechar{ᵘ}{\ifmmode^{u}\else\textsuperscript{\ensuremath{u}}\fi}
\newunicodechar{ʸ}{\ifmmode^{y}\else\textsuperscript{\ensuremath{y}}\fi}
\newunicodechar{ᵃ}{\ifmmode^{a}\else\textsuperscript{\ensuremath{a}}\fi}
\newunicodechar{ᵉ}{\ifmmode^{e}\else\textsuperscript{\ensuremath{e}}\fi}
\newunicodechar{ᵏ}{\ifmmode^{k}\else\textsuperscript{\ensuremath{k}}\fi}
\newunicodechar{ᵖ}{\ifmmode^{p}\else\textsuperscript{\ensuremath{p}}\fi}
\newunicodechar{ᴺ}{\ifmmode^{N}\else\textsuperscript{\ensuremath{N}}\fi}
\newunicodechar{ᶜ}{\ifmmode^{c}\else\textsuperscript{\ensuremath{c}}\fi}
\newunicodechar{ʰ}{\ifmmode^{h}\else\textsuperscript{\ensuremath{h}}\fi}
\newunicodechar{ᵠ}{\ifmmode^{q}\else\textsuperscript{\ensuremath{q}}\fi}
\newunicodechar{ₙ}{\ifmmode_{n}\else\textsubscript{\ensuremath{n}}\fi}
\newunicodechar{ₐ}{\ifmmode_{a}\else\textsubscript{\ensuremath{a}}\fi}
\newunicodechar{ᵢ}{\ifmmode_{i}\else\textsubscript{\ensuremath{i}}\fi}
\newunicodechar{ᵣ}{\ifmmode_{r}\else\textsubscript{\ensuremath{r}}\fi}
\newunicodechar{ₖ}{\ifmmode_{k}\else\textsubscript{\ensuremath{k}}\fi}
\newunicodechar{ᵦ}{\ifmmode_{\beta}\else\textsubscript{\ensuremath{\beta}}\fi}
\newunicodechar{ₓ}{\ifmmode_{x}\else\textsubscript{\ensuremath{x}}\fi}
\newfontfamily\popdisplay{ArchivoBlack-Regular.ttf}[Path=../../../fonts/]
\newfontfamily\poplogo{SpaceGrotesk-Bold.ttf}[Path=../../../fonts/]
\newfontfamily\popsemi{PlusJakartaSans-SemiBold.ttf}[Path=../../../fonts/]
\newfontfamily\popxbold{PlusJakartaSans-ExtraBold.ttf}[Path=../../../fonts/]
\newfontfamily\popxboldls{PlusJakartaSans-ExtraBold.ttf}[Path=../../../fonts/,LetterSpace=3.0]

\setlist[enumerate]{leftmargin=1.7em,itemsep=5pt,topsep=4pt}
\setlist[itemize]{leftmargin=1.7em,itemsep=4pt,topsep=4pt,label={\color{poporange}\textbullet}}
\setlength{\parindent}{0pt}
\setlength{\parskip}{5pt}
\renewcommand{\arraystretch}{1.25}

% pgfplots : style Pop par défaut
\pgfplotsset{
  every axis/.append style={axis line style={popink,line width=1pt},tick style={popink},
    grid style={popline,line width=0.5pt},label style={font=\small},tick label style={font=\footnotesize}},
  popcurve/.style={poporange,line width=1.8pt,no markers,smooth},
}
% Même style disponible dans un \draw TikZ simple (hors pgfplots)
\tikzset{popcurve/.style={poporange,line width=1.8pt}}

% ============================================================
% Pill / squiggle
% ============================================================
\newcommand{\poppill}[3]{%
  \tikz[baseline=-0.55ex]\node[draw=popink,line width=1.2pt,rounded corners=2.6pt,%
    fill=#2,inner xsep=6.5pt,inner ysep=3pt]{%
    \popxboldls\fontsize{7}{7}\selectfont\color{#3}\MakeUppercase{#1}};%
}
\newcommand{\popsquiggle}[1]{%
  \tikz[baseline=-0.4ex]{
    \draw[#1,line width=2.6pt,line cap=round]
      plot [smooth] coordinates {(0,0.09) (0.34,0.01) (0.68,0.21) (1.02,0.01) (1.36,0.10)};
  }%
}
% Mot-clé surligné (marqueur orange sous le texte)
\newcommand{\cle}[1]{{\popxbold\color{popdark}#1}}

% ============================================================
% En-tête / pied
% ============================================================
\pagestyle{fancy}
\fancyhf{}
\renewcommand{\headrulewidth}{0pt}
\renewcommand{\footrulewidth}{0pt}
\lhead{\raisebox{-0.15ex}{\poplogo\fontsize{13}{13}\selectfont\color{popink}OnBuch\color{poporange}+}}
\rhead{\poppill{\DOCMATIERE\ \textbullet\ \DOCNIVEAU\ \textbullet\ Leçon \DOCLECON}{white}{popink}}
\lfoot{\popxbold\fontsize{7.6}{7.6}\selectfont\color{popmuted}\MakeUppercase{Cours signé OnBuch+}\ \textbullet\ {\popsemi\DOCTITRE}}
\rfoot{\tikz[baseline=-0.6ex]\node[rounded corners=8.5pt,fill=popink,minimum height=17pt,minimum width=17pt,inner xsep=5pt,inner sep=0]{\popxbold\fontsize{8}{8}\selectfont\color{popcream}\thepage\,/\,\pageref*{LastPage}};}

% ============================================================
% Titres
% ============================================================
\newcommand{\chaptitre}[2]{%
  \begin{tcolorbox}[enhanced,colback=poporange,colframe=popink,boxrule=2.6pt,arc=11pt,
      drop shadow=popink,left=15pt,right=15pt,top=12pt,bottom=11pt,before skip=3pt,after skip=18pt]
    \poppill{#2}{popink}{popcream}\par\vspace{9pt}
    {\raggedright\hyphenpenalty=10000\exhyphenpenalty=10000\popdisplay\fontsize{22}{24}\selectfont\color{popcream}\MakeUppercase{#1}\par}
  \end{tcolorbox}
}
% Section numérotée : pastille orange avec le numéro + titre Archivo
\newcounter{popsec}
\newcommand{\coursec}[1]{%
  \stepcounter{popsec}\setcounter{popsub}{0}%
  \par\vspace{16pt}\noindent
  \begin{minipage}{\linewidth}
  \tikz[baseline=-0.7ex]\node[draw=popink,line width=1.6pt,fill=poporange,rounded corners=4pt,minimum size=22pt,inner sep=0]{\popdisplay\fontsize{12}{12}\selectfont\color{popcream}\thepopsec};%
  \hspace{8pt}{\popdisplay\fontsize{15}{17}\selectfont\color{popink}\MakeUppercase{#1}}\par
  \vspace{4pt}\noindent{\color{poporange}\rule{36pt}{3.6pt}}
  \end{minipage}\par\nopagebreak\vspace{8pt}
}
\newcounter{popsub}
\newcommand{\courssub}[1]{%
  \stepcounter{popsub}%
  \par\vspace{9pt}\noindent{\popxbold\fontsize{11.5}{13}\selectfont\color{popink}{\color{poporange}\thepopsec.\thepopsub}\hspace{6pt}#1}\par\nopagebreak\vspace{4pt}
}

% ============================================================
% Boîtes pédagogiques — même recette : fond blanc, bordure épaisse colorée,
% ombre dure encre, badge plein. Titre optionnel [#1].
% ============================================================
\tcbset{popbase/.style={breakable,enhanced,colback=white,boxrule=2.3pt,arc=9pt,
  shadow={3pt}{-3pt}{0pt}{popink},left=14pt,right=14pt,top=11pt,bottom=11pt,before skip=11pt,after skip=15pt,
  fontupper=\popsemi\fontsize{10.6}{14.5}\selectfont\color{popink},
  fontlower=\popsemi\fontsize{10.6}{14.5}\selectfont\color{popink}}}
% Coche dessinée (le glyphe ✓ n'existe pas dans les polices du document)
\newcommand{\popcheck}[1]{\tikz[baseline=-0.5ex]\draw[#1,line width=1.6pt,line cap=round,line join=round](-0.13,0.0)--(-0.04,-0.09)--(0.13,0.1);}
\newcommand{\popboxhead}[3]{\poppill{#1}{#2}{white}\ifx\relax#3\relax\else\hspace{6pt}{\popxbold\fontsize{10.6}{12}\selectfont\color{popink}#3}\fi\par\vspace{7pt}}

\newtcolorbox{aretenir}[1][]{popbase,colframe=popgreen,before upper={\popboxhead{À retenir}{popgreen}{#1}}}
\newtcolorbox{exemplebox}[1][]{popbase,colframe=poporange,before upper={\popboxhead{Exemple}{poporange}{#1}}}
\newtcolorbox{definition}[1][]{popbase,colframe=popblue,before upper={\popboxhead{Définition}{popblue}{#1}}}
\newtcolorbox{propriete}[1][]{popbase,colframe=poppurple,before upper={\popboxhead{Propriété}{poppurple}{#1}}}
\newtcolorbox{methode}[1][]{popbase,colframe=popgold,before upper={\popboxhead{Méthode}{popgold}{#1}}}
\newtcolorbox{attention}[1][]{popbase,colframe=poppink,before upper={\popboxhead{Attention}{poppink}{#1}}}
\newtcolorbox{experience}[1][]{popbase,colframe=popblue,colback=popblueL,before upper={\popboxhead{Expérience}{popblue}{#1}}}
% « Le savais-tu ? » : carte violette pleine (contexte, culture, Cameroun)
\newtcolorbox{savaistu}[1][]{popbase,colback=poppurple,colframe=popink,
  fontupper=\popsemi\fontsize{10.4}{14.2}\selectfont\color{white},
  before upper={\poppill{Le savais-tu\,?}{white}{poppurple}\ifx\relax#1\relax\else\hspace{6pt}{\popxbold\color{white}#1}\fi\par\vspace{7pt}}}
% Exercice résolu : énoncé en haut, \tcblower puis la solution sur fond crème
\newcounter{popexo}\newcounter{popexr}
\newtcolorbox{exoresolu}[1][]{popbase,colframe=poporange,colbacklower=popcream,
  segmentation style={popink,line width=1.2pt,dashed},
  before upper={\stepcounter{popexr}\popboxhead{Exercice résolu \thepopexr}{poporange}{#1}},
  before lower={\poppill{Solution}{popink}{popcream}\par\vspace{6pt}}}
% Exercice (non résolu en place) — la correction est dans la partie Corrigés
\newtcolorbox{exercice}[1][]{popbase,colframe=popink,
  before upper={\stepcounter{popexo}\popboxhead{Exercice \thepopexo}{popink}{#1}}}
% Corrigé
\newtcolorbox{corrige}[1][]{popbase,colframe=popgreen,colback=popgreenL,
  before upper={\popboxhead{Corrigé}{popgreen}{#1}}}
% Objectifs / prérequis / bilan
\newtcolorbox{objectifs}{popbase,colframe=popink,colback=poporangeL,
  before upper={\popboxhead{Objectifs de la leçon}{popink}{}}}
\newtcolorbox{prerequis}{popbase,colframe=popink,colback=popblueL,
  before upper={\popboxhead{Prérequis}{popblue}{}}}
\newtcolorbox{bilan}{popbase,colframe=popink,colback=white,boxrule=2.8pt,
  before upper={\popboxhead{Fiche bilan}{poporange}{L'essentiel à connaître pour l'examen}}}
% Figure encadrée avec légende
\newtcolorbox{popfigure}[1][]{enhanced,breakable=false,colback=white,colframe=popline,boxrule=1.4pt,arc=8pt,
  left=10pt,right=10pt,top=10pt,bottom=8pt,before skip=10pt,after skip=12pt,halign=center,
  lower separated=false,halign lower=center,
  fontlower=\popsemi\fontsize{9}{11.5}\selectfont\color{popmuted},
  #1}
\newcommand{\legende}[1]{\par\vspace{6pt}{\popsemi\fontsize{9}{11.5}\selectfont\color{popmuted}#1}}

% ============================================================
% Signature OnBuch+
% ============================================================
\newcommand{\poplogoinline}[1]{{\poplogo\fontsize{#1}{#1}\selectfont\color{popink}OnBuch\color{poporange}+}}
\newcommand{\signatureonbuch}{%
  \begin{tcolorbox}[enhanced,colback=popink,colframe=popink,boxrule=2.2pt,arc=10pt,
      drop shadow=poporange,left=14pt,right=14pt,top=10pt,bottom=10pt,before skip=0pt,after skip=0pt]
    \tikz[baseline=-0.7ex]\node[circle,fill=popgreen,draw=popcream,line width=1.3pt,minimum size=22pt,inner sep=0]{\popcheck{white}};%
    \hspace{9pt}\begin{minipage}[c]{0.62\linewidth}
      {\popxbold\fontsize{12}{13}\selectfont\color{popcream}Signé\ \textbullet\ {\poplogo OnBuch\color{poporange}+}}\par\vspace{2pt}
      {\raggedright\popsemi\fontsize{8}{10.5}\selectfont\color{popline}\MakeUppercase{Équipe pédagogique OnBuch+}\\\MakeUppercase{Conforme au programme officiel MINESEC}\par}
    \end{minipage}\hfill
    {\poplogo\fontsize{22}{22}\selectfont\color{popcream}OnBuch\color{poporange}+}
  \end{tcolorbox}%
}

% ============================================================
% Page de garde
% #1 titre · #2 matière · #3 niveau · #4 module · #5 n° leçon · #6 durée ·
% #7 description / objectifs (texte ou itemize)
% ============================================================
\newcommand{\pagegarde}[7]{%
  \thispagestyle{empty}
  \newgeometry{a4paper,margin=1.7cm,top=1.6cm,bottom=1.4cm}%
  \begin{tikzpicture}[remember picture,overlay]
    \fill[poporange!18] ([xshift=-3.2cm,yshift=-3.6cm]current page.north east) circle (4.6cm);
    \fill[poppurple!14] ([xshift=2.4cm,yshift=7.6cm]current page.south west) circle (3.8cm);
  \end{tikzpicture}%
  {\poplogo\fontsize{30}{30}\selectfont\color{popink}OnBuch\color{poporange}+}\hfill
  \raisebox{0.6ex}{\poppill{Cours signé \textbullet\ Premium}{popink}{popcream}}\par
  \vspace{1pt}\popsquiggle{poppurple}\par
  \vspace{8pt}
  \begin{tcolorbox}[enhanced,colback=poporange,colframe=popink,boxrule=2.8pt,arc=13pt,
      drop shadow=popink,left=18pt,right=18pt,top=12pt,bottom=13pt,after skip=10pt]
    \poppill{#2 \textbullet\ #3}{popink}{popcream}\hspace{5pt}\poppill{#4}{white}{popink}\par\vspace{8pt}
    {\popdisplay\fontsize{14}{14}\selectfont\color{popink}LEÇON #5}\par\vspace{4pt}
    {\raggedright\hyphenpenalty=10000\exhyphenpenalty=10000\popdisplay\fontsize{25}{27}\selectfont\color{popcream}\MakeUppercase{#1}\par}
    \vspace{8pt}
    {\popxbold\fontsize{9.5}{9.5}\selectfont\color{popink}\MakeUppercase{Durée conseillée : #6}}
  \end{tcolorbox}
  \begin{tcolorbox}[enhanced,colback=white,colframe=popink,boxrule=2.2pt,arc=10pt,
      drop shadow=popink,left=16pt,right=16pt,top=10pt,bottom=10pt,after skip=8pt]
    \poppill{Dans cette leçon}{poppurple}{white}\par\vspace{5pt}
    \popsemi\fontsize{9.3}{12.6}\selectfont\color{popink}#7
  \end{tcolorbox}
  \noindent\begin{minipage}[t]{0.487\linewidth}
    \begin{tcolorbox}[enhanced,colback=popgreen,colframe=popink,boxrule=2.2pt,arc=10pt,
        drop shadow=popink,left=11pt,right=11pt,top=10pt,bottom=10pt,height=3.1cm,valign=center]
      \tcbox[colback=white,colframe=popink,boxrule=1.3pt,arc=5pt,boxsep=0pt,nobeforeafter,
        left=3pt,right=3pt,top=3pt,bottom=3pt,valign=center]{\qrcode[height=1.9cm]{https://play.google.com/store/apps/details?id=com.luvvix.onbuch&hl=fr}}%
      \hspace{8pt}\begin{minipage}[c]{0.5\linewidth}
        {\popdisplay\fontsize{11}{12.5}\selectfont\color{white}\MakeUppercase{Télécharge OnBuch}}\par\vspace{4pt}
        {\raggedright\popsemi\fontsize{8.4}{11}\selectfont\color{white}Gratuit \textbullet\ Google Play\\Tuteur IA, annales, concours, résultats\par}
      \end{minipage}
    \end{tcolorbox}
  \end{minipage}\hfill
  \begin{minipage}[t]{0.487\linewidth}
    \begin{tcolorbox}[enhanced,colback=poppurple,colframe=popink,boxrule=2.2pt,arc=10pt,
        drop shadow=popink,left=11pt,right=11pt,top=10pt,bottom=10pt,height=3.1cm,valign=center]
      \tikz[baseline=-0.7ex]\node[circle,fill=white,draw=popink,line width=1.3pt,minimum size=1.6cm,inner sep=0]{\popdisplay\fontsize{24}{24}\selectfont\color{poppurple}+};%
      \hspace{8pt}\begin{minipage}[c]{0.6\linewidth}
        {\popdisplay\fontsize{11}{12.5}\selectfont\color{white}\MakeUppercase{Passe à OnBuch+}}\par\vspace{4pt}
        {\raggedright\popsemi\fontsize{8.4}{11}\selectfont\color{white}Tous les cours signés, Tuteur IA illimité, fiches PDF — onglet OnBuch+ de l'app\par}
      \end{minipage}
    \end{tcolorbox}
  \end{minipage}
  \vspace{8pt}
  \popcontenu
  \vfill
  \signatureonbuch
  \restoregeometry
  \newpage
}

% Rangée « Ce cours contient » (page de garde)
\newcommand{\poptile}[3]{% #1 couleur #2 gros texte #3 légende
  \begin{tcolorbox}[enhanced,colback=#1,colframe=popink,boxrule=2pt,arc=9pt,shadow={2.5pt}{-2.5pt}{0pt}{popink},
      left=6pt,right=6pt,top=9pt,bottom=9pt,halign=center,valign=center,height=1.75cm,width=\linewidth,nobeforeafter]
    {\popdisplay\fontsize{12.5}{13}\selectfont\color{white}\MakeUppercase{#2}}\par\vspace{3pt}
    {\centering\popsemi\fontsize{7.8}{9.5}\selectfont\color{white}#3\par}
  \end{tcolorbox}}
\newcommand{\popcontenu}{%
  \par\noindent{\popxboldls\fontsize{7.5}{7.5}\selectfont\color{popmuted}CE COURS SIGNÉ CONTIENT}\par\vspace{7pt}
  \noindent\begin{minipage}[t]{0.235\linewidth}\poptile{popblue}{Cours}{complet, illustré, pas à pas}\end{minipage}\hfill
  \begin{minipage}[t]{0.235\linewidth}\poptile{poporange}{Méthodes}{et exercices résolus}\end{minipage}\hfill
  \begin{minipage}[t]{0.235\linewidth}\poptile{poppink}{Exercices}{type examen + corrigés}\end{minipage}\hfill
  \begin{minipage}[t]{0.235\linewidth}\poptile{popgreen}{Fiche bilan}{l'essentiel à retenir}\end{minipage}\par}

% Sommaire
\newtcolorbox{sommaire}{enhanced,colback=white,colframe=popink,boxrule=2.2pt,arc=10pt,
  drop shadow=popink,left=18pt,right=18pt,top=13pt,bottom=13pt,before skip=6pt,after skip=20pt,
  before upper={\poppill{Sommaire}{poppurple}{white}\par\vspace{9pt}\popsemi\fontsize{10.6}{14.5}\selectfont\color{popink}}}

% Signature de fin de cours — poussée tout en bas de la dernière page
\newcommand{\findecours}{%
  \par\vfill\nopagebreak
  \begin{center}\popsquiggle{poporange}\end{center}\vspace{4pt}
  \signatureonbuch
}
\AtBeginDocument{\def\llbracket{\mathopen{[\mkern-3.5mu[}}\def\rrbracket{\mathclose{]\mkern-3.5mu]}}} % intervalles d'entiers (glyphe ⟦ absent de la police)
% Glyphes absents de Fira Math : redéfinis à partir de glyphes disponibles
\AtBeginDocument{%
  \def\setminus{\mathbin{\backslash}}\let\smallsetminus\setminus
  \def\vdots{\mathord{\vbox{\baselineskip3.2pt\lineskiplimit0pt\kern4pt\hbox{.}\hbox{.}\hbox{.}}}}%
  \def\ddots{\mathinner{\mkern1mu\raise6.5pt\hbox{.}\mkern2mu\raise3.6pt\hbox{.}\mkern2mu\raise0.7pt\hbox{.}\mkern1mu}}%
  \def\triangle{\mathord{\scalebox{0.9}{$\symup{\Delta}$}}}%
  \def\nabla{\mathord{\raisebox{\depth}{\scalebox{1}[-1]{$\symup{\Delta}$}}}}%
  \def\checkmark{\popcheck{popdark}}%
  \def\star{\mathbin{\ast}}\def\bigstar{\mathord{\ast}}%
  \def\diamond{\mathbin{\scalebox{0.6}{$\Diamond$}}}%
}

%% FIGFIX-BEGIN v3 — correctifs globaux des figures (docs/onbuch-generation/tools/figfix)
\usepackage{adjustbox}\usepackage{etoolbox}
% toute figure TikZ/pgfplots plus large que la ligne est réduite à la largeur disponible
\BeforeBeginEnvironment{tikzpicture}{\begin{adjustbox}{max width=\linewidth}}
\AfterEndEnvironment{tikzpicture}{\end{adjustbox}}
% légendes pgfplots lisibles par-dessus les courbes
\pgfplotsset{every axis/.append style={legend style={fill=white,fill opacity=0.92,text opacity=1,draw=black!15,inner sep=3pt}}}
% étiquettes d'arêtes (syntaxe edge["..."]) avec fond blanc
\tikzset{every edge quotes/.append style={fill=white,inner sep=1.2pt}}
% bulles de cartes mentales : texte à la ligne dans la bulle, bulle un peu plus grande
\tikzset{every concept/.append style={text width=2.0cm,align=center,minimum size=2.3cm,inner sep=2pt}}
%% FIGFIX-END
\begin{document}

\pagegarde{Nombres complexes : forme algébrique}{Mathématiques}{Tle D}{Relations et opérations fondamentales dans l'ensemble des nombres réels}{02}{9 h}{Cette leçon introduit l'ensemble ℂ des nombres complexes et leur forme algébrique a + ib. On y apprend à effectuer les opérations fondamentales, à utiliser le conjugué et le module, à extraire les racines carrées d'un complexe et à résoudre dans ℂ des équations polynomiales de degrés 2 et 3.
\vspace{4pt}
\begin{multicols}{2}\small
\begin{itemize}
\item À la fin de cette leçon, je sais définir l'ensemble ℂ, l'unité imaginaire i et la forme algébrique d'un nombre complexe.
\item À la fin de cette leçon, je sais effectuer la somme, le produit, l'opposé, l'inverse et le quotient de deux nombres complexes sous forme algébrique.
\item À la fin de cette leçon, je sais calculer le conjugué d'un nombre complexe et utiliser ses propriétés (notamment z·z̄ = |z|²).
\item À la fin de cette leçon, je sais calculer le module d'un nombre complexe et appliquer ses propriétés opératoires.
\item À la fin de cette leçon, je sais déterminer les racines carrées d'un nombre complexe donné sous forme algébrique.
\item À la fin de cette leçon, je sais résoudre dans ℂ une équation du second degré à coefficients réels, y compris lorsque le discriminant est négatif.
\end{itemize}\end{multicols}}

\begin{sommaire}
\begin{multicols}{2}
\begin{enumerate}
\item L'ensemble ℂ et la forme algébrique
\item Opérations dans ℂ
\item Conjugué d'un nombre complexe
\item Module d'un nombre complexe
\item Racines carrées d'un nombre complexe
\item Équations du second degré et polynômes de degré 3 dans ℂ
\item Activité d'intégration
\item Méthodes et exercices résolus
\item Exercices
\item Corrigés des exercices
\item Fiche bilan
\end{enumerate}
\end{multicols}
\end{sommaire}

\begin{objectifs}
\begin{itemize}
\item À la fin de cette leçon, je sais définir l'ensemble ℂ, l'unité imaginaire i et la forme algébrique d'un nombre complexe.
\item À la fin de cette leçon, je sais effectuer la somme, le produit, l'opposé, l'inverse et le quotient de deux nombres complexes sous forme algébrique.
\item À la fin de cette leçon, je sais calculer le conjugué d'un nombre complexe et utiliser ses propriétés (notamment z·z̄ = |z|²).
\item À la fin de cette leçon, je sais calculer le module d'un nombre complexe et appliquer ses propriétés opératoires.
\item À la fin de cette leçon, je sais déterminer les racines carrées d'un nombre complexe donné sous forme algébrique.
\item À la fin de cette leçon, je sais résoudre dans ℂ une équation du second degré à coefficients réels, y compris lorsque le discriminant est négatif.
\item À la fin de cette leçon, je sais factoriser un polynôme de degré 3 à coefficients complexes connaissant une racine, et résoudre l'équation associée.
\item À la fin de cette leçon, je sais interpréter géométriquement un complexe, son conjugué et son module dans le plan muni d'un repère orthonormé (complément utile au Bac).
\end{itemize}
\end{objectifs}

\begin{prerequis}
\begin{itemize}
\item Opérations dans ℝ : développement, factorisation, identités remarquables
\item Résolution dans ℝ de l'équation du second degré (discriminant Δ, cas Δ ≥ 0)
\item Racine carrée d'un réel positif et calculs sur les radicaux
\item Repère orthonormé du plan, coordonnées d'un point et distance entre deux points
\item Division euclidienne de polynômes et méthode des coefficients indéterminés (Première)
\end{itemize}
\end{prerequis}

\newpage

\coursec{L'ensemble $\mathbb{C}$ et la forme algébrique}

Au lycée de Ngaoundéré, un élève technicien en électricité doit calculer l'impédance d'un circuit alimenté par le groupe électrogène du quartier. En manipulant les grandeurs alternatives (tension, courant), il découvre que les formules font intervenir un nombre étrange dont le carré vaut $-1$ ! De son côté, le mathématicien se heurte depuis longtemps à un mur : l'équation $x^2 + 1 = 0$, pourtant si simple, n'a \emph{aucune} solution dans $\mathbb{R}$, car le carré d'un nombre réel est toujours positif ou nul. Cette impossibilité bloque la résolution de nombreux problèmes d'ingénierie, de physique et d'informatique.

Plutôt que de renoncer, les mathématiciens ont fait ce qu'ils ont toujours fait : ils ont \emph{inventé} un ensemble plus vaste. De $\mathbb{N}$ à $\mathbb{Z}$ pour soustraire, de $\mathbb{Z}$ à $\mathbb{Q}$ pour diviser, de $\mathbb{Q}$ à $\mathbb{R}$ pour mesurer les longueurs... et maintenant de $\mathbb{R}$ à $\mathbb{C}$ pour extraire les racines des nombres négatifs. Comment construire cet ensemble $\mathbb{C}$ ? Comment y calculer ? Comment y résoudre des équations ? C'est l'objet de cette leçon.

\courssub{Pourquoi élargir $\mathbb{R}$ ? Une nécessité historique et pratique}

\begin{experience}[Une équation sans solution]
Cherchons à résoudre dans $\mathbb{R}$ l'équation $x^2 + 1 = 0$.
\begin{enumerate}
\item Cette équation s'écrit $x^2 = -1$.
\item Or, pour tout réel $x$, on a $x^2 \geq 0$ (produit de deux nombres de même signe).
\item Donc $x^2$ ne peut jamais être égal à $-1$ : l'équation n'a \textbf{aucune solution réelle}.
\end{enumerate}
Conclusion : si l'on veut que cette équation ait des solutions, il faut \emph{créer} un nouveau nombre, noté $\mathrm{i}$, dont le carré vaut $-1$, et un ensemble qui le contient.
\end{experience}

Cette démarche n'a rien d'arbitraire. Elle suit exactement le même schéma que les extensions de nombres que tu connais depuis l'enfance : impossible de calculer $3 - 5$ dans $\mathbb{N}$ ? On invente $\mathbb{Z}$. Impossible de diviser $1$ par $3$ dans $\mathbb{Z}$ ? On invente $\mathbb{Q}$. Impossible de mesurer la diagonale d'un carré de côté $1$ dans $\mathbb{Q}$ ? On invente $\mathbb{R}$. Impossible de résoudre $x^2 = -1$ dans $\mathbb{R}$ ? On invente $\mathbb{C}$.

\begin{savaistu}[Les nombres « imaginaires » ont mis trois siècles à être acceptés]
Le nombre $\mathrm{i}$ est apparu au XVI\textsuperscript{e} siècle dans les travaux des algébristes italiens \textbf{Jérôme Cardan} (1501--1576) et \textbf{Rafael Bombelli} (1526--1572). En résolvant des équations du troisième degré, ils ont remarqué que des racines carrées de nombres négatifs apparaissaient dans les calculs intermédiaires, puis disparaissaient pour donner des solutions bien réelles ! Descartes les qualifiera méprisamment d'« imaginaires » en 1637. Il faudra attendre Euler, Gauss et Argand (fin XVIII\textsuperscript{e} -- début XIX\textsuperscript{e} siècle) pour que ces nombres soient pleinement acceptés. Aujourd'hui, ils sont indispensables : électricité, traitement du signal (ton téléphone MTN ou Orange utilise les nombres complexes pour moduler les ondes !), mécanique quantique, aéronautique.
\end{savaistu}

\courssub{Le théorème d'existence et l'unité imaginaire}

On admet le résultat fondateur suivant, qui garantit que notre construction est cohérente.

\begin{propriete}[Théorème d'existence (admis)]
Il existe un ensemble, noté $\mathbb{C}$, appelé \cle{ensemble des nombres complexes}, qui vérifie les propriétés suivantes :
\begin{enumerate}
\item $\mathbb{C}$ contient $\mathbb{R}$ : tout nombre réel est un nombre complexe ;
\item $\mathbb{C}$ est muni d'une \cle{addition} et d'une \cle{multiplication} qui prolongent celles de $\mathbb{R}$ et qui en conservent les règles de calcul (commutativité, associativité, distributivité, existence d'opposés et d'inverses) ;
\item $\mathbb{C}$ contient un élément noté $\mathrm{i}$, appelé \cle{unité imaginaire}, tel que
\[ \mathrm{i}^2 = -1 \; ; \]
\item tout élément $z$ de $\mathbb{C}$ s'écrit de manière \textbf{unique} sous la forme
\[ z = a + \mathrm{i}\,b \quad \text{avec } (a, b) \in \mathbb{R}^2. \]
\end{enumerate}
\end{propriete}

\begin{attention}[Le nombre $\mathrm{i}$ n'est pas un réel !]
On écrit $\mathrm{i}^2 = -1$, mais il serait faux d'écrire $\mathrm{i} = \sqrt{-1}$ au sens de la racine carrée des réels : la fonction $\sqrt{\cdot}$ n'est définie que sur $[0\,; +\infty[$. Le symbole $\mathrm{i}$ est un \emph{nouvel objet}, dont la seule propriété de départ est que son carré vaut $-1$. Toutes les règles de calcul habituelles (sur les sommes et les produits) restent valables, mais attention : la formule $\sqrt{a}\times\sqrt{b} = \sqrt{ab}$ ne s'étend \textbf{pas} aux nombres négatifs. Par exemple, $\mathrm{i}^2 = \mathrm{i} \times \mathrm{i} = -1$, alors qu'une application naïve de la règle des racines donnerait $\sqrt{(-1)(-1)} = \sqrt{1} = 1$ : contradiction !
\end{attention}

\courssub{Forme algébrique, partie réelle, partie imaginaire}

\begin{definition}[Forme algébrique, parties réelle et imaginaire]
Soit $z$ un nombre complexe. L'écriture unique
\[ z = a + \mathrm{i}\,b \quad \text{avec } (a,b) \in \mathbb{R}^2 \]
est appelée \cle{forme algébrique} de $z$.
\begin{itemize}
\item Le réel $a$ est appelé \cle{partie réelle} de $z$, notée $\operatorname{Re}(z)$ ;
\item le réel $b$ est appelé \cle{partie imaginaire} de $z$, notée $\operatorname{Im}(z)$.
\end{itemize}
On a donc : $z = \operatorname{Re}(z) + \mathrm{i}\,\operatorname{Im}(z)$.
\end{definition}

\begin{definition}[Imaginaires purs et inclusion de $\mathbb{R}$ dans $\mathbb{C}$]
\begin{itemize}
\item Si $\operatorname{Im}(z) = 0$, alors $z = a$ est un \textbf{nombre réel}. Ainsi $\mathbb{R} \subset \mathbb{C}$ : tout réel est un complexe de partie imaginaire nulle.
\item Si $\operatorname{Re}(z) = 0$, alors $z = \mathrm{i}\,b$ est appelé \cle{imaginaire pur}. L'ensemble des imaginaires purs est noté $\mathrm{i}\mathbb{R}$.
\end{itemize}
\end{definition}

\begin{exemplebox}[Identifier parties réelle et imaginaire]
\begin{itemize}
\item $z = 3 - 5\mathrm{i}$ : $\operatorname{Re}(z) = 3$ et $\operatorname{Im}(z) = -5$ (attention au signe !) ;
\item $z = 7$ : $\operatorname{Re}(z) = 7$ et $\operatorname{Im}(z) = 0$ ; c'est un nombre réel ;
\item $z = -4\mathrm{i}$ : $\operatorname{Re}(z) = 0$ et $\operatorname{Im}(z) = -4$ ; c'est un imaginaire pur ;
\item $z = \mathrm{i}$ : $\operatorname{Re}(z) = 0$ et $\operatorname{Im}(z) = 1$.
\end{itemize}
\end{exemplebox}

\begin{attention}[La partie imaginaire est un réel !]
Erreur classique au Baccalauréat : pour $z = a + \mathrm{i}b$, la partie imaginaire est le \textbf{réel} $b$, et \emph{non} $\mathrm{i}b$. Ainsi, pour $z = 3 - 5\mathrm{i}$, on écrit $\operatorname{Im}(z) = -5$ et jamais $\operatorname{Im}(z) = -5\mathrm{i}$. Une réponse du type $\operatorname{Im}(z) = -5\mathrm{i}$ est sanctionnée.
\end{attention}

\courssub{Égalité de deux nombres complexes}

L'unicité de la forme algébrique a une conséquence capitale, utilisée dans presque tous les exercices.

\begin{propriete}[Égalité de deux complexes]
Soient $z = a + \mathrm{i}b$ et $z' = a' + \mathrm{i}b'$ deux nombres complexes, avec $a, b, a', b'$ réels. Alors :
\[ z = z' \iff \begin{cases} a = a' \\ b = b' \end{cases} \quad \text{c'est-à-dire} \quad \begin{cases} \operatorname{Re}(z) = \operatorname{Re}(z') \\ \operatorname{Im}(z) = \operatorname{Im}(z') \end{cases} \]
En particulier : $\; z = 0 \iff \operatorname{Re}(z) = 0 \;\text{et}\; \operatorname{Im}(z) = 0$.
\end{propriete}

\begin{experience}[Pourquoi l'écriture $a + \mathrm{i}b$ est-elle unique ?]
Montrons que si $a + \mathrm{i}b = a' + \mathrm{i}b'$ avec $a, b, a', b'$ réels, alors $a = a'$ et $b = b'$.
\begin{enumerate}
\item Supposons $a + \mathrm{i}b = a' + \mathrm{i}b'$. Alors $a - a' = \mathrm{i}(b' - b)$.
\item Supposons $b' - b \neq 0$. On pourrait diviser : $\mathrm{i} = \dfrac{a - a'}{b' - b}$, qui est un quotient de deux réels, donc un réel.
\item Or $\mathrm{i}$ n'est pas réel (sinon $\mathrm{i}^2 \geq 0$, ce qui contredit $\mathrm{i}^2 = -1$). Contradiction.
\item Donc $b' - b = 0$, c'est-à-dire $b = b'$, et par suite $a - a' = 0$, donc $a = a'$. $\blacksquare$
\end{enumerate}
Cette démonstration par l'absurde est formatrice : elle montre que l'unicité repose entièrement sur le fait que $\mathrm{i}$ n'est pas réel.
\end{experience}

\begin{methode}[Exploiter une égalité entre complexes]
Pour déterminer des inconnues réelles $x$ et $y$ vérifiant une égalité entre nombres complexes :
\begin{enumerate}
\item regrouper chaque membre sous la forme $a + \mathrm{i}b$ (parties réelles d'un côté, termes en $\mathrm{i}$ de l'autre) ;
\item identifier les parties réelles entre elles et les parties imaginaires entre elles ;
\item résoudre le système de deux équations à deux inconnues obtenu.
\end{enumerate}
\end{methode}

\begin{exoresolu}[Déterminer deux réels à partir d'une égalité]
\textbf{Énoncé.} Déterminer les réels $x$ et $y$ tels que : $(2x + 1) + \mathrm{i}(y - 3) = 5 + 2\mathrm{i}$.
\tcblower
\textbf{Solution.} Les deux membres sont sous forme algébrique. Par identification des parties réelles et imaginaires :
\[ \begin{cases} 2x + 1 = 5 \\ y - 3 = 2 \end{cases} \iff \begin{cases} 2x = 4 \\ y = 5 \end{cases} \iff \begin{cases} x = 2 \\ y = 5 \end{cases} \]
Vérification : $(2\times 2 + 1) + \mathrm{i}(5 - 3) = 5 + 2\mathrm{i}$. La solution est le couple $(2\,; 5)$.
\end{exoresolu}

\courssub{Interprétation géométrique : le plan complexe}

Voici l'idée qui a rendu les nombres complexes définitivement respectables : puisqu'un complexe $z = a + \mathrm{i}b$ est entièrement déterminé par le couple de réels $(a\,; b)$, on peut le \emph{représenter} par un point du plan !

\begin{definition}[Plan complexe, image, affixe]
Le plan est muni d'un repère orthonormé direct $(O\,; \vec{u}, \vec{v})$, appelé alors \cle{plan complexe}.
\begin{itemize}
\item À tout nombre complexe $z = a + \mathrm{i}b$, on associe le point $M(a\,; b)$, appelé \cle{point image} de $z$, et le vecteur $\overrightarrow{OM} = a\vec{u} + b\vec{v}$, appelé \cle{vecteur image} de $z$.
\item Réciproquement, le nombre $z = a + \mathrm{i}b$ est appelé \cle{affixe} du point $M(a\,; b)$ et du vecteur $\overrightarrow{OM}$. On note $z_M = a + \mathrm{i}b$ ou $\operatorname{aff}(M) = a + \mathrm{i}b$.
\end{itemize}
L'axe $(O\,; \vec{u})$ est l'\cle{axe des réels} (affixes réelles) ; l'axe $(O\,; \vec{v})$ est l'\cle{axe des imaginaires purs}.
\end{definition}

\begin{exemplebox}[Placer des points dans le plan complexe]
Plaçons les points images de $z_1 = 2 + 3\mathrm{i}$, $z_2 = -1 + \mathrm{i}$ et $z_3 = -2\mathrm{i}$ :
\begin{itemize}
\item $M_1(2\,; 3)$ est l'image de $z_1$ : on lit $\operatorname{Re}(z_1) = 2$ en abscisse, $\operatorname{Im}(z_1) = 3$ en ordonnée ;
\item $M_2(-1\,; 1)$ est l'image de $z_2$ ;
\item $M_3(0\,; -2)$ est l'image de $z_3$ : c'est un imaginaire pur, donc son point image est sur l'axe $(O\,; \vec{v})$.
\end{itemize}
\end{exemplebox}

\begin{popfigure}
\begin{center}
\begin{tikzpicture}[scale=1.1,>=Stealth]
\draw[popline!60] (-3.4,-2.9) grid (3.4,3.9);
\draw[->,thick,popdark] (-3.4,0) -- (3.4,0) node[below right] {axe des réels};
\draw[->,thick,popdark] (0,-2.9) -- (0,3.9) node[above left] {axe des imaginaires purs};
\draw[->,very thick,popblue] (0,0) -- (1,0) node[midway,below] {$\vec{u}$};
\draw[->,very thick,popblue] (0,0) -- (0,1) node[midway,left] {$\vec{v}$};
\node[below left,popdark] at (0,0) {$O$};
\draw[->,very thick,poporange] (0,0) -- (2,3);
\fill[poporange] (2,3) circle (2.2pt) node[above right,popdark] {$M_1$ : $z_1 = 2 + 3\mathrm{i}$};
\draw[->,very thick,popgreen] (0,0) -- (-1,1);
\fill[popgreen] (-1,1) circle (2.2pt) node[above left,popdark] {$M_2$ : $z_2 = -1 + \mathrm{i}$};
\draw[->,very thick,poppurple] (0,0) -- (0,-2);
\fill[poppurple] (0,-2) circle (2.2pt) node[below right,popdark] {$M_3$ : $z_3 = -2\mathrm{i}$};
\draw[dashed,popmuted] (2,0) -- (2,3) -- (0,3);
\draw[dashed,popmuted] (-1,0) -- (-1,1) -- (0,1);
\end{tikzpicture}
\end{center}
\legende{Plan complexe $(O\,; \vec{u}, \vec{v})$ : points images et vecteurs images de $z_1 = 2 + 3\mathrm{i}$, $z_2 = -1 + \mathrm{i}$, $z_3 = -2\mathrm{i}$.}
\end{popfigure}

\begin{aretenir}[Correspondance fondamentale]
Dans le plan complexe :
\begin{itemize}
\item chaque nombre complexe $z = a + \mathrm{i}b$ correspond à un unique point $M(a\,; b)$ et à un unique vecteur $\overrightarrow{OM}$ ;
\item $z$ réel $\iff$ $M$ est sur l'axe des abscisses ; $z$ imaginaire pur $\iff$ $M$ est sur l'axe des ordonnées ;
\item deux complexes sont égaux si et seulement si leurs points images coïncident.
\end{itemize}
Cette correspondance géométrie $\leftrightarrow$ algèbre est un outil puissant du Baccalauréat : elle permet de traduire des problèmes géométriques en calculs sur les complexes, et inversement.
\end{aretenir}

\courssub{La tour des ensembles de nombres}

L'ensemble $\mathbb{C}$ couronne la construction des ensembles de nombres étudiés au lycée. Chaque étage a été créé pour résoudre des équations impossibles à l'étage inférieur.

\begin{popfigure}
\begin{center}
\begin{tikzpicture}[scale=0.95]
\draw[fill=poppurpleL,draw=poppurple,thick] (0,0) ellipse (6.2 and 3.6);
\node[poppurple,font=\Large\bfseries] at (-4.6,2.9) {$\mathbb{C}$};
\draw[fill=popblueL,draw=popblue,thick] (-0.4,-0.2) ellipse (4.9 and 2.7);
\node[popblue,font=\large\bfseries] at (-3.6,2.0) {$\mathbb{R}$};
\draw[fill=popgreenL,draw=popgreen,thick] (-0.8,-0.4) ellipse (3.6 and 1.9);
\node[popgreen,font=\large\bfseries] at (-2.7,1.2) {$\mathbb{Q}$};
\draw[fill=poporangeL,draw=poporange,thick] (-1.1,-0.6) ellipse (2.5 and 1.25);
\node[poporange,font=\normalsize\bfseries] at (-1.9,0.35) {$\mathbb{Z}$};
\draw[fill=popgoldL,draw=popgold,thick] (-1.35,-0.75) ellipse (1.35 and 0.7);
\node[popdark,font=\normalsize\bfseries] at (-1.35,-0.75) {$\mathbb{N}$};
\node[popdark,align=left,font=\small] at (3.4,2.6) {$\mathrm{i}$, $2+3\mathrm{i}$};
\node[popdark,align=left,font=\small] at (2.6,1.35) {$\sqrt{2}$, $\pi$};
\node[popdark,align=left,font=\small] at (1.9,0.35) {$\dfrac{3}{4}$, $-0{,}5$};
\node[popdark,font=\small] at (0.6,-0.75) {$-3$};
\node[popdark,font=\small] at (-1.35,-0.15) {$0, 1, 2$};
\end{tikzpicture}
\end{center}
\legende{Inclusions des ensembles de nombres : $\mathbb{N} \subset \mathbb{Z} \subset \mathbb{Q} \subset \mathbb{R} \subset \mathbb{C}$ (l'ensemble $\mathbb{D}$ des décimaux s'intercale entre $\mathbb{Z}$ et $\mathbb{Q}$).}
\end{popfigure}

\begin{aretenir}[Ce qu'il faut retenir de cette section]
\begin{itemize}
\item $\mathbb{C}$ est un ensemble contenant $\mathbb{R}$, muni d'une addition et d'une multiplication aux règles habituelles, et d'un élément $\mathrm{i}$ tel que $\mathrm{i}^2 = -1$.
\item Tout complexe s'écrit de manière \textbf{unique} $z = a + \mathrm{i}b$ (forme algébrique) avec $a = \operatorname{Re}(z)$ et $b = \operatorname{Im}(z)$ réels.
\item $z = z' \iff \operatorname{Re}(z) = \operatorname{Re}(z')$ et $\operatorname{Im}(z) = \operatorname{Im}(z')$ ; en particulier $z = 0 \iff a = 0$ et $b = 0$.
\item $z$ réel $\iff \operatorname{Im}(z) = 0$ ; $z$ imaginaire pur $\iff \operatorname{Re}(z) = 0$.
\item Dans le plan complexe, $z = a + \mathrm{i}b$ est l'affixe du point $M(a\,; b)$ et du vecteur $\overrightarrow{OM}$.
\end{itemize}
\end{aretenir}

\begin{exercice}[Pour t'entraîner]
\begin{enumerate}
\item Donner la partie réelle et la partie imaginaire de : $z_1 = -2 + 7\mathrm{i}$ ; $z_2 = \dfrac{1 - \mathrm{i}}{2}$ ; $z_3 = 5\mathrm{i}$ ; $z_4 = -\sqrt{3}$.
\item Déterminer les réels $x$ et $y$ tels que : $(x - 2y) + \mathrm{i}(2x + y) = 4 - 3\mathrm{i}$.
\item Placer dans le plan complexe les points $A$, $B$, $C$ d'affixes respectives $z_A = 1 + 2\mathrm{i}$, $z_B = -3$, $z_C = \dfrac{3}{2}\mathrm{i}$. Que peut-on dire des points $B$ et $C$ ?
\end{enumerate}
\end{exercice}

\begin{corrige}[Exercice]
\begin{enumerate}
\item $\operatorname{Re}(z_1) = -2$, $\operatorname{Im}(z_1) = 7$ ; $z_2 = \dfrac{1}{2} - \dfrac{1}{2}\mathrm{i}$ donc $\operatorname{Re}(z_2) = \dfrac{1}{2}$, $\operatorname{Im}(z_2) = -\dfrac{1}{2}$ ; $\operatorname{Re}(z_3) = 0$, $\operatorname{Im}(z_3) = 5$ ; $\operatorname{Re}(z_4) = -\sqrt{3}$, $\operatorname{Im}(z_4) = 0$.
\item Par identification : $\begin{cases} x - 2y = 4 \\ 2x + y = -3 \end{cases}$. De la seconde équation, $y = -3 - 2x$ ; en reportant : $x - 2(-3 - 2x) = 4$, soit $5x + 6 = 4$, $x = -\dfrac{2}{5}$, puis $y = -3 + \dfrac{4}{5} = -\dfrac{11}{5}$.
\item $A(1\,; 2)$ ; $B(-3\,; 0)$ est sur l'axe des réels car $z_B$ est réel ; $C(0\,; \tfrac{3}{2})$ est sur l'axe des ordonnées car $z_C$ est imaginaire pur.
\end{enumerate}
\end{corrige}

\coursec{Opérations dans $\mathbb{C}$}

Dans la section précédente, nous avons construit l'ensemble $\mathbb{C}$ des nombres complexes et donné à chacun d'eux une \emph{carte d'identité} : la forme algébrique $z = a + ib$, avec $a$ et $b$ réels et $i^2 = -1$. Mais un ensemble de nombres ne sert à rien si l'on ne peut pas \emph{calculer} avec. Imagine un opérateur comme MTN ou Orange qui ne pourrait que stocker les numéros de ses abonnés sans jamais pouvoir additionner les durées d'appels ni multiplier un tarif par un nombre de minutes : inutile ! De la même façon, $\mathbb{C}$ ne devient un outil puissant que lorsqu'on y définit une addition et une multiplication. C'est précisément l'objet de cette section : tu vas découvrir que les règles de calcul que tu pratiques depuis le collège dans $\mathbb{R}$ s'étendent à $\mathbb{C}$, à condition de respecter une seule règle nouvelle, toute simple : \cle{$i^2 = -1$}.

\courssub{Somme de deux nombres complexes}

\textbf{Intuition.} Additionner deux complexes, c'est additionner « étage par étage » : les parties réelles entre elles, les parties imaginaires entre elles. C'est exactement comme lorsque tu additionnes deux sommes d'argent exprimées en billets de deux types différents : tu regroupes les billets de même nature.

\begin{definition}[Somme de deux nombres complexes]
Soient $z = a + ib$ et $z' = a' + ib'$ deux nombres complexes ($a, b, a', b' \in \mathbb{R}$). On définit leur \cle{somme} par :
\[ z + z' = (a + ib) + (a' + ib') = (a + a') + i(b + b'). \]
Autrement dit : $\mathrm{Re}(z + z') = \mathrm{Re}(z) + \mathrm{Re}(z')$ et $\mathrm{Im}(z + z') = \mathrm{Im}(z) + \mathrm{Im}(z')$.
\end{definition}

\begin{exemplebox}[Addition de deux complexes]
Avec $z_1 = 3 + i$ et $z_2 = 1 + 2i$ :
\[ z_1 + z_2 = (3 + 1) + i(1 + 2) = 4 + 3i. \]
Avec $z = 5 - 2i$ et $z' = -3 + 7i $ : $z + z' = (5 - 3) + i(-2 + 7) = 2 + 5i$.
\end{exemplebox}

\begin{definition}[Opposé d'un nombre complexe]
L'\cle{opposé} de $z = a + ib$ est le complexe noté $-z$ défini par :
\[ -z = -a - ib. \]
Il vérifie $z + (-z) = 0$. La \cle{différence} de deux complexes est alors définie par $z - z' = z + (-z') = (a - a') + i(b - b')$.
\end{definition}

\textbf{Interprétation géométrique : la règle du parallélogramme.} Plaçons-nous dans le plan complexe muni du repère orthonormé $(O\,; \vec{u}, \vec{v})$. Si $M$ et $M'$ sont les points d'affixes $z$ et $z'$, alors le point $S$ d'affixe $z + z'$ est le quatrième sommet du parallélogramme construit sur $O$, $M$ et $M'$ : c'est la célèbre relation de Chasles vectorielle $\overrightarrow{OS} = \overrightarrow{OM} + \overrightarrow{OM'}$. Cette lecture géométrique de l'addition est un complément très utile au Baccalauréat.

\begin{popfigure}
\begin{center}
\begin{tikzpicture}[scale=1.1, >=Stealth]
  % Axes
  \draw[->, popmuted] (-0.5,0) -- (5.2,0) node[below left] {$\mathrm{Re}$};
  \draw[->, popmuted] (0,-0.5) -- (0,4.2) node[below left] {$\mathrm{Im}$};
  % Points
  \coordinate (O) at (0,0);
  \coordinate (M1) at (3,1);
  \coordinate (M2) at (1,2);
  \coordinate (S) at (4,3);
  % Parallélogramme (pointillés)
  \draw[dashed, popblue] (M1) -- (S);
  \draw[dashed, popblue] (M2) -- (S);
  % Vecteurs
  \draw[->, very thick, popink] (O) -- (M1) node[midway, below right] {$z_1 = 3 + i$};
  \draw[->, very thick, poporange] (O) -- (M2) node[midway, above left] {$z_2 = 1 + 2i$};
  \draw[->, very thick, popgreen] (O) -- (S) node[midway, above right, popdark] {$z_1 + z_2 = 4 + 3i$};
  % Points marqués
  \fill[popink] (M1) circle (1.6pt) node[below right] {$M_1$};
  \fill[poporange] (M2) circle (1.6pt) node[above left] {$M_2$};
  \fill[popgreen] (S) circle (1.6pt) node[above right] {$S$};
  \fill[popdark] (O) circle (1.4pt) node[below left] {$O$};
\end{tikzpicture}
\end{center}
\legende{Règle du parallélogramme : l'affixe de $S$ est $z_1 + z_2 = 4 + 3i$, car $\overrightarrow{OS} = \overrightarrow{OM_1} + \overrightarrow{OM_2}$.}
\end{popfigure}

\courssub{Produit de deux nombres complexes}

\textbf{Intuition.} Pour multiplier deux complexes, on applique la double distributivité comme pour le produit $(a + b)(c + d)$, puis on remplace chaque $i^2$ par $-1$. Toute la « magie » des complexes tient dans ce remplacement.

\begin{definition}[Produit de deux nombres complexes]
Soient $z = a + ib$ et $z' = a' + ib'$. On définit leur \cle{produit} par :
\[ z \times z' = (a + ib)(a' + ib') = (aa' - bb') + i(ab' + a'b). \]
\end{definition}

\begin{experience}[D'où vient la formule du produit ?]
Développons comme au collège, en utilisant la distributivité :
\begin{align*}
(a + ib)(a' + ib') &= aa' + iab' + ia'b + i^2 bb' \\
&= aa' + i(ab' + a'b) + (-1)\,bb' \quad \text{car } i^2 = -1 \\
&= (aa' - bb') + i(ab' + a'b).
\end{align*}
Retiens la démarche, pas la formule : \textbf{on développe, on remplace $i^2$ par $-1$, on regroupe}. Cette démarche est exigible au Bac et elle évite toute erreur de mémorisation.
\end{experience}

\begin{attention}[Le piège du signe dans le produit]
L'erreur la plus fréquente consiste à oublier que $i^2 = -1$ \emph{change le signe} du terme $bb'$. Ainsi $(a + ib)(a' + ib')$ donne $aa' \mathbf{-} bb'$ pour la partie réelle, et non $aa' + bb'$. Par exemple :
\[ (2 + 3i)(1 - 4i) \neq 2 - 12 + \ldots \quad \text{faux calcul du terme réel !} \]
Le bon calcul : $aa' - bb' = 2\times 1 - 3\times(-4) = 2 + 12 = 14$. Vérifie toujours ce signe en premier.
\end{attention}

\begin{propriete}[Règles de calcul dans $\mathbb{C}$]
L'ensemble $(\mathbb{C}, +, \times)$ hérite de toutes les règles de calcul de $\mathbb{R}$. Pour tous complexes $z$, $z'$, $z''$ :
\begin{itemize}
\item \cle{Commutativité} : $z + z' = z' + z$ et $zz' = z'z$ ;
\item \cle{Associativité} : $(z + z') + z'' = z + (z' + z'')$ et $(zz')z'' = z(z'z'')$ ;
\item \cle{Distributivité} : $z(z' + z'') = zz' + zz''$ ;
\item Éléments neutres : $z + 0 = z$ et $z \times 1 = z$ ;
\item Les \cle{identités remarquables} restent valables dans $\mathbb{C}$ :
\[ (z + z')^2 = z^2 + 2zz' + z'^2, \quad (z - z')^2 = z^2 - 2zz' + z'^2, \quad (z - z')(z + z') = z^2 - z'^2. \]
\end{itemize}
\end{propriete}

\begin{exemplebox}[Identité remarquable en action]
Calculons $(1 + i)^2$ :
\[ (1 + i)^2 = 1^2 + 2\times 1 \times i + i^2 = 1 + 2i - 1 = 2i. \]
Résultat remarquable à retenir : $(1 + i)^2 = 2i$ et, de même, $(1 - i)^2 = -2i$.
\end{exemplebox}

\courssub{Puissances de $i$ : un cycle de période 4}

Calculons les premières puissances de $i$ :
\[ i^0 = 1, \qquad i^1 = i, \qquad i^2 = -1, \qquad i^3 = i^2 \times i = -i, \qquad i^4 = i^2 \times i^2 = (-1)(-1) = 1. \]
Et ensuite ? $i^5 = i^4 \times i = i$, $i^6 = -1$... On retombe sur les mêmes valeurs : les puissances de $i$ sont \cle{périodiques de période 4}.

\begin{propriete}[Cycle des puissances de $i$]
Pour tout entier naturel $n$, $i^n$ ne dépend que du reste $r$ de la division euclidienne de $n$ par $4$ :
\[ i^n = \begin{cases} 1 & \text{si } n \equiv 0 \pmod 4 \\ i & \text{si } n \equiv 1 \pmod 4 \\ -1 & \text{si } n \equiv 2 \pmod 4 \\ -i & \text{si } n \equiv 3 \pmod 4 \end{cases} \]
En particulier, $i^{4k} = 1$ pour tout $k \in \mathbb{N}$.
\end{propriete}

\begin{popfigure}
\begin{center}
\begin{tikzpicture}[scale=1.6, >=Stealth]
  \coordinate (A) at (90:1.4);
  \coordinate (B) at (180:1.4);
  \coordinate (C) at (270:1.4);
  \coordinate (D) at (0:1.4);
  \node[circle, draw=popblue, fill=popblueL, thick, minimum size=9mm] (n1) at (D) {$1$};
  \node[circle, draw=poporange, fill=poporangeL, thick, minimum size=9mm] (n2) at (A) {$i$};
  \node[circle, draw=popink, fill=poppinkL, thick, minimum size=9mm] (n3) at (B) {$-1$};
  \node[circle, draw=poppurple, fill=poppurpleL, thick, minimum size=9mm] (n4) at (C) {$-i$};
  \draw[->, very thick, popdark] (n1) to[bend left=25] node[midway, above right] {$\times i$} (n2);
  \draw[->, very thick, popdark] (n2) to[bend left=25] node[midway, above left] {$\times i$} (n3);
  \draw[->, very thick, popdark] (n3) to[bend left=25] node[midway, below left] {$\times i$} (n4);
  \draw[->, very thick, popdark] (n4) to[bend left=25] node[midway, below right] {$\times i$} (n1);
\end{tikzpicture}
\end{center}
\legende{Le cycle des puissances de $i$ : chaque multiplication par $i$ fait tourner d'un quart de tour. Après 4 multiplications, on revient à $1$.}
\end{popfigure}

\begin{methode}[Calculer $i^n$ pour un grand exposant]
\begin{enumerate}
\item Effectue la division euclidienne de $n$ par $4$ : $n = 4q + r$ avec $r \in \{0, 1, 2, 3\}$.
\item Alors $i^n = i^{4q + r} = (i^4)^q \times i^r = 1^q \times i^r = i^r$.
\item Conclus avec $i^0 = 1$, $i^1 = i$, $i^2 = -1$, $i^3 = -i$.
\end{enumerate}
\end{methode}

\begin{exemplebox}[Calcul de $i^{2025}$]
Division euclidienne : $2025 = 4 \times 506 + 1$, donc le reste est $r = 1$. Par conséquent :
\[ i^{2025} = i^{4 \times 506 + 1} = (i^4)^{506} \times i^1 = 1^{506} \times i = i. \]
\end{exemplebox}

\courssub{Inverse d'un nombre complexe non nul}

\textbf{Intuition.} Dans $\mathbb{R}$, diviser par $2$ revient à multiplier par $\tfrac{1}{2}$. Dans $\mathbb{C}$, on voudrait de même un « inverse » $z^{-1}$ tel que $z \times z^{-1} = 1$. Le problème : comment rendre réel un dénominateur imaginaire ? L'astuce consiste à exploiter l'identité remarquable $(a + ib)(a - ib) = a^2 - (ib)^2 = a^2 + b^2$, qui est un \emph{nombre réel positif}. Le complexe $a - ib$, appelé \cle{conjugué} de $z$ (nous l'étudierons en détail dans la section suivante), est la clé de voûte de cette technique.

\begin{propriete}[Inverse d'un complexe non nul]
Tout nombre complexe non nul $z = a + ib$ admet un unique \cle{inverse}, noté $\dfrac{1}{z}$ ou $z^{-1}$, donné par :
\[ \frac{1}{z} = \frac{1}{a + ib} = \frac{a - ib}{a^2 + b^2} = \frac{a}{a^2 + b^2} - i\,\frac{b}{a^2 + b^2}. \]
\end{propriete}

\begin{experience}[Pourquoi cette formule fonctionne-t-elle ?]
Multiplions le numérateur et le dénominateur de $\dfrac{1}{a + ib}$ par $a - ib$ (ce qui ne change pas la valeur du quotient) :
\[ \frac{1}{a + ib} = \frac{1 \times (a - ib)}{(a + ib)(a - ib)} = \frac{a - ib}{a^2 - (ib)^2} = \frac{a - ib}{a^2 - i^2 b^2} = \frac{a - ib}{a^2 + b^2}, \]
car $i^2 = -1$ transforme le $-$ en $+$. Le dénominateur $a^2 + b^2$ est bien un réel strictement positif (puisque $z \neq 0$), donc le résultat est bien sous forme algébrique. Vérification : $(a + ib) \times \dfrac{a - ib}{a^2 + b^2} = \dfrac{a^2 + b^2}{a^2 + b^2} = 1$. $\blacksquare$
\end{experience}

\begin{exemplebox}[Inverse de $3 + i$]
\[ \frac{1}{3 + i} = \frac{3 - i}{(3 + i)(3 - i)} = \frac{3 - i}{9 + 1} = \frac{3 - i}{10} = \frac{3}{10} - \frac{1}{10}i. \]
\end{exemplebox}

\courssub{Quotient de deux nombres complexes}

\begin{definition}[Quotient de deux complexes]
Pour $z$ et $z'$ complexes avec $z' \neq 0$, le \cle{quotient} de $z$ par $z'$ est défini par :
\[ \frac{z}{z'} = z \times \frac{1}{z'}. \]
\end{definition}

\begin{methode}[Mettre un quotient sous forme algébrique]
Pour écrire $\dfrac{z}{z'}$ sous la forme $a + ib$ :
\begin{enumerate}
\item Identifie le conjugué du dénominateur : si $z' = a' + ib'$, son conjugué est $a' - ib'$.
\item Multiplie le numérateur \emph{et} le dénominateur par ce conjugué.
\item Le dénominateur devient le réel $a'^2 + b'^2$ ; développe le numérateur avec $i^2 = -1$.
\item Sépare partie réelle et partie imaginaire : $\dfrac{\mathrm{Re}}{a'^2 + b'^2} + i\,\dfrac{\mathrm{Im}}{a'^2 + b'^2}$.
\end{enumerate}
\end{methode}

\begin{exoresolu}[Quatre calculs complets]
\textbf{1.} Calculer $(2 + 3i)(1 - 4i)$. \quad \textbf{2.} Calculer $\dfrac{3 + i}{1 - 2i}$. \quad \textbf{3.} Calculer $i^{2025}$. \quad \textbf{4.} Calculer $(2 - i)^2$.
\tcblower
\textbf{1. Produit.} On développe puis on remplace $i^2$ par $-1$ :
\begin{align*}
(2 + 3i)(1 - 4i) &= 2\times 1 - 8i + 3i - 12i^2 = 2 - 5i - 12\times(-1) \\
&= 2 - 5i + 12 = 14 - 5i.
\end{align*}
\textbf{2. Quotient.} Le conjugué du dénominateur $1 - 2i$ est $1 + 2i$ :
\begin{align*}
\frac{3 + i}{1 - 2i} &= \frac{(3 + i)(1 + 2i)}{(1 - 2i)(1 + 2i)} = \frac{3 + 6i + i + 2i^2}{1^2 + 2^2} = \frac{3 + 7i - 2}{5} = \frac{1 + 7i}{5} = \frac{1}{5} + \frac{7}{5}i.
\end{align*}
\textbf{3. Puissance de $i$.} $2025 = 4 \times 506 + 1$, donc $i^{2025} = i^1 = i$.

\textbf{4. Carré.} Avec l'identité remarquable : $(2 - i)^2 = 4 - 4i + i^2 = 4 - 4i - 1 = 3 - 4i$.
\end{exoresolu}

\begin{attention}[Interdiction formelle : la racine carrée d'un produit de négatifs]
La formule $\sqrt{ab} = \sqrt{a}\times\sqrt{b}$, valable pour $a, b$ \emph{positifs ou nuls}, devient \textbf{fausse} dès que $a$ et $b$ sont négatifs. Par exemple :
\[ \sqrt{-1} \times \sqrt{-1} = i \times i = i^2 = -1, \quad \text{alors que} \quad \sqrt{(-1)(-1)} = \sqrt{1} = 1. \]
On obtiendrait l'absurdité $-1 = 1$ ! C'est d'ailleurs ce « paradoxe » qui a longtemps fait hésiter les mathématiciens avant d'accepter les nombres imaginaires. Retiens : dans $\mathbb{C}$, on ne manipule \emph{jamais} le symbole $\sqrt{\phantom{x}}$ sur des nombres négatifs ; on écrit $\sqrt{-1} = i$ et on calcule avec $i^2 = -1$.
\end{attention}

\begin{savaistu}[Pourquoi les ingénieurs aiment tant $i$ ?]
En électrotechnique, les courants alternatifs — comme celui que distribue ENEO dans nos maisons — sont modélisés par des nombres complexes : la partie réelle représente la puissance active, la partie imaginaire la puissance réactive. Les ingénieurs additionnent et multiplient ces complexes exactement avec les règles de cette section pour dimensionner les transformateurs et éviter les délestages. Petit détail amusant : les électriciens notent l'unité imaginaire $j$ au lieu de $i$, pour ne pas la confondre avec l'intensité du courant !
\end{savaistu}

\begin{exercice}[À toi de jouer]
Soient $z = 2 - 3i$ et $z' = 1 + i$.
\begin{enumerate}
\item Calculer $z + z'$, $z - z'$ et $zz'$ sous forme algébrique.
\item Calculer $\dfrac{z}{z'}$ sous forme algébrique.
\item Calculer $i^{2026}$ et $(1 - i)^2$.
\end{enumerate}
\end{exercice}

\begin{corrige}[Exercice « À toi de jouer »]
\textbf{1.} $z + z' = (2 + 1) + i(-3 + 1) = 3 - 2i$ ; \quad $z - z' = (2 - 1) + i(-3 - 1) = 1 - 4i$ ;
\[ zz' = (2 - 3i)(1 + i) = 2 + 2i - 3i - 3i^2 = 2 - i + 3 = 5 - i. \]
\textbf{2.} On multiplie par le conjugué $1 - i$ :
\[ \frac{z}{z'} = \frac{(2 - 3i)(1 - i)}{(1 + i)(1 - i)} = \frac{2 - 2i - 3i + 3i^2}{1 + 1} = \frac{2 - 5i - 3}{2} = -\frac{1}{2} - \frac{5}{2}i. \]
\textbf{3.} $2026 = 4 \times 506 + 2$, donc $i^{2026} = i^2 = -1$. Et $(1 - i)^2 = 1 - 2i + i^2 = -2i$.
\end{corrige}

\begin{aretenir}[L'essentiel de la section]
\begin{itemize}
\item \textbf{Somme} : $(a + ib) + (a' + ib') = (a + a') + i(b + b')$ — géométriquement, règle du parallélogramme.
\item \textbf{Produit} : on développe et on remplace $i^2$ par $-1$ : $(a + ib)(a' + ib') = (aa' - bb') + i(ab' + a'b)$.
\item Toutes les règles de calcul de $\mathbb{R}$ (commutativité, associativité, distributivité, identités remarquables) restent valables dans $\mathbb{C}$.
\item \textbf{Puissances de $i$} : cycle $1 \to i \to -1 \to -i$ de période $4$ ; $i^n = i^r$ où $r$ est le reste de $n$ modulo $4$.
\item \textbf{Inverse et quotient} : on multiplie numérateur et dénominateur par le conjugué du dénominateur : $\dfrac{1}{a + ib} = \dfrac{a - ib}{a^2 + b^2}$.
\item Jamais $\sqrt{a}\sqrt{b} = \sqrt{ab}$ avec $a, b < 0$.
\end{itemize}
\end{aretenir}

Maintenant que tu maîtrises les quatre opérations, il est temps de donner un nom et un statut officiel à ce fameux $a - ib$ qui nous a sauvés dans les quotients : c'est le \emph{conjugué}, héros de la prochaine section.

\coursec{Conjugué d'un nombre complexe}

Dans la section précédente, nous avons appris à additionner, multiplier et diviser les nombres complexes. Tu as peut-être remarqué un phénomène curieux lors du calcul des quotients : pour écrire $\dfrac{1}{a+ib}$ sous forme algébrique, nous multiplions numérateur et dénominateur par $a-ib$, et le dénominateur devient alors le réel positif $a^2+b^2$. Ce nombre $a-ib$, obtenu en changeant le signe de la partie imaginaire, joue un rôle si central qu'il mérite un nom et une étude à part entière : c'est le \cle{conjugué} de $z$. Toute cette section lui est consacrée.

\courssub{Définition et premières propriétés}

\textbf{Intuition.} Un nombre complexe $z = a+ib$ est entièrement déterminé par deux réels : sa partie réelle $a$ et sa partie imaginaire $b$. L'opération la plus simple que l'on puisse imaginer consiste à garder la partie réelle et à changer le signe de la partie imaginaire. Géométriquement, cela revient à « plier » le plan le long de l'axe des réels : le point image de $z$ est envoyé sur son symétrique. Cette symétrie, d'une grande régularité, va se révéler compatible avec toutes les opérations de $\mathbb{C}$.

\begin{definition}[Conjugué d'un nombre complexe]
Soit $z = a + ib$ un nombre complexe écrit sous forme algébrique, avec $a \in \mathbb{R}$ et $b \in \mathbb{R}$. On appelle \cle{conjugué} de $z$, et on note $\bar{z}$, le nombre complexe
\[ \bar{z} = a - ib. \]
Autrement dit : $\operatorname{Re}(\bar{z}) = \operatorname{Re}(z)$ et $\operatorname{Im}(\bar{z}) = -\operatorname{Im}(z)$.
\end{definition}

\begin{exemplebox}[Lecture directe de conjugués]
Le conjugué de $z = 2+3i$ est $\bar{z} = 2-3i$ ; celui de $z = -5+i$ est $\bar{z} = -5-i$ ; celui de $z = 4i$ (imaginaire pur) est $\bar{z} = -4i$ ; celui de $z = 7$ (réel) est $\bar{z} = 7$ : un réel est son propre conjugué.
\end{exemplebox}

\begin{experience}[Découverte : somme, différence et produit de $z$ et $\bar{z}$]
Prenons $z = 3 + 2i$, donc $\bar{z} = 3 - 2i$, et calculons :
\begin{align*}
z + \bar{z} &= (3+2i) + (3-2i) = 6 = 2\times 3 = 2\operatorname{Re}(z) ;\\
z - \bar{z} &= (3+2i) - (3-2i) = 4i = 2i\times 2 = 2i\operatorname{Im}(z) ;\\
z \times \bar{z} &= (3+2i)(3-2i) = 9 - 6i + 6i - 4i^2 = 9 + 4 = 13 = 3^2 + 2^2.
\end{align*}
Les parties imaginaires s'annihilent dans la somme et le produit : on obtient des réels. Vérifie avec $z = 1 - 5i$ que le phénomène se reproduit. Nous allons maintenant le démontrer en toute généralité.
\end{experience}

\begin{propriete}[Somme, différence et produit de $z$ et $\bar{z}$]
Pour tout nombre complexe $z = a+ib$ ($a, b \in \mathbb{R}$) :
\[ z + \bar{z} = 2\operatorname{Re}(z) = 2a, \qquad z - \bar{z} = 2i\operatorname{Im}(z) = 2ib, \qquad z\,\bar{z} = a^2 + b^2 \in \mathbb{R}_{+}. \]
En particulier, $z\,\bar{z}$ est un réel positif, nul si et seulement si $z = 0$.
\end{propriete}

\begin{experience}[Démonstration guidée]
Posons $z = a+ib$, donc $\bar{z} = a - ib$.
\begin{enumerate}
\item \textbf{Somme.} $z + \bar{z} = (a+ib) + (a-ib) = 2a + i(b-b) = 2a$. Les parties imaginaires, opposées, se détruisent.
\item \textbf{Différence.} $z - \bar{z} = (a+ib) - (a-ib) = (a-a) + i(b+b) = 2ib$.
\item \textbf{Produit.} Développons :
\[ z\,\bar{z} = (a+ib)(a-ib) = a^2 - iab + iab - i^2 b^2 = a^2 - (-1)b^2 = a^2 + b^2. \]
On reconnaît l'identité remarquable $(x-y)(x+y) = x^2 - y^2$ appliquée avec $x = a$ et $y = ib$.
\item \textbf{Positivité.} Comme $a$ et $b$ sont réels, $a^2 \geq 0$ et $b^2 \geq 0$, donc $a^2 + b^2 \geq 0$. De plus, $a^2+b^2 = 0$ impose $a = 0$ et $b = 0$, c'est-à-dire $z = 0$.
\end{enumerate}
\end{experience}

\begin{aretenir}[Le réflexe du quotient]
L'égalité $z\,\bar{z} = a^2+b^2$ est la clé de tous les calculs d'inverses et de quotients : pour tout $z \neq 0$,
\[ \frac{1}{z} = \frac{\bar{z}}{z\,\bar{z}} = \frac{\bar{z}}{a^2+b^2} \qquad \text{et} \qquad \frac{z'}{z} = \frac{z'\,\bar{z}}{z\,\bar{z}}. \]
Multiplier par le conjugué « réalise » le dénominateur. C'est exactement la technique que nous avons utilisée dans la section précédente : elle trouve ici sa justification définitive.
\end{aretenir}

\begin{propriete}[Conjugué du conjugué]
Pour tout $z \in \mathbb{C}$, on a $\overline{\bar{z}} = z$ : la conjugaison est une \cle{involution}.
\end{propriete}

En effet, si $z = a+ib$, alors $\bar{z} = a-ib$ et $\overline{\bar{z}} = a - (-b)i = a + ib = z$. Changer deux fois de signe revient à ne rien changer.

\courssub{Caractérisation des réels et des imaginaires purs}

Voici l'une des idées les plus fécondes du chapitre : au lieu de montrer qu'un complexe est réel en calculant sa partie imaginaire (ce qui peut être long), on peut comparer $z$ et $\bar{z}$.

\begin{propriete}[Caractérisations par le conjugué]
Soit $z$ un nombre complexe.
\begin{itemize}
\item $z$ est \cle{réel} si et seulement si $z = \bar{z}$ ;
\item $z$ est \cle{imaginaire pur} si et seulement si $z = -\bar{z}$.
\end{itemize}
\end{propriete}

\begin{experience}[Démonstration guidée]
Écrivons $z = a+ib$ avec $a, b \in \mathbb{R}$, donc $\bar{z} = a-ib$.
\begin{enumerate}
\item \textbf{Réels.} L'égalité $z = \bar{z}$ s'écrit $a+ib = a-ib$, soit $2ib = 0$, soit $b = 0$. Or $b = 0$ signifie exactement que $z$ est réel. On a bien : $z$ réel $\iff z = \bar{z}$.
\item \textbf{Imaginaires purs.} L'égalité $z = -\bar{z}$ s'écrit $a+ib = -a+ib$, soit $2a = 0$, soit $a = 0$. Or $a = 0$ signifie exactement que $z$ est imaginaire pur. On a bien : $z$ imaginaire pur $\iff z = -\bar{z}$.
\end{enumerate}
\textbf{Remarque.} Le nombre $0$ vérifie à la fois $0 = \bar{0}$ et $0 = -\bar{0}$ : il est à la fois réel et imaginaire pur, c'est le seul.
\end{experience}

\begin{methode}[Prouver qu'un complexe est réel ou imaginaire pur]
Pour montrer qu'un nombre complexe $z$ (donné par une expression compliquée) est réel ou imaginaire pur :
\begin{enumerate}
\item Calculer $\bar{z}$ en utilisant les propriétés opératoires de la conjugaison (voir plus bas) — sans développer $z$ lui-même.
\item Comparer avec $z$ : si $\bar{z} = z$, alors $z$ est réel ; si $\bar{z} = -z$, alors $z$ est imaginaire pur.
\item Conclure, et si besoin calculer la valeur exacte de $z$ pour vérifier.
\end{enumerate}
\end{methode}

\courssub{Conjugaison et opérations}

La conjugaison « passe à travers » toutes les opérations : c'est ce qui rend la méthode précédente si puissante.

\begin{propriete}[Propriétés opératoires de la conjugaison]
Pour tous nombres complexes $z$ et $z'$, avec $z' \neq 0$ pour le quotient, et pour tout entier $n \in \mathbb{N}$ :
\[ \overline{z + z'} = \bar{z} + \bar{z}', \qquad \overline{z \times z'} = \bar{z} \times \bar{z}', \qquad \overline{\left(\frac{z}{z'}\right)} = \frac{\bar{z}}{\bar{z}'}, \qquad \overline{z^n} = \bar{z}^{\,n}. \]
\end{propriete}

\begin{experience}[Démonstration guidée à partir de la forme algébrique]
Posons $z = a+ib$ et $z' = c+id$ avec $a, b, c, d \in \mathbb{R}$.
\begin{enumerate}
\item \textbf{Somme.} $z + z' = (a+c) + i(b+d)$, donc
\[ \overline{z+z'} = (a+c) - i(b+d) = (a-ib) + (c-id) = \bar{z} + \bar{z}'. \]
\item \textbf{Produit.} $z\,z' = (ac - bd) + i(ad + bc)$, donc $\overline{z\,z'} = (ac-bd) - i(ad+bc)$. D'autre part :
\[ \bar{z}\,\bar{z}' = (a-ib)(c-id) = ac - iad - ibc + i^2 bd = (ac - bd) - i(ad+bc). \]
Les deux expressions coïncident : $\overline{z\,z'} = \bar{z}\,\bar{z}'$.
\item \textbf{Quotient.} Posons $w = \dfrac{z}{z'}$, de sorte que $z = w\,z'$. Par la propriété du produit : $\bar{z} = \overline{w\,z'} = \bar{w}\,\bar{z}'$. Comme $z' \neq 0$, on a $\bar{z}' \neq 0$, d'où $\bar{w} = \dfrac{\bar{z}}{\bar{z}'}$, c'est-à-dire $\overline{\left(\dfrac{z}{z'}\right)} = \dfrac{\bar{z}}{\bar{z}'}$.
\item \textbf{Puissance.} On raisonne par récurrence sur $n$. Pour $n = 0$ et $n = 1$, c'est immédiat. Si $\overline{z^n} = \bar{z}^{\,n}$, alors par la propriété du produit :
\[ \overline{z^{n+1}} = \overline{z^n \times z} = \overline{z^n} \times \bar{z} = \bar{z}^{\,n} \times \bar{z} = \bar{z}^{\,n+1}. \]
La propriété est donc vraie pour tout $n \in \mathbb{N}$.
\end{enumerate}
\end{experience}

\begin{attention}[Erreurs fréquentes autour de la conjugaison]
\begin{itemize}
\item Devant $\overline{(2+i)(3-4i)}$, beaucoup d'élèves développent d'abord le produit, puis conjuguent : c'est long et source d'erreurs de signe. Utilise directement les propriétés opératoires :
\[ \overline{(2+i)(3-4i)} = \overline{2+i}\;\times\;\overline{3-4i} = (2-i)(3+4i). \]
\item La conjugaison ne change que le signe de la partie \emph{imaginaire} : écrire $\overline{2+3i} = -2-3i$ est faux, on obtient $2-3i$.
\item Attention : $\bar{z} \times z$ n'est pas $z^2$ ! On a $z\,\bar{z} = a^2+b^2$, réel positif, alors que $z^2 = (a+ib)^2 = a^2 - b^2 + 2iab$ est en général un complexe non réel.
\item Pour le quotient, n'oublie pas la condition $z' \neq 0$ avant d'écrire $\overline{\left(\dfrac{z}{z'}\right)} = \dfrac{\bar{z}}{\bar{z}'}$.
\end{itemize}
\end{attention}

\begin{exoresolu}[Montrer qu'un complexe est un imaginaire pur]
On considère le nombre complexe $z = \dfrac{1+i}{1-i}$. Montrer, sans calculer $z$ sous forme algébrique, que $z$ est un imaginaire pur, puis déterminer sa valeur.
\tcblower
\textbf{Étape 1 : calcul de $\bar{z}$ par les propriétés opératoires.} Par conjugaison d'un quotient puis d'une somme :
\[ \bar{z} = \overline{\left(\frac{1+i}{1-i}\right)} = \frac{\overline{1+i}}{\overline{1-i}} = \frac{1-i}{1+i}. \]

\textbf{Étape 2 : comparaison de $\bar{z}$ et $-z$.} On a $-z = -\dfrac{1+i}{1-i} = \dfrac{-(1+i)}{1-i} = \dfrac{1+i}{i-1}$. Multiplions numérateur et dénominateur de $\bar{z}$ par $-1$ sous la forme qui arrange : écrivons
\[ \bar{z} = \frac{1-i}{1+i} = \frac{-(i-1)}{1+i}\cdot\frac{1}{1} \quad \text{... plutôt, comparons par produits en croix.} \]
Plus simplement : $\bar{z} = \dfrac{1-i}{1+i}$. Multiplions numérateur et dénominateur par $i$ :
\[ \bar{z} = \frac{(1-i)\,i}{(1+i)\,i} = \frac{i - i^2}{i + i^2} = \frac{i+1}{i-1} = \frac{1+i}{i-1} = -\frac{1+i}{1-i} = -z. \]
Comme $\bar{z} = -z$, la caractérisation des imaginaires purs s'applique : $z$ est un \textbf{imaginaire pur}.

\textbf{Étape 3 : valeur exacte de $z$.} On multiplie par le conjugué du dénominateur :
\[ z = \frac{1+i}{1-i} = \frac{(1+i)(1+i)}{(1-i)(1+i)} = \frac{(1+i)^2}{1^2+1^2} = \frac{1 + 2i + i^2}{2} = \frac{2i}{2} = i. \]
On vérifie au passage : $\bar{z} = -i = -z$, conformément à l'étape 2.

\textbf{Conclusion.} $z$ est un imaginaire pur et $z = i$.
\end{exoresolu}

\begin{exemplebox}[Application des propriétés opératoires]
Soit $z = \dfrac{(3+i)^2}{2-i}$. Sans calculer $z$, calculons $\bar{z}$ :
\[ \bar{z} = \frac{\overline{(3+i)^2}}{\overline{2-i}} = \frac{(\overline{3+i})^2}{2+i} = \frac{(3-i)^2}{2+i}. \]
Vérification par le calcul direct : $(3+i)^2 = 9 + 6i + i^2 = 8 + 6i$, donc
\[ z = \frac{8+6i}{2-i} = \frac{(8+6i)(2+i)}{(2-i)(2+i)} = \frac{16 + 8i + 12i + 6i^2}{5} = \frac{10 + 20i}{5} = 2 + 4i. \]
Et effectivement $\bar{z} = 2 - 4i$, ce que redonne le calcul direct :
\[ \frac{(3-i)^2}{2+i} = \frac{8-6i}{2+i} = \frac{(8-6i)(2-i)}{(2+i)(2-i)} = \frac{16 - 8i - 12i + 6i^2}{5} = \frac{10 - 20i}{5} = 2 - 4i. \]
\end{exemplebox}

\courssub{Interprétation géométrique}

Plaçons-nous dans le plan muni d'un repère orthonormé $(O\,;\vec{u},\vec{v})$. Au complexe $z = a+ib$ on associe le point $M$ de coordonnées $(a\,;b)$, appelé \cle{image} de $z$. Le point image de $\bar{z} = a-ib$ a pour coordonnées $(a\,;-b)$ : c'est le symétrique de $M$ par rapport à l'axe des abscisses, c'est-à-dire l'\cle{axe des réels}.

\begin{propriete}[Symétrie associée à la conjugaison]
Dans le plan complexe, les points images de $z$ et de $\bar{z}$ sont symétriques par rapport à l'axe des réels $(Ox)$.
\end{propriete}

Cette observation géométrique éclaire tous les résultats de la section : un réel est un point \emph{sur} l'axe des réels, donc invariant par la symétrie ($z = \bar{z}$) ; un imaginaire pur est un point sur l'axe des ordonnées, envoyé sur son opposé ($\bar{z} = -z$) ; et la distance $OM$ est la même pour $z$ et $\bar{z}$, ce qui annonce l'égalité des modules que nous verrons dans la section suivante.

\begin{popfigure}
\begin{center}
\begin{tikzpicture}[scale=1.05,>=Stealth]
% Grille
\draw[popline,very thin] (-1.5,-3.8) grid[step=1] (4.5,3.8);
% Axes
\draw[->,thick,popdark] (-1.5,0) -- (4.5,0) node[below right] {axe des réels};
\draw[->,thick,popdark] (0,-3.8) -- (0,3.8) node[above left] {axe des imaginaires};
\draw[popdark] (0,0) node[below left] {$O$};
% Graduations
\foreach \x in {1,2,3,4} \draw[popdark] (\x,0.08) -- (\x,-0.08) node[below,scale=0.8] {\x};
\foreach \y in {-3,-2,-1,1,2,3} \draw[popdark] (0.08,\y) -- (-0.08,\y) node[left,scale=0.8] {\y};
% Points z et zbar
\coordinate (M) at (2,3);
\coordinate (Mb) at (2,-3);
% Pointillés de construction
\draw[dashed,popmuted] (M) -- (Mb);
\draw[dashed,popmuted] (2,3) -- (0,3);
\draw[dashed,popmuted] (2,-3) -- (0,-3);
% Vecteurs
\draw[->,very thick,popblue] (0,0) -- (M);
\draw[->,very thick,poporange] (0,0) -- (Mb);
% Points
\fill[popblue] (M) circle (2.5pt) node[above right,popblue] {$M(z=2+3i)$};
\fill[poporange] (Mb) circle (2.5pt) node[below right,poporange] {$M'(\bar{z}=2-3i)$};
% Annotation symétrie
\draw[<->,popgreen,thick] (2.6,1.2) -- (2.6,-1.2) node[midway,right,popgreen,scale=0.9] {symétrie / $(Ox)$};
\end{tikzpicture}
\end{center}
\legende{Les images de $z = 2+3i$ et de $\bar{z} = 2-3i$ sont symétriques par rapport à l'axe des réels : la conjugaison est une symétrie orthogonale d'axe $(Ox)$.}
\end{popfigure}

\begin{savaistu}[Pourquoi le mot « conjugué » ?]
Le mot vient du latin \emph{conjugare}, « unir par le joug », « marier ensemble » : $z$ et $\bar{z}$ forment un couple inséparable. Cette idée de couple est omniprésente : quand un polynôme à coefficients réels admet une racine complexe non réelle $z_0$, alors $\bar{z}_0$ est automatiquement racine aussi — les racines non réelles « viennent par paires ». C'est ce phénomène que tu exploiteras pour résoudre les équations du second degré à discriminant négatif : les deux solutions seront toujours conjuguées l'une de l'autre. Au Cameroun comme ailleurs, les ingénieurs en télécommunications (chez MTN, Orange ou CAMTEL) manipulent quotidiennement des signaux représentés par des nombres complexes, et la conjugaison y est l'outil de base du calcul de puissance d'un signal : la puissance se calcule précisément via le produit $z\,\bar{z}$.
\end{savaistu}

\begin{exercice}[À toi de jouer]
\begin{enumerate}
\item Soit $z = \dfrac{2+3i}{1+2i}$. Calculer $\bar{z}$ de deux façons : en calculant d'abord $z$, puis en utilisant les propriétés opératoires.
\item Montrer que $Z = \dfrac{1+2i}{1-2i} - \dfrac{1-2i}{1+2i}$ est un imaginaire pur, puis le calculer.
\item Montrer que pour tout $z \in \mathbb{C}$, le nombre $z + \bar{z}$ est réel et le nombre $z - \bar{z}$ est imaginaire pur.
\end{enumerate}
\end{exercice}

\begin{corrige}[Exercice]
\begin{enumerate}
\item \textbf{Première façon.} $z = \dfrac{(2+3i)(1-2i)}{(1+2i)(1-2i)} = \dfrac{2 - 4i + 3i - 6i^2}{1+4} = \dfrac{8 - i}{5} = \dfrac{8}{5} - \dfrac{1}{5}i$. Donc $\bar{z} = \dfrac{8}{5} + \dfrac{1}{5}i$.

\textbf{Deuxième façon.} $\bar{z} = \dfrac{\overline{2+3i}}{\overline{1+2i}} = \dfrac{2-3i}{1-2i} = \dfrac{(2-3i)(1+2i)}{(1-2i)(1+2i)} = \dfrac{2+4i-3i-6i^2}{5} = \dfrac{8+i}{5}$. Même résultat.
\item On a $\bar{Z} = \overline{\dfrac{1+2i}{1-2i}} - \overline{\dfrac{1-2i}{1+2i}} = \dfrac{1-2i}{1+2i} - \dfrac{1+2i}{1-2i} = -\left(\dfrac{1+2i}{1-2i} - \dfrac{1-2i}{1+2i}\right) = -Z$. Donc $Z$ est imaginaire pur. Calcul : $\dfrac{1+2i}{1-2i} = \dfrac{(1+2i)^2}{(1-2i)(1+2i)} = \dfrac{1+4i+4i^2}{5} = \dfrac{-3+4i}{5}$ et de même $\dfrac{1-2i}{1+2i} = \dfrac{-3-4i}{5}$, d'où $Z = \dfrac{-3+4i}{5} - \dfrac{-3-4i}{5} = \dfrac{8i}{5}$.
\item $\overline{z+\bar{z}} = \bar{z} + \overline{\bar{z}} = \bar{z} + z = z + \bar{z}$ : le nombre est égal à son conjugué, donc réel. De même $\overline{z - \bar{z}} = \bar{z} - \overline{\bar{z}} = \bar{z} - z = -(z - \bar{z})$ : le nombre est l'opposé de son conjugué, donc imaginaire pur. On retrouve $z + \bar{z} = 2\operatorname{Re}(z)$ et $z - \bar{z} = 2i\operatorname{Im}(z)$.
\end{enumerate}
\end{corrige}

\begin{aretenir}[L'essentiel de la section]
\begin{itemize}
\item Le conjugué de $z = a+ib$ est $\bar{z} = a - ib$ ; géométriquement, c'est la symétrie par rapport à l'axe des réels.
\item $z + \bar{z} = 2\operatorname{Re}(z)$, $z - \bar{z} = 2i\operatorname{Im}(z)$, $z\,\bar{z} = a^2 + b^2 \in \mathbb{R}_+$.
\item $z$ réel $\iff z = \bar{z}$ ; $z$ imaginaire pur $\iff z = -\bar{z}$.
\item La conjugaison respecte somme, produit, quotient et puissances : conjugue d'abord, calcule ensuite.
\end{itemize}
Dans la section suivante, nous donnerons un nom au réel $\sqrt{a^2+b^2} = \sqrt{z\,\bar{z}}$ : le \emph{module} de $z$, qui mesure la distance de l'image de $z$ à l'origine.
\end{aretenir}

\coursec{Module d'un nombre complexe}

Dans la section précédente, nous avons découvert le conjugué d'un nombre complexe et une formule remarquable : pour $z = a + ib$, le produit $z\,\bar{z} = a^2 + b^2$ est toujours un \cle{réel positif ou nul}. Cette quantité $a^2 + b^2$ n'est pas apparue par hasard : si tu places le point $M(a\,;\,b)$ image de $z$ dans le plan muni d'un repère orthonormé $(O\,;\vec{u},\vec{v})$, alors le théorème de Pythagore te donne $OM^2 = a^2 + b^2$. Autrement dit, $z\,\bar{z} = OM^2$ ! Le conjugué nous offre donc un moyen \emph{algébrique} de mesurer une \emph{distance géométrique}. C'est exactement cette idée que formalise la notion de \cle{module}, l'un des outils les plus puissants du chapitre : tu le retrouveras dans presque tous les sujets de Baccalauréat portant sur les nombres complexes.

\courssub{Définition du module}

\courssub{Une intuition géométrique : mesurer une distance}

Dans $\mathbb{R}$, la valeur absolue $|x|$ mesure la distance du point d'abscisse $x$ à l'origine de la droite graduée. Dans $\mathbb{C}$, le plan remplace la droite : il nous faut un outil qui mesure la distance d'un point $M$ à l'origine $O$ du repère. Cet outil, c'est le module.

\begin{definition}[Module d'un nombre complexe]
Soit $z = a + ib$ un nombre complexe, avec $a$ et $b$ réels. On appelle \cle{module} de $z$, noté $|z|$, le réel positif ou nul défini par :
\[ |z| = \sqrt{a^2 + b^2}. \]
Si $M$ est le point d'affixe $z$ dans le plan muni du repère orthonormé $(O\,;\vec{u},\vec{v})$, alors :
\[ |z| = OM. \]
\end{definition}

\begin{attention}[Le module est un réel positif]
Le module d'un nombre complexe est toujours un \cle{nombre réel positif ou nul}, jamais un complexe. Écrire $|z| = 3 + 4i$ n'a aucun sens. De plus, la racine carrée porte sur la \emph{somme} $a^2 + b^2$ toute entière : $|z| = \sqrt{a^2+b^2}$ et non $\sqrt{a^2} + \sqrt{b^2}$.
\end{attention}

\begin{exemplebox}[Premiers calculs de modules]
\begin{itemize}
\item $|3 + 4i| = \sqrt{3^2 + 4^2} = \sqrt{9 + 16} = \sqrt{25} = 5$ ;
\item $|-2i| = \sqrt{0^2 + (-2)^2} = \sqrt{4} = 2$ ;
\item $|7| = \sqrt{7^2 + 0^2} = 7$ : pour un réel, le module coïncide avec la valeur absolue, ce qui justifie la notation $|z|$ ;
\item $\left|\dfrac{1}{2} - i\dfrac{\sqrt{3}}{2}\right| = \sqrt{\dfrac{1}{4} + \dfrac{3}{4}} = \sqrt{1} = 1$.
\end{itemize}
\end{exemplebox}

\begin{popfigure}
\begin{center}
\begin{tikzpicture}[scale=0.9,>=Stealth]
% Axes
\draw[->,popmuted] (-0.8,0) -- (4.6,0) node[below left] {$\vec{u}$};
\draw[->,popmuted] (0,-0.8) -- (0,4.8) node[below left] {$\vec{v}$};
% Graduations
\foreach \x in {1,2,3,4} \draw[popmuted] (\x,0.06) -- (\x,-0.06) node[below] {\small $\x$};
\foreach \y in {1,2,3,4} \draw[popmuted] (0.06,\y) -- (-0.06,\y) node[left] {\small $\y$};
% Projeté
\draw[dashed,popblue] (3,0) -- (3,4);
\draw[dashed,popblue] (0,4) -- (3,4);
% Triangle rectangle : angle droit
\draw[popblue] (3,0.3) -- (2.7,0.3) -- (2.7,0);
% Segment OM
\draw[very thick,poporange] (0,0) -- (3,4) node[midway,above left,popdark] {$OM = |z| = 5$};
% Points
\fill[popink] (0,0) circle (2pt) node[below left] {$O$};
\fill[popink] (3,4) circle (2.5pt) node[above right] {$M(3\,;\,4)$ image de $z = 3+4i$};
\fill[popblue] (3,0) circle (2pt) node[below right] {\small $H(3\,;\,0)$};
% Annotations
\node[popblue] at (3.35,2) {\small $4$};
\node[popblue] at (1.5,-0.45) {\small $3$};
\end{tikzpicture}
\end{center}
\legende{Le point $M$ image de $z = 3+4i$. Le triangle $OHM$ est rectangle en $H$ : Pythagore donne $OM = \sqrt{3^2+4^2} = 5 = |z|$.}
\end{popfigure}

\courssub{Module et distance entre deux points}

Le module mesure plus généralement la distance entre \emph{deux} points du plan, ce qui en fait un outil géométrique précieux au Baccalauréat.

\begin{propriete}[Interprétation géométrique : distance]
Soient $A$ et $B$ deux points d'affixes respectives $z_A$ et $z_B$. Alors :
\[ |z_B - z_A| = AB. \]
En particulier, $|z| = |z - 0| = OM$ où $M$ est l'image de $z$.
\end{propriete}

\begin{experience}[Démonstration guidée]
Posons $z_A = a + ib$ et $z_B = c + id$ avec $a, b, c, d$ réels.
\begin{enumerate}
\item Écris $z_B - z_A$ sous forme algébrique : $z_B - z_A = (c - a) + i(d - b)$.
\item Applique la définition du module : $|z_B - z_A| = \sqrt{(c-a)^2 + (d-b)^2}$.
\item Reconnais la formule de la distance entre deux points dans un repère orthonormé : $AB = \sqrt{(x_B - x_A)^2 + (y_B - y_A)^2}$.
\item Conclus : $|z_B - z_A| = AB$.
\end{enumerate}
\end{experience}

\begin{exemplebox}[Distance entre deux points]
Soient $A$ d'affixe $z_A = 1 - 2i$ et $B$ d'affixe $z_B = 4 + 2i$. Alors :
\[ AB = |z_B - z_A| = |(4-1) + (2-(-2))i| = |3 + 4i| = \sqrt{9+16} = 5. \]
\end{exemplebox}

\begin{savaistu}[Le module au quotidien]
Quand ton téléphone affiche la distance « à vol d'oiseau » entre Douala et Yaoundé sur une carte, le logiciel calcule exactement un module : chaque ville est repérée par ses coordonnées, et la distance est $\sqrt{(\Delta x)^2 + (\Delta y)^2}$. Les ingénieurs de CAMTEL ou de MTN utilisent le même principe pour estimer la portée d'une antenne relais : la zone couverte est un disque, c'est-à-dire l'ensemble des points $M$ tels que $|z_M - z_A| \leq R$, où $A$ est l'antenne et $R$ sa portée.
\end{savaistu}

\courssub{La propriété fondamentale : $z\,\bar{z} = |z|^2$}

Voici le pont définitif entre le conjugué (section précédente) et le module. C'est la formule la plus utilisée du chapitre.

\begin{propriete}[Propriété fondamentale — démonstration exigible]
Pour tout nombre complexe $z$ :
\[ z\,\bar{z} = |z|^2. \]
\textbf{Démonstration.} Posons $z = a + ib$ avec $a, b \in \mathbb{R}$. Alors $\bar{z} = a - ib$, donc :
\begin{align*}
z\,\bar{z} &= (a + ib)(a - ib) \\
&= a^2 - (ib)^2 \quad \text{(identité remarquable } (x-y)(x+y) = x^2 - y^2\text{)} \\
&= a^2 - i^2 b^2 = a^2 + b^2 \quad \text{car } i^2 = -1 \\
&= |z|^2.
\end{align*}
\end{propriete}

\begin{aretenir}[Une conséquence précieuse : l'inverse]
Si $z \neq 0$, alors $|z|^2 \neq 0$ et on peut diviser par $z\,\bar{z}$. De $z\,\bar{z} = |z|^2$, on tire :
\[ \frac{1}{z} = \frac{\bar{z}}{|z|^2}. \]
C'est exactement la technique vue dans la section sur les opérations : pour rendre réel le dénominateur de $\dfrac{1}{a+ib}$, on multiplie par le conjugué, et le dénominateur devient $|z|^2 = a^2 + b^2$.
\end{aretenir}

\courssub{Propriétés opératoires du module}

\courssub{Module d'un produit, d'un quotient, d'une puissance}

\begin{propriete}[Module d'un produit — démonstration exigible]
Pour tous nombres complexes $z$ et $z'$ :
\[ |z \cdot z'| = |z| \times |z'|. \]
\textbf{Démonstration.} Plutôt que de développer $(a+ib)(c+id)$ (long et pénible), utilisons la propriété fondamentale et les propriétés du conjugué :
\begin{align*}
|z z'|^2 &= (zz')\,\overline{(zz')} & \text{(propriété fondamentale appliquée à } zz'\text{)} \\
&= zz' \times \bar{z}\,\bar{z}' & \text{(conjugué d'un produit)} \\
&= (z\bar{z})(z'\bar{z}') & \text{(commutativité dans } \mathbb{C}\text{)} \\
&= |z|^2 \, |z'|^2 = \left(|z|\,|z'|\right)^2.
\end{align*}
Les deux nombres $|zz'|$ et $|z|\,|z'|$ sont positifs et ont le même carré : ils sont donc égaux.
\end{propriete}

\begin{propriete}[Autres propriétés opératoires]
Pour tous complexes $z$ et $z'$ (avec $z' \neq 0$ pour le quotient), et pour tout entier naturel $n$ :
\[ \left|\frac{z}{z'}\right| = \frac{|z|}{|z'|} \qquad ; \qquad |z^n| = |z|^n \qquad ; \qquad |\bar{z}| = |z| \qquad ; \qquad |-z| = |z|. \]
\textbf{Justifications rapides.}
\begin{itemize}
\item \textbf{Quotient :} de $z = \dfrac{z}{z'} \times z'$, on tire $|z| = \left|\dfrac{z}{z'}\right| \times |z'|$ par la propriété du produit, d'où le résultat en divisant par $|z'| \neq 0$.
\item \textbf{Puissance :} récurrence immédiate à partir de $|zz'| = |z|\,|z'|$.
\item \textbf{Conjugué :} si $z = a + ib$, alors $|\bar{z}| = \sqrt{a^2 + (-b)^2} = \sqrt{a^2+b^2} = |z|$.
\item \textbf{Opposé :} $|-z| = |-1 \times z| = |-1| \times |z| = 1 \times |z| = |z|$.
\end{itemize}
\end{propriete}

\begin{attention}[Pièges fréquents avec le module]
\begin{itemize}
\item \textbf{Le module ne se distribue pas sur la somme :} en général, $|z + z'| \neq |z| + |z'|$. Contre-exemple : $z = 1$ et $z' = -1$ donnent $|z + z'| = 0$ mais $|z| + |z'| = 2$.
\item \textbf{Le module n'est pas la valeur absolue de chaque partie :} $|3 + 4i| \neq |3| + |4i| = 3 + 4 = 7$ ; la bonne valeur est $5$. On ne calcule jamais $|a| + |b|$, mais $\sqrt{a^2 + b^2}$.
\item \textbf{Ne confonds pas $|z|^2$ et $z^2$ :} $|z|^2 = z\bar{z}$ est un réel positif, alors que $z^2$ est un complexe quelconque. Par exemple, pour $z = i$ : $|z|^2 = 1$ mais $z^2 = -1$.
\end{itemize}
\end{attention}

\courssub{L'inégalité triangulaire}

Si le module ne se distribue pas sur la somme, que peut-on dire de $|z + z'|$ ? La géométrie répond : dans un triangle, la longueur d'un côté est toujours inférieure ou égale à la somme des longueurs des deux autres.

\begin{propriete}[Inégalité triangulaire — admise]
Pour tous nombres complexes $z$ et $z'$ :
\[ |z + z'| \leq |z| + |z'|. \]
Il y a égalité si et seulement si $z$ et $z'$ ont le même argument (images alignées avec $O$ du même côté), c'est-à-dire si l'un est un multiple positif de l'autre.
\end{propriete}

\textbf{Interprétation géométrique.} Soient $M$ et $M'$ les images de $z$ et $z'$. En construisant le point $S$ d'affixe $z + z'$ par la relation de Chasles vectorielle $\overrightarrow{OS} = \overrightarrow{OM} + \overrightarrow{OM'}$, on obtient un parallélogramme (ou un triangle) dans lequel $OS \leq OM + OM'$, c'est-à-dire $|z+z'| \leq |z| + |z'|$ : le chemin direct de $O$ à $S$ est plus court que le détour par $M$.

\begin{exemplebox}[Vérification sur un exemple]
Prenons $z = 3 + 4i$ et $z' = 1 - 2i$. Alors $z + z' = 4 + 2i$, donc :
\[ |z + z'| = \sqrt{16 + 4} = \sqrt{20} \approx 4{,}47 \quad \text{et} \quad |z| + |z'| = 5 + \sqrt{5} \approx 7{,}24. \]
On a bien $|z + z'| \leq |z| + |z'|$, avec inégalité stricte car $z$ et $z'$ ne sont pas alignés avec $O$.
\end{exemplebox}

\courssub{Cas particuliers remarquables}

\begin{propriete}[Module nul et module unité]
\begin{itemize}
\item $|z| = 0 \iff z = 0$. En effet, $|z| = 0 \iff a^2 + b^2 = 0 \iff a = 0$ \emph{et} $b = 0$ (une somme de carrés de réels n'est nulle que si chaque terme est nul).
\item $|z| = 1 \iff z\,\bar{z} = 1 \iff \bar{z} = \dfrac{1}{z}$ (pour $z \neq 0$). Géométriquement, les points images des complexes de module $1$ forment le \cle{cercle trigonométrique} : le cercle de centre $O$ et de rayon $1$.
\end{itemize}
\end{propriete}

\begin{popfigure}
\begin{center}
\begin{tikzpicture}[scale=2.2,>=Stealth]
% Axes
\draw[->,popmuted] (-1.5,0) -- (1.6,0) node[below left] {$\vec{u}$};
\draw[->,popmuted] (0,-1.4) -- (0,1.5) node[below left] {$\vec{v}$};
% Cercle trigonométrique
\draw[very thick,popblue] (0,0) circle (1);
% Points de module 1
\fill[popink] (1,0) circle (1.5pt) node[below right] {\small $1$};
\fill[popink] (0,1) circle (1.5pt) node[above right] {\small $i$};
\fill[popink] (-1,0) circle (1.5pt) node[below left] {\small $-1$};
\fill[popink] (0,-1) circle (1.5pt) node[below right] {\small $-i$};
\fill[poporange] (0.5,0.866) circle (1.5pt) node[above right] {\small $\dfrac{1}{2} + i\dfrac{\sqrt{3}}{2}$};
\fill[poporange] (-0.707,0.707) circle (1.5pt) node[above left] {\small $-\dfrac{\sqrt{2}}{2} + i\dfrac{\sqrt{2}}{2}$};
% Rayons
\draw[poporange,thick] (0,0) -- (0.5,0.866);
\draw[poporange,thick] (0,0) -- (-0.707,0.707);
\fill[popink] (0,0) circle (1.2pt) node[below left] {\small $O$};
\end{tikzpicture}
\end{center}
\legende{Le cercle trigonométrique : ensemble des points images des complexes de module $1$. Tous les rayons (en orange) ont pour longueur $1$.}
\end{popfigure}

\begin{savaistu}[Pourquoi le module 1 est-il si important ?]
Les complexes de module $1$ sont les « rotations en puissance » : nous verrons dans la suite du cours (forme trigonométrique) que tout complexe de module $1$ s'écrit $\cos\theta + i\sin\theta$, et que multiplier par un tel nombre revient à \emph{tourner} autour de $O$ sans changer les distances. C'est le principe utilisé en télécommunications (modulation de phase des signaux) et en infographie pour faire pivoter les images.
\end{savaistu}

\courssub{Application : calculer un module sans forme algébrique}

Voici l'un des savoir-faire les plus rentables au Baccalauréat : les propriétés opératoires permettent de calculer le module d'une expression compliquée \emph{sans jamais la développer}.

\begin{methode}[Calculer le module d'un produit ou d'un quotient]
\begin{enumerate}
\item \textbf{Ne développe pas} l'expression : ce serait long et source d'erreurs.
\item Applique les propriétés : $|zz'| = |z|\,|z'|$ et $\left|\dfrac{z}{z'}\right| = \dfrac{|z|}{|z'|}$ pour éclater le module en modules de facteurs simples.
\item Calcule chaque petit module avec $|a + ib| = \sqrt{a^2+b^2}$.
\item Simplifie le résultat (les racines carrées se combinent souvent agréablement).
\end{enumerate}
\end{methode}

\begin{exemplebox}[Un calcul spectaculaire]
Calculons $\left|\dfrac{(2-i)(1+3i)}{1+i}\right|$ \textbf{sans} mettre le quotient sous forme algébrique :
\begin{align*}
\left|\dfrac{(2-i)(1+3i)}{1+i}\right| &= \dfrac{|2-i| \times |1+3i|}{|1+i|} = \dfrac{\sqrt{2^2+(-1)^2} \times \sqrt{1^2+3^2}}{\sqrt{1^2+1^2}} \\
&= \dfrac{\sqrt{5} \times \sqrt{10}}{\sqrt{2}} = \sqrt{\dfrac{5 \times 10}{2}} = \sqrt{25} = 5.
\end{align*}
Compare avec la méthode « brutale » : développer $(2-i)(1+3i) = 5 + 5i$, puis calculer $\dfrac{5+5i}{1+i}$, multiplier par le conjugué... pour finalement trouver $5$ et $|5| = 5$. La méthode par propriétés est plus rapide, plus sûre, et très appréciée des correcteurs.
\end{exemplebox}

\begin{exoresolu}[Module et distance au Baccalauréat]
Le plan est muni d'un repère orthonormé $(O\,;\vec{u},\vec{v})$. On considère les points $A$, $B$ et $C$ d'affixes respectives $z_A = 1 + i$, $z_B = 3 - i$ et $z_C = 2 + 2i$.
\begin{enumerate}
\item Calculer les modules de $z_A$, $z_B - z_A$ et $z_C - z_A$.
\item En déduire la nature du triangle $ABC$.
\item Déterminer l'ensemble des points $M$ d'affixe $z$ tels que $|z - z_A| = |z - z_C|$.
\end{enumerate}
\tcblower
\textbf{Solution détaillée.}
\begin{enumerate}
\item $|z_A| = |1+i| = \sqrt{1^2+1^2} = \sqrt{2}$.
\[ |z_B - z_A| = |(3-1) + (-1-1)i| = |2 - 2i| = \sqrt{4+4} = \sqrt{8} = 2\sqrt{2}. \]
\[ |z_C - z_A| = |(2-1) + (2-1)i| = |1 + i| = \sqrt{2}. \]
\item D'après l'interprétation géométrique : $AB = |z_B - z_A| = 2\sqrt{2}$ et $AC = |z_C - z_A| = \sqrt{2}$. Calculons aussi $BC = |z_C - z_B| = |-1 + 3i| = \sqrt{1+9} = \sqrt{10}$. Vérifions Pythagore :
\[ AB^2 + AC^2 = 8 + 2 = 10 = BC^2. \]
Par la réciproque du théorème de Pythagore, le triangle $ABC$ est \textbf{rectangle en $A$}.
\item La condition $|z - z_A| = |z - z_C|$ s'écrit $AM = CM$ : $M$ est équidistant de $A$ et de $C$. L'ensemble cherché est donc la \textbf{médiatrice du segment $[AC]$}.
\end{enumerate}
\end{exoresolu}

\begin{exercice}[À toi de jouer]
\begin{enumerate}
\item Calculer les modules suivants \textbf{sans développer} : $\left|(3-4i)(1+i\sqrt{2})\right|$ ; $\left|\dfrac{5-i}{2+3i}\right|$ ; $\left|(1-i)^6\right|$.
\item Soit $z = \dfrac{3+4i}{5}$. Montrer que $|z| = 1$ et en déduire $\bar{z}$ en fonction de $z$.
\item Soient $A(2 + i)$ et $B(-1 + 5i)$. Déterminer la distance $AB$, puis l'ensemble des points $M$ d'affixe $z$ vérifiant $|z - z_A| = 3$.
\end{enumerate}
\end{exercice}

\begin{corrige}[Exercice : À toi de jouer]
\begin{enumerate}
\item $\left|(3-4i)(1+i\sqrt{2})\right| = |3-4i| \times |1+i\sqrt{2}| = \sqrt{9+16} \times \sqrt{1+2} = 5\sqrt{3}$.
\[ \left|\dfrac{5-i}{2+3i}\right| = \dfrac{|5-i|}{|2+3i|} = \dfrac{\sqrt{25+1}}{\sqrt{4+9}} = \dfrac{\sqrt{26}}{\sqrt{13}} = \sqrt{2}. \]
$\left|(1-i)^6\right| = |1-i|^6 = (\sqrt{2})^6 = 2^3 = 8$.
\item $|z| = \dfrac{|3+4i|}{|5|} = \dfrac{5}{5} = 1$. Comme $|z| = 1$, on a $z\bar{z} = 1$, donc $\bar{z} = \dfrac{1}{z}$.
\item $AB = |z_B - z_A| = |-3 + 4i| = \sqrt{9+16} = 5$. La condition $|z - z_A| = 3$ signifie $AM = 3$ : l'ensemble cherché est le \textbf{cercle de centre $A$ et de rayon $3$}.
\end{enumerate}
\end{corrige}

\begin{aretenir}[L'essentiel du module]
\begin{itemize}
\item $|a + ib| = \sqrt{a^2+b^2}$ : un réel positif, égal à la distance $OM$.
\item $|z_B - z_A| = AB$ : le module traduit les distances, donc la géométrie.
\item $z\bar{z} = |z|^2$ : la formule-pont entre conjugué et module.
\item $|zz'| = |z|\,|z'|$, $\left|\dfrac{z}{z'}\right| = \dfrac{|z|}{|z'|}$, $|z^n| = |z|^n$, $|\bar{z}| = |-z| = |z|$.
\item $|z + z'| \leq |z| + |z'|$ (inégalité triangulaire) ; jamais d'égalité automatique !
\item $|z| = 1$ : $M$ appartient au cercle trigonométrique.
\end{itemize}
\end{aretenir}

Le module mesure la \emph{taille} d'un nombre complexe, sa distance à l'origine. Mais un point du plan n'est pas seulement repéré par sa distance à $O$ : il faut aussi une \emph{direction}, un angle. Cette seconde information, appelée \emph{argument}, sera au cœur de la forme trigonométrique que nous étudierons dans une prochaine leçon. Auparavant, la section suivante te montrera comment le module et le conjugué permettent de résoudre un problème classique et redouté : extraire les \emph{racines carrées} d'un nombre complexe.

\coursec{Racines carrées d'un nombre complexe}

Dans la section précédente, nous avons défini le \cle{module} d'un nombre complexe $z = a + ib$ par $|z| = \sqrt{a^2 + b^2}$, et nous avons vu qu'il mesure la distance de l'image de $z$ à l'origine. Ce module va maintenant nous rendre un service remarquable : il va transformer un problème apparemment difficile --- extraire la « racine carrée » d'un nombre complexe --- en un simple système d'équations que tu sais résoudre depuis la Seconde. C'est toute la beauté de cette section : une astuce, un système, et le tour est joué.

\courssub{Qu'est-ce qu'une racine carrée ? Intuition et définition}

Dans $\mathbb{R}$, tu sais depuis longtemps que $3$ et $-3$ sont les racines carrées de $9$, car $3^2 = 9$ et $(-3)^2 = 9$. Tu sais aussi qu'un réel \emph{négatif} n'a pas de racine carrée réelle : aucun réel élevé au carré ne donne $-4$. Mais dans $\mathbb{C}$, tout change : $(2i)^2 = 4i^2 = -4$, donc $-4$ possède des racines carrées dans $\mathbb{C}$, à savoir $2i$ et $-2i$. Mieux encore : \emph{tout} nombre complexe non nul possède des racines carrées dans $\mathbb{C}$. C'est l'une des grandes richesses de cet ensemble.

\begin{definition}[Racine carrée d'un nombre complexe]
Soit $z$ un nombre complexe. On appelle \cle{racine carrée} de $z$ tout nombre complexe $\delta$ tel que :
\[ \delta^2 = z. \]
\end{definition}

\begin{propriete}[Nombre de racines carrées]
Tout nombre complexe \textbf{non nul} admet \textbf{exactement deux} racines carrées dans $\mathbb{C}$, et ces deux racines sont \textbf{opposées} : si $\delta$ est une racine carrée de $z$, alors l'autre est $-\delta$. Le complexe $0$ admet une seule racine carrée : $0$ lui-même.
\end{propriete}

\begin{experience}[Pourquoi les deux racines sont-elles opposées ?]
Raisonnons. Supposons que $\delta$ et $\delta'$ soient deux racines carrées du même complexe non nul $z$. Alors :
\[ \delta^2 = z \quad \text{et} \quad \delta'^2 = z \quad \Longrightarrow \quad \delta^2 = \delta'^2 \quad \Longrightarrow \quad \delta^2 - \delta'^2 = 0. \]
En factorisant : $(\delta - \delta')(\delta + \delta') = 0$. Or dans $\mathbb{C}$, un produit est nul si et seulement si l'un des facteurs est nul. Donc $\delta = \delta'$ ou $\delta = -\delta'$. Autrement dit, à part $\delta$ elle-même, la seule autre racine carrée possible est $-\delta$. Vérifions qu'elle convient : $(-\delta)^2 = \delta^2 = z$. Elle convient bien, et elle est distincte de $\delta$ puisque $z \neq 0$ entraîne $\delta \neq 0$. Nous avons donc exactement deux racines carrées opposées. Leur existence effective sera établie par le calcul dans la méthode qui suit.
\end{experience}

\begin{attention}[On n'écrit pas $\sqrt{z}$ !]
Pour un réel positif, le symbole $\sqrt{a}$ désigne \emph{la} racine carrée \emph{positive} : ce choix est possible car $\mathbb{R}$ est ordonné. Mais $\mathbb{C}$ n'est \textbf{pas ordonné} : il n'y a pas de notion de complexe « positif ». On ne peut donc pas choisir l'une des deux racines de façon naturelle. \textbf{Le symbole $\sqrt{z}$ est interdit pour un complexe quelconque.} On dit : « \emph{les racines carrées de $z$ sont...} ». Écrire $\sqrt{3+4i} = 2+i$ est une faute de langage sanctionnée au Bac.
\end{attention}

\courssub{Rappel : racines carrées d'un nombre réel, vues dans $\mathbb{C}$}

Avant de traiter le cas général, réglons le cas où $z$ est un réel $a$.

\begin{propriete}[Racines carrées d'un réel dans $\mathbb{C}$]
Soit $a$ un nombre réel.
\begin{itemize}
\item Si $a > 0$ : les racines carrées de $a$ sont $\sqrt{a}$ et $-\sqrt{a}$ (réelles).
\item Si $a = 0$ : la seule racine carrée est $0$.
\item Si $a < 0$ : les racines carrées de $a$ sont $i\sqrt{-a}$ et $-i\sqrt{-a}$ (imaginaires pures).
\end{itemize}
\end{propriete}

\begin{exemplebox}[Racines carrées de réels dans $\mathbb{C}$]
\begin{itemize}
\item Racines carrées de $25$ : $5$ et $-5$.
\item Racines carrées de $-9$ : $3i$ et $-3i$, car $(3i)^2 = 9i^2 = -9$.
\item Racines carrées de $-2$ : $i\sqrt{2}$ et $-i\sqrt{2}$.
\end{itemize}
\end{exemplebox}

\courssub{La méthode algébrique : le système fondamental}

Venons-en au cœur de la section. Soit $z = a + ib$ un complexe donné (avec $b \neq 0$, sinon on est ramené au rappel précédent). On cherche $\delta = x + iy$, avec $x$ et $y$ \textbf{réels}, tel que $\delta^2 = z$.

Développons le carré :
\[ \delta^2 = (x + iy)^2 = x^2 + 2ixy + (iy)^2 = (x^2 - y^2) + i(2xy). \]
Par identification des parties réelle et imaginaire avec $z = a + ib$, on obtient :
\[ x^2 - y^2 = a \qquad \text{et} \qquad 2xy = b. \]
Ce système de deux équations à deux inconnues est délicat à résoudre directement. C'est ici qu'intervient \textbf{l'astuce du module}. Si $\delta^2 = z$, alors en passant aux modules : $|\delta^2| = |z|$, c'est-à-dire $|\delta|^2 = |z|$ (propriété du module d'un produit). Or $|\delta|^2 = x^2 + y^2$. D'où une troisième équation :
\[ x^2 + y^2 = |z| = \sqrt{a^2 + b^2}. \]
En \textbf{additionnant} et en \textbf{soustrayant} les équations $x^2 - y^2 = a$ et $x^2 + y^2 = |z|$, on isole immédiatement $x^2$ et $y^2$.

\begin{methode}[Déterminer les racines carrées de $z = a + ib$]
\begin{enumerate}
\item Calculer le module $|z| = \sqrt{a^2 + b^2}$.
\item Poser $\delta = x + iy$ ($x, y$ réels) et écrire le système :
\[ \begin{cases} x^2 - y^2 = a \\ 2xy = b \\ x^2 + y^2 = |z| \end{cases} \]
\item En additionnant et soustrayant la première et la troisième équation :
\[ x^2 = \frac{|z| + a}{2} \qquad \text{et} \qquad y^2 = \frac{|z| - a}{2}. \]
Ces deux quantités sont toujours positives ou nulles (car $|z| \geq |a|$), ce qui garantit l'existence des solutions.
\item Déterminer les valeurs possibles de $x$ (deux signes) et de $y$ (deux signes).
\item \textbf{Utiliser la condition $2xy = b$ pour choisir les signes relatifs} : si $b > 0$, $x$ et $y$ sont de même signe ; si $b < 0$, ils sont de signes contraires. On obtient ainsi exactement deux couples $(x, y)$ opposés.
\item \textbf{Vérifier} en élevant au carré chaque candidat trouvé.
\end{enumerate}
\end{methode}

\begin{attention}[Le piège classique : oublier la condition $2xy = b$]
Les étapes 3 et 4 fournissent \textbf{quatre} candidats : deux valeurs de $x$ et deux valeurs de $y$ se combinent en quatre couples $(x, y)$. Mais seuls \textbf{deux} de ces quatre couples conviennent ! C'est la condition $2xy = b$ qui fait le tri. Par exemple, si $x^2 = 4$ et $y^2 = 1$ avec $b = 4 > 0$, les candidats sont $(2, 1)$, $(2, -1)$, $(-2, 1)$, $(-2, -1)$, mais seuls $(2, 1)$ et $(-2, -1)$ vérifient $2xy > 0$. Oublier cette étape, c'est rendre quatre réponses dont deux fausses : erreur très pénalisée au Bac.
\end{attention}

\courssub{Exemples types du Baccalauréat}

\begin{exoresolu}[Racines carrées de $z = 3 + 4i$]
Déterminer les racines carrées du nombre complexe $z = 3 + 4i$.
\tcblower
\textbf{Solution détaillée.}

\textbf{Étape 1 : module de $z$.}
\[ |z| = \sqrt{3^2 + 4^2} = \sqrt{9 + 16} = \sqrt{25} = 5. \]

\textbf{Étape 2 : le système.} On cherche $\delta = x + iy$ tel que $\delta^2 = z$ :
\[ \begin{cases} x^2 - y^2 = 3 \\ 2xy = 4 \\ x^2 + y^2 = 5 \end{cases} \]

\textbf{Étape 3 : calcul de $x^2$ et $y^2$.} En additionnant puis soustrayant la première et la troisième équation :
\[ x^2 = \frac{5 + 3}{2} = 4 \qquad \text{et} \qquad y^2 = \frac{5 - 3}{2} = 1. \]
Donc $x = 2$ ou $x = -2$, et $y = 1$ ou $y = -1$.

\textbf{Étape 4 : choix des signes.} Comme $2xy = 4 > 0$, $x$ et $y$ sont de \textbf{même signe}. Les couples convenables sont $(2, 1)$ et $(-2, -1)$.

\textbf{Étape 5 : vérification.}
\[ (2 + i)^2 = 4 + 4i + i^2 = 3 + 4i \checkmark \qquad (-2 - i)^2 = (2+i)^2 = 3 + 4i \checkmark \]

\textbf{Conclusion :} les racines carrées de $3 + 4i$ sont $\boxed{2 + i}$ et $\boxed{-2 - i}$.
\end{exoresolu}

\begin{exoresolu}[Racines carrées de $z = -5 + 12i$]
Déterminer les racines carrées de $z = -5 + 12i$.
\tcblower
\textbf{Solution détaillée.}

\textbf{Étape 1 :} $|z| = \sqrt{(-5)^2 + 12^2} = \sqrt{25 + 144} = \sqrt{169} = 13$.

\textbf{Étapes 2 et 3 :} le système $\begin{cases} x^2 - y^2 = -5 \\ x^2 + y^2 = 13 \end{cases}$ donne :
\[ x^2 = \frac{13 - 5}{2} = 4 \qquad \text{et} \qquad y^2 = \frac{13 + 5}{2} = 9. \]
Donc $x = \pm 2$ et $y = \pm 3$.

\textbf{Étape 4 :} $2xy = 12 > 0$, donc $x$ et $y$ de même signe : couples $(2, 3)$ et $(-2, -3)$.

\textbf{Étape 5 : vérification.} $(2 + 3i)^2 = 4 + 12i + 9i^2 = -5 + 12i$ \checkmark

\textbf{Conclusion :} les racines carrées de $-5 + 12i$ sont $\boxed{2 + 3i}$ et $\boxed{-2 - 3i}$.
\end{exoresolu}

\begin{exoresolu}[Racines carrées de $i$]
Déterminer les racines carrées de $z = i$.
\tcblower
\textbf{Solution détaillée.} Ici $a = 0$ et $b = 1$.

\textbf{Étape 1 :} $|i| = \sqrt{0^2 + 1^2} = 1$.

\textbf{Étapes 2 et 3 :}
\[ x^2 = \frac{1 + 0}{2} = \frac{1}{2} \qquad \text{et} \qquad y^2 = \frac{1 - 0}{2} = \frac{1}{2}. \]
Donc $x = \pm \dfrac{1}{\sqrt{2}} = \pm \dfrac{\sqrt{2}}{2}$ et $y = \pm \dfrac{\sqrt{2}}{2}$.

\textbf{Étape 4 :} $2xy = 1 > 0$, donc $x$ et $y$ de même signe.

\textbf{Étape 5 : vérification.}
\[ \left(\frac{\sqrt{2}}{2} + i\frac{\sqrt{2}}{2}\right)^2 = \frac{1}{2} + 2 \times \frac{1}{2}\,i - \frac{1}{2} = i \checkmark \]

\textbf{Conclusion :} les racines carrées de $i$ sont $\boxed{\dfrac{\sqrt{2}}{2} + i\dfrac{\sqrt{2}}{2}}$ et $\boxed{-\dfrac{\sqrt{2}}{2} - i\dfrac{\sqrt{2}}{2}}$. Retiens ce résultat : il revient très souvent au Bac.
\end{exoresolu}

\courssub{Interprétation géométrique}

Plaçons dans le plan complexe les images de $z = 3 + 4i$ et de ses deux racines carrées $\delta = 2 + i$ et $-\delta = -2 - i$. Puisque les deux racines sont opposées, leurs images sont \textbf{symétriques par rapport à l'origine} $O$ : c'est une vérification visuelle immédiate de la cohérence de tes calculs.

\begin{popfigure}
\begin{center}
\begin{tikzpicture}[scale=0.9,>=Stealth]
\draw[popline] (-4.5,0) -- (4.5,0);
\draw[popline] (0,-2.5) -- (0,4.5);
\draw[->,thick,popdark] (-4.2,0) -- (4.4,0) node[below left] {Ré};
\draw[->,thick,popdark] (0,-2.3) -- (0,4.3) node[below left] {Im};
\foreach \x in {-4,-3,-2,-1,1,2,3} \draw (\x,0.08) -- (\x,-0.08) node[below,scale=0.8] {\x};
\foreach \y in {-2,-1,1,2,3} \draw (0.08,\y) -- (-0.08,\y) node[left,scale=0.8] {\y};
\node[below left,scale=0.9] at (0,0) {$O$};
% vecteurs
\draw[->,very thick,popblue] (0,0) -- (3,4);
\draw[->,very thick,popgreen] (0,0) -- (2,1);
\draw[->,very thick,poporange] (0,0) -- (-2,-1);
% points
\fill[popblue] (3,4) circle (2.2pt) node[above right] {$M(z=3+4i)$};
\fill[popgreen] (2,1) circle (2.2pt) node[above right] {$A(\delta=2+i)$};
\fill[poporange] (-2,-1) circle (2.2pt) node[below left] {$A'(-\delta=-2-i)$};
% symétrie
\draw[dashed,popmuted] (-2,-1) -- (2,1);
\end{tikzpicture}
\end{center}
\legende{Images de $z = 3+4i$ et de ses deux racines carrées $\delta = 2+i$ et $-\delta = -2-i$ : les points $A$ et $A'$ sont symétriques par rapport à $O$.}
\end{popfigure}

\begin{aretenir}[L'essentiel de la section]
\begin{itemize}
\item $\delta$ est une \cle{racine carrée} de $z$ si $\delta^2 = z$ ; tout complexe non nul en admet exactement deux, opposées.
\item Pour $z = a + ib$, on résout le système $\begin{cases} x^2 - y^2 = a \\ 2xy = b \\ x^2 + y^2 = |z| \end{cases}$ d'où $x^2 = \dfrac{|z|+a}{2}$ et $y^2 = \dfrac{|z|-a}{2}$.
\item La condition $2xy = b$ fixe les signes relatifs de $x$ et $y$ : même signe si $b > 0$, signes contraires si $b < 0$.
\item On \textbf{vérifie} toujours le résultat en élevant au carré.
\item On n'écrit \textbf{jamais} $\sqrt{z}$ pour un complexe.
\end{itemize}
\end{aretenir}

\begin{savaistu}[Et si on connaissait l'angle ?]
Tu as remarqué que la méthode algébrique, bien que sûre, est un peu lourde : système, astuce du module, discussion des signes. Il existe un point de vue beaucoup plus élégant. Géométriquement, élever un complexe au carré revient à \emph{élever son module au carré} et à \emph{doubler son angle} avec l'axe des réels. Extraire une racine carrée revient donc à prendre la racine du module et à \emph{diviser l'angle par deux} : par exemple, $i$ a pour module $1$ et pour angle $90^{\circ}$, ses racines ont module $1$ et angles $45^{\circ}$ et $225^{\circ}$ --- d'où le $\frac{\sqrt{2}}{2} = \cos 45^{\circ}$ de notre exemple ! C'est la \textbf{forme trigonométrique}, objet de la leçon suivante, qui rendra ces calculs quasi immédiats. Voilà pourquoi le module, que tu as appris à manipuler dans la section précédente, est la clé de voûte de tout ce chapitre : plus tu le maîtrises maintenant, plus la suite sera facile.
\end{savaistu}

\begin{exercice}[Pour t'entraîner]
Déterminer les racines carrées des nombres complexes suivants, puis vérifier par élévation au carré :
\begin{enumerate}
\item $z_1 = 8 - 6i$ ;
\item $z_2 = -3 - 4i$ ;
\item $z_3 = -i$ ;
\item $z_4 = 5 + 12i$.
\end{enumerate}
\end{exercice}

\begin{corrige}[Exercice d'entraînement]
\begin{enumerate}
\item $|z_1| = \sqrt{64+36} = 10$ ; $x^2 = \frac{10+8}{2} = 9$, $y^2 = \frac{10-8}{2} = 1$ ; $2xy = -6 < 0$ donc signes contraires : $\delta = 3 - i$ ou $\delta = -3 + i$. Vérification : $(3-i)^2 = 9 - 6i + i^2 = 8 - 6i$ \checkmark
\item $|z_2| = 5$ ; $x^2 = \frac{5-3}{2} = 1$, $y^2 = \frac{5+3}{2} = 4$ ; $2xy = -4 < 0$ : $\delta = 1 - 2i$ ou $\delta = -1 + 2i$. Vérification : $(1-2i)^2 = 1 - 4i - 4 = -3 - 4i$ \checkmark
\item $|z_3| = 1$ ; $x^2 = y^2 = \frac{1}{2}$ ; $2xy = -1 < 0$ : $\delta = \frac{\sqrt{2}}{2} - i\frac{\sqrt{2}}{2}$ ou $\delta = -\frac{\sqrt{2}}{2} + i\frac{\sqrt{2}}{2}$.
\item $|z_4| = 13$ ; $x^2 = \frac{13+5}{2} = 9$, $y^2 = \frac{13-5}{2} = 4$ ; $2xy = 12 > 0$ : $\delta = 3 + 2i$ ou $\delta = -3 - 2i$. Vérification : $(3+2i)^2 = 9 + 12i - 4 = 5 + 12i$ \checkmark
\end{enumerate}
\end{corrige}

\coursec{Équations du second degré et polynômes de degré 3 dans $\mathbb{C}$}

Dans la section précédente, nous avons appris à extraire les racines carrées d'un nombre complexe. Nous allons maintenant récolter le fruit de tout ce travail : résoudre dans $\mathbb{C}$ les équations polynomiales qui résistaient dans $\mathbb{R}$. Tu sais depuis la classe de Première qu'une équation du second degré à discriminant négatif n'a « pas de solution »... dans $\mathbb{R}$. Cette phrase va enfin trouver son accomplissement : dans $\mathbb{C}$, \emph{toute} équation du second degré à coefficients réels possède des solutions, et nous saurons même traiter les équations de degré $3$. C'est l'un des théorèmes les plus beaux et les plus rentables du programme de Terminale.

\courssub{Équation du second degré à coefficients réels dans $\mathbb{C}$}

\textbf{Intuition et rappel.} Considérons l'équation $z^2 - 2z + 5 = 0$. Son discriminant est $\Delta = 4 - 20 = -16 < 0$. En Première, on concluait : « pas de solution réelle ». Mais nous savons désormais que $-16$ possède des racines carrées dans $\mathbb{C}$ : ce sont $4i$ et $-4i$, puisque $(4i)^2 = 16i^2 = -16$. Tout le mécanisme de la résolution va donc pouvoir se dérouler, à condition de remplacer l'écriture interdite $\sqrt{\Delta}$ par $i\sqrt{-\Delta}$.

\begin{definition}[Discriminant]
Soit l'équation $az^2 + bz + c = 0$ où $a$, $b$, $c$ sont des \cle{réels} avec $a \neq 0$. On appelle \cle{discriminant} de l'équation le réel
\[ \Delta = b^2 - 4ac. \]
\end{definition}

\begin{propriete}[Théorème — Résolution de $az^2+bz+c=0$ à coefficients réels]
Soit $az^2 + bz + c = 0$ avec $a, b, c \in \mathbb{R}$, $a \neq 0$, et $\Delta = b^2 - 4ac$.
\begin{itemize}
\item Si $\Delta > 0$ : l'équation admet \textbf{deux solutions réelles distinctes}
\[ z_1 = \frac{-b - \sqrt{\Delta}}{2a} \quad \text{et} \quad z_2 = \frac{-b + \sqrt{\Delta}}{2a}. \]
\item Si $\Delta = 0$ : l'équation admet \textbf{une solution réelle double} $z_0 = -\dfrac{b}{2a}$.
\item Si $\Delta < 0$ : l'équation admet \textbf{deux solutions complexes conjuguées}
\[ z_1 = \frac{-b - i\sqrt{-\Delta}}{2a} \quad \text{et} \quad z_2 = \frac{-b + i\sqrt{-\Delta}}{2a} = \overline{z_1}. \]
\end{itemize}
\emph{Démonstration exigible au Baccalauréat (cas $\Delta < 0$).}
\end{propriete}

\begin{experience}[Démonstration par la forme canonique]
Partons du membre de gauche et mettons-le sous \cle{forme canonique}, exactement comme en Première :
\begin{align*}
az^2 + bz + c &= a\left(z^2 + \frac{b}{a}z\right) + c = a\left[\left(z + \frac{b}{2a}\right)^2 - \frac{b^2}{4a^2}\right] + c \\
&= a\left[\left(z + \frac{b}{2a}\right)^2 - \frac{b^2 - 4ac}{4a^2}\right] = a\left[\left(z + \frac{b}{2a}\right)^2 - \frac{\Delta}{4a^2}\right].
\end{align*}
L'équation $az^2 + bz + c = 0$ équivaut donc, puisque $a \neq 0$, à
\[ \left(z + \frac{b}{2a}\right)^2 = \frac{\Delta}{4a^2}. \]
\textbf{Cas $\Delta < 0$.} Le point décisif : comme $\Delta < 0$, on a $-\Delta > 0$, et l'on peut écrire
\[ \Delta = -(-\Delta) = i^2 \left(\sqrt{-\Delta}\right)^2 = \left(i\sqrt{-\Delta}\right)^2. \]
L'équation devient alors
\[ \left(z + \frac{b}{2a}\right)^2 = \left(\frac{i\sqrt{-\Delta}}{2a}\right)^2. \]
Un nombre complexe a exactement deux racines carrées opposées, donc
\[ z + \frac{b}{2a} = \frac{i\sqrt{-\Delta}}{2a} \quad \text{ou} \quad z + \frac{b}{2a} = -\frac{i\sqrt{-\Delta}}{2a}, \]
ce qui donne $z_1 = \dfrac{-b - i\sqrt{-\Delta}}{2a}$ et $z_2 = \dfrac{-b + i\sqrt{-\Delta}}{2a}$. Comme $a$, $b$ et $\sqrt{-\Delta}$ sont réels, ces deux nombres sont bien \cle{conjugués} l'un de l'autre : $z_2 = \overline{z_1}$. $\blacksquare$

\textbf{Cas $\Delta > 0$ et $\Delta = 0$.} Ce sont les démonstrations de Première : si $\Delta > 0$, $\left(z + \frac{b}{2a}\right)^2 = \left(\frac{\sqrt{\Delta}}{2a}\right)^2$ donne les deux racines réelles ; si $\Delta = 0$, on obtient $\left(z + \frac{b}{2a}\right)^2 = 0$, d'où la racine double.
\end{experience}

\begin{attention}[Le piège du discriminant négatif]
Quand $\Delta < 0$, il est \textbf{interdit} d'écrire $\sqrt{\Delta}$ : le symbole $\sqrt{\phantom{x}}$ n'est défini que pour les réels \emph{positifs}. La bonne démarche est d'écrire $\Delta = \left(i\sqrt{-\Delta}\right)^2$, c'est-à-dire que « les racines carrées de $\Delta$ » sont $\pm i\sqrt{-\Delta}$. Autre oubli fréquent : ne donner qu'une seule solution. Quand $\Delta < 0$, il y a toujours \textbf{deux} solutions, et elles sont \textbf{conjuguées} : si tu trouves $1 + 2i$, l'autre est obligatoirement $1 - 2i$. Profite-en pour vérifier ton calcul !
\end{attention}

\begin{exoresolu}[Résolution de $z^2 - 2z + 5 = 0$]
Résoudre dans $\mathbb{C}$ l'équation $z^2 - 2z + 5 = 0$.
\tcblower
Les coefficients $a = 1$, $b = -2$, $c = 5$ sont réels. Calculons le discriminant :
\[ \Delta = b^2 - 4ac = (-2)^2 - 4 \times 1 \times 5 = 4 - 20 = -16. \]
Comme $\Delta < 0$, l'équation admet deux solutions complexes conjuguées. On écrit $\Delta = -16 = (4i)^2$, c'est-à-dire $i\sqrt{-\Delta} = i\sqrt{16} = 4i$. D'où :
\[ z_1 = \frac{-b - i\sqrt{-\Delta}}{2a} = \frac{2 - 4i}{2} = 1 - 2i, \qquad z_2 = \overline{z_1} = 1 + 2i. \]
\textbf{Vérification} (toujours profitable) : $(1+2i)^2 - 2(1+2i) + 5 = 1 + 4i + 4i^2 - 2 - 4i + 5 = 1 - 4 - 2 + 5 = 0$. $\checkmark$

Conclusion : $S = \{1 - 2i \; ; \; 1 + 2i\}$.
\end{exoresolu}

\courssub{Somme et produit des racines, factorisation}

\begin{propriete}[Somme et produit des racines]
Si $z_1$ et $z_2$ sont les solutions (réelles ou complexes, distinctes ou confondues) de $az^2 + bz + c = 0$ à coefficients réels, alors
\[ S = z_1 + z_2 = -\frac{b}{a} \qquad \text{et} \qquad P = z_1 \times z_2 = \frac{c}{a}. \]
De plus, dans \textbf{tous les cas} (quel que soit le signe de $\Delta$), on a la \cle{factorisation} :
\[ az^2 + bz + c = a(z - z_1)(z - z_2). \]
\end{propriete}

\textbf{Justification.} D'après la forme canonique, $az^2 + bz + c = a\left(z + \frac{b}{2a} - \frac{\delta}{2a}\right)\left(z + \frac{b}{2a} + \frac{\delta}{2a}\right)$ où $\delta$ désigne une racine carrée de $\Delta$ (réelle si $\Delta \geq 0$, imaginaire pure $i\sqrt{-\Delta}$ si $\Delta < 0$) : c'est exactement $a(z - z_1)(z - z_2)$. En développant $a(z - z_1)(z - z_2) = az^2 - a(z_1+z_2)z + az_1z_2$ et en identifiant avec $az^2 + bz + c$, on obtient $-aS = b$ et $aP = c$, d'où les formules.

\begin{exemplebox}[Vérification sur l'exemple précédent]
Pour $z^2 - 2z + 5 = 0$ de solutions $z_1 = 1 - 2i$ et $z_2 = 1 + 2i$ :
\[ S = (1-2i) + (1+2i) = 2 = -\frac{b}{a} = -\frac{-2}{1}, \qquad P = (1-2i)(1+2i) = 1^2 + 2^2 = 5 = \frac{c}{a}. \checkmark \]
Et l'on peut écrire $z^2 - 2z + 5 = (z - 1 + 2i)(z - 1 - 2i)$. Remarque au passage : le produit de deux complexes conjugués est toujours un réel positif, égal au carré du module.
\end{exemplebox}

\begin{methode}[Retrouver une équation connaissant ses racines]
Pour fabriquer une équation du second degré dont on connaît les racines $z_1$ et $z_2$ :
\begin{enumerate}
\item Calculer $S = z_1 + z_2$ et $P = z_1 z_2$.
\item Écrire l'équation $z^2 - Sz + P = 0$.
\end{enumerate}
Exemple : les racines $3 + i$ et $3 - i$ donnent $S = 6$, $P = 9 + 1 = 10$, donc l'équation $z^2 - 6z + 10 = 0$.
\end{methode}

\courssub{Polynômes de degré 3 : factorisation connaissant une racine}

\textbf{Intuition.} Au Baccalauréat, on te demandera souvent de résoudre une équation de degré $3$ : $P(z) = 0$. L'idée directrice est simple : on ne sait résoudre « en formules » que le degré $2$. Il faut donc \emph{abaisser le degré} : si l'on connaît (ou si l'on devine) une racine $\alpha$, on « sort » le facteur $(z - \alpha)$ et il reste un quotient de degré $2$, que l'on sait traiter.

\begin{propriete}[Factorisation par $(z - \alpha)$]
Soit $P$ un polynôme de degré $3$ et $\alpha$ un nombre complexe. Si $\alpha$ est \cle{racine} de $P$, c'est-à-dire si $P(\alpha) = 0$, alors il existe un polynôme $Q$ de degré $2$ tel que
\[ P(z) = (z - \alpha)\,Q(z) \qquad \text{pour tout } z \in \mathbb{C}. \]
Les solutions de $P(z) = 0$ sont alors $\alpha$ et les solutions de $Q(z) = 0$.
\end{propriete}

\begin{methode}[Déterminer $Q$ : deux techniques au choix]
\textbf{Technique 1 — Coefficients indéterminés (la plus sûre).} Si $P(z) = az^3 + bz^2 + cz + d$ et $P(\alpha) = 0$ :
\begin{enumerate}
\item Poser $P(z) = (z - \alpha)(pz^2 + qz + r)$ où $p$, $q$, $r$ sont à déterminer.
\item Développer : $(z-\alpha)(pz^2 + qz + r) = pz^3 + (q - \alpha p)z^2 + (r - \alpha q)z - \alpha r$.
\item Identifier les coefficients avec ceux de $P$ et résoudre le système obtenu.
\item Vérifier la racine évidente au préalable : calculer $P(\alpha)$ et constater qu'il vaut $0$.
\end{enumerate}
\textbf{Technique 2 — Division euclidienne.} Poser la division de $P(z)$ par $(z - \alpha)$ comme une division de nombres entiers : le reste est nul (car $\alpha$ est racine) et le quotient est $Q(z)$.
\end{methode}

\begin{exoresolu}[Résolution de $z^3 - 2z^2 + z - 2 = 0$]
On considère $P(z) = z^3 - 2z^2 + z - 2$.
\begin{enumerate}
\item Vérifier que $2$ est racine de $P$.
\item Résoudre dans $\mathbb{C}$ l'équation $P(z) = 0$.
\end{enumerate}
\tcblower
\textbf{1.} $P(2) = 8 - 2 \times 4 + 2 - 2 = 8 - 8 + 0 = 0$. Donc $2$ est bien racine de $P$.

\textbf{2.} On peut donc factoriser : $P(z) = (z - 2)(az^2 + bz + c)$. Développons :
\[ (z-2)(az^2 + bz + c) = az^3 + (b - 2a)z^2 + (c - 2b)z - 2c. \]
Par identification avec $z^3 - 2z^2 + z - 2$ :
\[ \begin{cases} a = 1 \\ b - 2a = -2 \\ c - 2b = 1 \\ -2c = -2 \end{cases} \quad \Longrightarrow \quad a = 1,\quad b = 0,\quad c = 1. \]
Donc $P(z) = (z - 2)(z^2 + 1)$. Résolvons $z^2 + 1 = 0$ : $\Delta = -4 = (2i)^2 < 0$, d'où $z = \pm i$ (on reconnaît d'ailleurs directement $z^2 = -1$).

Conclusion : $S = \{2 \; ; \; i \; ; \; -i\}$. On remarque que les deux racines non réelles sont conjuguées : ce n'est pas un hasard, voir la propriété ci-dessous.
\end{exoresolu}

\begin{propriete}[Coefficients réels et racines conjuguées (admise)]
Si $P$ est un polynôme \textbf{à coefficients réels} et si le complexe non réel $z_0$ est racine de $P$, alors $\overline{z_0}$ est aussi racine de $P$. Autrement dit, les racines non réelles d'un polynôme à coefficients réels sont \cle{deux à deux conjuguées}.
\end{propriete}

\textbf{Idée de la preuve.} Elle repose sur les propriétés de la conjugaison : $\overline{z + z'} = \bar z + \bar z'$ et $\overline{zz'} = \bar z \, \bar z'$. Si $P(z) = az^3 + bz^2 + cz + d$ avec $a, b, c, d$ réels, alors
\[ P(\overline{z_0}) = a\overline{z_0}^3 + b\overline{z_0}^2 + c\overline{z_0} + d = \overline{a z_0^3 + b z_0^2 + c z_0 + d} = \overline{P(z_0)} = \overline{0} = 0, \]
car un réel est égal à son conjugué. Donc $\overline{z_0}$ est racine. Conséquence pratique : un polynôme de degré $3$ à coefficients réels possède soit trois racines réelles, soit une racine réelle et deux racines complexes conjuguées.

\begin{popfigure}
\begin{tikzpicture}[
  grow=right, level distance=3.4cm,
  every node/.style={draw, rounded corners=2pt, align=center, font=\small, inner sep=4pt},
  level 1/.style={sibling distance=2.6cm},
  level 2/.style={sibling distance=1.6cm},
  edge from parent/.style={draw=popblue, thick, -{Stealth}}
]
\node[fill=popblueL, align=center] {$P(z) = 0$\\ (degré 3)}
  child { node[fill=poporangeL, align=center] {Racine évidente $\alpha$\\ vérifier $P(\alpha)=0$}
    child { node[fill=popgreenL, align=center] {$P(z)=(z-\alpha)Q(z)$\\ coefficients indéterminés}
      child { node[fill=poppurpleL] {$\Delta$ de $Q$}
        child { node[fill=popgoldL] {$\Delta \geq 0$ : racines réelles} }
        child { node[fill=poppinkL] {$\Delta < 0$ : racines conjuguées} }
      }
    }
  };
\end{tikzpicture}
\legende{Stratégie générale de résolution d'une équation de degré 3 : racine évidente, factorisation, puis discriminant du quotient.}
\end{popfigure}

\begin{savaistu}[Les racines cubiques de l'unité]
L'équation $z^3 = 1$, c'est-à-dire $z^3 - 1 = 0$, se résout par notre méthode : $1$ est racine évidente, et $z^3 - 1 = (z - 1)(z^2 + z + 1)$. Le discriminant du quotient est $\Delta = 1 - 4 = -3$, d'où deux racines conjuguées notées traditionnellement
\[ j = -\frac{1}{2} + i\frac{\sqrt{3}}{2} \qquad \text{et} \qquad j^2 = \overline{j} = -\frac{1}{2} - i\frac{\sqrt{3}}{2}. \]
Ces trois nombres $1$, $j$, $j^2$ sont les \cle{racines cubiques de l'unité} ; leurs images forment un triangle équilatéral inscrit dans le cercle trigonométrique. Elles interviennent partout : en électrotechnique (systèmes triphasés qui alimentent les quartiers de Douala et Yaoundé), en traitement du signal, et dans la célèbre formule de Cardan pour les équations de degré 3.
\end{savaistu}

\begin{popfigure}
\begin{tikzpicture}[scale=2.6]
  \draw[popline, -{Stealth}] (-1.5,0) -- (1.6,0) node[below left] {\small Réel};
  \draw[popline, -{Stealth}] (0,-1.4) -- (0,1.5) node[below left] {\small Imaginaire};
  \draw[popmuted, dashed] (0,0) circle (1);
  \coordinate (A) at (1,0);
  \coordinate (B) at (-0.5,0.866);
  \coordinate (C) at (-0.5,-0.866);
  \draw[thick, popblue, fill=popblueL, fill opacity=0.35] (A) -- (B) -- (C) -- cycle;
  \fill[popink] (A) circle (1.4pt) node[above right] {$1$};
  \fill[poporange] (B) circle (1.4pt) node[above left] {$j = -\frac12 + i\frac{\sqrt3}{2}$};
  \fill[poppurple] (C) circle (1.4pt) node[below left] {$j^2 = \bar j = -\frac12 - i\frac{\sqrt3}{2}$};
  \fill[popdark] (0,0) circle (1pt) node[below right] {\small $O$};
\end{tikzpicture}
\legende{Les trois racines de $z^3 - 1 = 0$ dans le plan complexe : $1$, $j$ et $j^2 = \bar j$, sommets d'un triangle équilatéral inscrit dans le cercle unité.}
\end{popfigure}

\begin{attention}[Conditions et rédaction]
\begin{itemize}
\item Le théorème « $\Delta < 0 \Rightarrow$ racines conjuguées » et la propriété « racines non réelles deux à deux conjuguées » exigent des \textbf{coefficients réels}. Pour une équation à coefficients complexes (hors programme mais parfois évoquée), ces résultats sont faux.
\item Avant de factoriser par $(z - \alpha)$, \textbf{montre explicitement} que $P(\alpha) = 0$ : c'est ce calcul qui justifie l'existence de la factorisation, et le correcteur du Bac l'attend.
\item Dans la méthode des coefficients indéterminés, n'oublie pas que le système comporte autant d'équations que de coefficients (4 équations pour un degré 3) : utilise-les toutes pour vérifier la cohérence.
\end{itemize}
\end{attention}

\begin{exercice}
\begin{enumerate}
\item Résoudre dans $\mathbb{C}$ : a) $z^2 + 4z + 13 = 0$ ; b) $2z^2 - 2z + 5 = 0$.
\item Déterminer une équation du second degré à coefficients réels dont une racine est $2 - 3i$.
\item Soit $P(z) = z^3 - z^2 + 4z - 4$. Vérifier que $1$ est racine de $P$, puis résoudre $P(z) = 0$ dans $\mathbb{C}$.
\end{enumerate}
\end{exercice}

\begin{corrige}[Exercice]
\textbf{1. a)} $\Delta = 16 - 52 = -36 = (6i)^2 < 0$ ; $z_1 = \dfrac{-4 - 6i}{2} = -2 - 3i$, $z_2 = -2 + 3i$. $S = \{-2 - 3i \,;\, -2 + 3i\}$.

\textbf{1. b)} $\Delta = 4 - 40 = -36 = (6i)^2$ ; $z_1 = \dfrac{2 - 6i}{4} = \dfrac{1}{2} - \dfrac{3}{2}i$, $z_2 = \dfrac{1}{2} + \dfrac{3}{2}i$. $S = \left\{\dfrac{1 - 3i}{2} \,;\, \dfrac{1 + 3i}{2}\right\}$.

\textbf{2.} Les coefficients étant réels, l'autre racine est $\overline{2 - 3i} = 2 + 3i$. $S = 4$, $P = (2)^2 + (3)^2 = 13$. Équation : $z^2 - 4z + 13 = 0$.

\textbf{3.} $P(1) = 1 - 1 + 4 - 4 = 0$ : $1$ est racine. On pose $P(z) = (z - 1)(z^2 + bz + c)$ ; en développant et identifiant : $z^3 + (b-1)z^2 + (c-b)z - c = z^3 - z^2 + 4z - 4$, d'où $b = 0$, $c = 4$. Donc $P(z) = (z-1)(z^2 + 4)$. Or $z^2 + 4 = 0 \Leftrightarrow z = \pm 2i$. Conclusion : $S = \{1 \,;\, 2i \,;\, -2i\}$.
\end{corrige}

\begin{aretenir}[L'essentiel de la section]
\begin{itemize}
\item Pour $az^2 + bz + c = 0$ à coefficients réels : si $\Delta < 0$, écrire $\Delta = \left(i\sqrt{-\Delta}\right)^2$ ; les deux solutions sont $z_{1,2} = \dfrac{-b \pm i\sqrt{-\Delta}}{2a}$, \textbf{conjuguées}.
\item Dans tous les cas : $z_1 + z_2 = -\dfrac{b}{a}$, $z_1 z_2 = \dfrac{c}{a}$, et $az^2 + bz + c = a(z - z_1)(z - z_2)$.
\item Degré 3 : vérifier la racine évidente $\alpha$, factoriser $P(z) = (z - \alpha)Q(z)$ par coefficients indéterminés ou division euclidienne, puis résoudre $Q(z) = 0$.
\item Coefficients réels : les racines non réelles vont par paires conjuguées.
\end{itemize}
\end{aretenir}

\newpage

\coursec{Activité d'intégration}

\coursec{Nombres complexes : forme algébrique}

\courssub{Objectifs de l'activité}

À l'issue de cette activité d'intégration, tu devras être capable de mobiliser \emph{simultanément} tous les savoir-faire de la leçon dans une situation concrète. Plus précisément, tu sauras :

\begin{itemize}
\item additionner, multiplier et diviser des nombres complexes et réduire le résultat sous \cle{forme algébrique} $a + ib$ ;
\item utiliser le \cle{conjugué} pour rendre un dénominateur réel et vérifier qu'un produit $z\bar{z}$ est un réel positif ;
\item calculer le \cle{module} d'un nombre complexe et l'interpréter dans un contexte physique ;
\item déterminer les \cle{racines carrées} d'un nombre complexe donné sous forme algébrique ;
\item résoudre dans $\mathbb{C}$ une \cle{équation du second degré} à coefficients réels de discriminant négatif ;
\item \cle{factoriser} un polynôme de degré 3 connaissant une racine évidente, puis résoudre l'équation associée ;
\item placer des nombres complexes dans le \cle{plan complexe} et observer géométriquement la conjugaison.
\end{itemize}

\courssub{Situation}

Dans le quartier Nkolbisson à Yaoundé, les coupures d'électricité sont fréquentes. L'entreprise ENEO dépêche un technicien, M. Etoundi, pour diagnostiquer une installation défectueuse alimentant un groupe de maisons. En électricité en courant alternatif, chaque élément d'un circuit (résistance, bobine, condensateur) est modélisé par une \emph{impédance}, qui n'est pas un nombre réel mais un \textbf{nombre complexe}, exprimé en ohms ($\Omega$).

M. Etoundi modélise deux branches du circuit du quartier par les impédances complexes :
\[ Z_1 = 3 + 4i \qquad \text{et} \qquad Z_2 = 1 - 2i \quad (\text{en } \Omega). \]

Il explique à son apprenti, un jeune bachelier passionné de mathématiques, les règles suivantes issues de la physique :

\begin{itemize}
\item Quand deux impédances sont montées \textbf{en série}, l'impédance totale est la somme : $Z_s = Z_1 + Z_2$.
\item Quand elles sont montées \textbf{en dérivation} (en parallèle), l'impédance équivalente est :
\[ Z_d = \dfrac{Z_1 Z_2}{Z_1 + Z_2}. \]
\item Le produit $Z_1 \overline{Z_1}$ donne le carré du module : il est toujours réel positif, et $|Z_1|$ s'interprète comme la \og résistance apparente \fg{} de la branche, celle que mesure l'appareil du technicien.
\item Pour un calcul de courant dans une bobine, M. Etoundi a besoin d'un nombre complexe $w$ tel que $w^2 = Z_1$.
\item Enfin, la fréquence de résonance du circuit est liée aux solutions de l'équation $z^2 - 2z + 10 = 0$, et la stabilité de l'ensemble du réseau du quartier dépend des solutions de l'équation $z^3 - 3z^2 + 4z - 2 = 0$, dont il a remarqué que $1$ est une solution évidente.
\end{itemize}

L'apprenti sort son cahier de mathématiques de Terminale : toutes ces questions relèvent de la leçon sur les nombres complexes. Par groupes de quatre élèves, aide M. Etoundi à mener tous les calculs, puis présente un rapport collectif au tableau.

\courssub{Consignes}

\textbf{Partie A : Impédances du circuit (opérations dans $\mathbb{C}$)}
\begin{enumerate}
\item Calculer l'impédance totale en série $Z_s = Z_1 + Z_2$ sous forme algébrique.
\item Calculer le produit $Z_1 Z_2$ sous forme algébrique.
\item En déduire l'impédance en dérivation $Z_d = \dfrac{Z_1 Z_2}{Z_1 + Z_2}$ sous forme algébrique. On pensera à multiplier le numérateur et le dénominateur par le conjugué du dénominateur.
\end{enumerate}

\textbf{Partie B : Conjugué et module}
\begin{enumerate}
\setcounter{enumi}{3}
\item Écrire le conjugué $\overline{Z_1}$, puis calculer le produit $Z_1 \overline{Z_1}$. Que constate-t-on ?
\item Calculer les modules $|Z_1|$ et $|Z_2|$. Quelle branche du circuit présente la plus grande \og résistance apparente \fg{} ?
\item Vérifier que $|Z_1 Z_2| = |Z_1| \times |Z_2|$ à l'aide des résultats précédents.
\end{enumerate}

\textbf{Partie C : Racines carrées de $Z_1$}
\begin{enumerate}
\setcounter{enumi}{6}
\item On cherche un nombre complexe $w = a + ib$ (avec $a$ et $b$ réels) tel que $w^2 = Z_1 = 3 + 4i$.
\begin{enumerate}
\item Montrer que $a$ et $b$ vérifient le système : $a^2 - b^2 = 3$ et $2ab = 4$.
\item En utilisant $|w^2| = |w|^2 = |Z_1|$, montrer que $a^2 + b^2 = 5$.
\item En déduire les valeurs de $a^2$ et $b^2$, puis les deux racines carrées de $Z_1$. Vérifier par le calcul.
\end{enumerate}
\end{enumerate}

\textbf{Partie D : Équations du réseau}
\begin{enumerate}
\setcounter{enumi}{7}
\item Résoudre dans $\mathbb{C}$ l'équation de résonance $z^2 - 2z + 10 = 0$.
\item Résolution de l'équation de stabilité $z^3 - 3z^2 + 4z - 2 = 0$.
\begin{enumerate}
\item Vérifier que $1$ est solution.
\item Déterminer les réels $a$, $b$, $c$ tels que $z^3 - 3z^2 + 4z - 2 = (z - 1)(az^2 + bz + c)$.
\item Achever la résolution dans $\mathbb{C}$.
\end{enumerate}
\item Placer dans le plan complexe muni d'un repère orthonormé $(O\,;\vec{u},\vec{v})$ les points images de toutes les solutions trouvées aux questions 8 et 9. Que remarque-t-on concernant les solutions non réelles de chaque équation ?
\end{enumerate}

\textbf{Bilan collectif :} chaque groupe présente une question au tableau ; la classe vérifie ensemble que les racines non réelles de chaque équation à coefficients réels sont deux à deux conjuguées, et discute de l'interprétation géométrique de la conjugaison (symétrie par rapport à l'axe des réels).

\courssub{Ressources à mobiliser}

\begin{aretenir}[Boîte à outils de la leçon]
\begin{itemize}
\item \textbf{Opérations} : $(a+ib) + (a'+ib') = (a+a') + i(b+b')$ ; $(a+ib)(a'+ib') = (aa' - bb') + i(ab' + a'b)$, en utilisant $i^2 = -1$.
\item \textbf{Quotient} : pour $z_2 \neq 0$, $\dfrac{z_1}{z_2} = \dfrac{z_1 \overline{z_2}}{z_2 \overline{z_2}} = \dfrac{z_1 \overline{z_2}}{|z_2|^2}$ : on multiplie par le conjugué du dénominateur.
\item \textbf{Conjugué} : si $z = a + ib$, alors $\overline{z} = a - ib$ et $z\overline{z} = a^2 + b^2 = |z|^2$.
\item \textbf{Module} : $|a + ib| = \sqrt{a^2 + b^2}$ ; propriétés : $|z_1 z_2| = |z_1||z_2|$ et $\left|\dfrac{z_1}{z_2}\right| = \dfrac{|z_1|}{|z_2|}$.
\item \textbf{Racines carrées de $z = a+ib$} : on pose $w = x + iy$, on traduit $w^2 = z$ par le système $x^2 - y^2 = a$, $2xy = b$, et on ajoute $x^2 + y^2 = |z|$.
\item \textbf{Second degré à coefficients réels} : si $\Delta < 0$, l'équation $az^2 + bz + c = 0$ admet deux solutions complexes conjuguées :
\[ z_1 = \dfrac{-b - i\sqrt{-\Delta}}{2a} \qquad \text{et} \qquad z_2 = \dfrac{-b + i\sqrt{-\Delta}}{2a} = \overline{z_1}. \]
\item \textbf{Factorisation de degré 3} : si $\alpha$ est racine de $P$, alors $P(z) = (z - \alpha)Q(z)$ avec $Q$ de degré 2, obtenu par identification des coefficients ou par division euclidienne.
\end{itemize}
\end{aretenir}

\courssub{Démarche et corrigé commenté}

\begin{exoresolu}[Rapport de l'apprenti de M. Etoundi : corrigé complet]
On donne $Z_1 = 3 + 4i$ et $Z_2 = 1 - 2i$. Il s'agit de répondre à toutes les consignes des parties A à D.
\tcblower
\textbf{Partie A : impédances du circuit.}

\textbf{1.} La somme s'effectue partie réelle avec partie réelle, partie imaginaire avec partie imaginaire :
\[ Z_s = Z_1 + Z_2 = (3 + 1) + i(4 - 2) = 4 + 2i \quad \SI{}{\ohm}. \]

\textbf{2.} Pour le produit, on développe et on remplace $i^2$ par $-1$ :
\begin{align*}
Z_1 Z_2 &= (3 + 4i)(1 - 2i) = 3 - 6i + 4i - 8i^2 \\
&= 3 - 2i + 8 = 11 - 2i \quad \SI{}{\ohm}.
\end{align*}

\textbf{3.} Pour le quotient, on multiplie numérateur et dénominateur par le conjugué du dénominateur, $\overline{4 + 2i} = 4 - 2i$ :
\[ Z_d = \dfrac{11 - 2i}{4 + 2i} = \dfrac{(11 - 2i)(4 - 2i)}{(4 + 2i)(4 - 2i)}. \]
Le dénominateur vaut $4^2 + 2^2 = 20$. Le numérateur vaut :
\[ (11 - 2i)(4 - 2i) = 44 - 22i - 8i + 4i^2 = 44 - 30i - 4 = 40 - 30i. \]
Donc $Z_d = \dfrac{40 - 30i}{20} = 2 - {1,5}i \ \SI{}{\ohm}$.

\textbf{Partie B : conjugué et module.}

\textbf{4.} $\overline{Z_1} = 3 - 4i$, et :
\[ Z_1 \overline{Z_1} = (3 + 4i)(3 - 4i) = 9 - 16i^2 = 9 + 16 = 25. \]
On constate que ce produit est un \textbf{réel positif} : c'est toujours le cas, car $z\overline{z} = |z|^2$.

\textbf{5.} $|Z_1| = \sqrt{3^2 + 4^2} = \sqrt{25} = 5 \ \SI{}{\ohm}$ et $|Z_2| = \sqrt{1^2 + (-2)^2} = \sqrt{5} \approx 2,24 \ \SI{}{\ohm}$. La branche 1 présente la plus grande résistance apparente.

\textbf{6.} D'une part $|Z_1 Z_2| = |11 - 2i| = \sqrt{121 + 4} = \sqrt{125} = 5\sqrt{5}$. D'autre part $|Z_1| \times |Z_2| = 5\sqrt{5}$. La propriété $|z_1 z_2| = |z_1||z_2|$ est bien vérifiée.

\textbf{Partie C : racines carrées de $Z_1 = 3 + 4i$.}

\textbf{7. a)} Soit $w = a + ib$ tel que $w^2 = 3 + 4i$. Or $w^2 = a^2 - b^2 + 2iab$. Par identification des parties réelles et imaginaires : $a^2 - b^2 = 3$ et $2ab = 4$.

\textbf{b)} De $|w^2| = |w|^2$ et $|w^2| = |Z_1| = 5$, on tire $a^2 + b^2 = 5$.

\textbf{c)} En additionnant : $2a^2 = 8$, donc $a^2 = 4$ ; en soustrayant : $2b^2 = 2$, donc $b^2 = 1$. Comme $2ab = 4 > 0$, $a$ et $b$ sont de même signe : $(a, b) = (2, 1)$ ou $(-2, -1)$. Les deux racines carrées sont :
\[ w_1 = 2 + i \qquad \text{et} \qquad w_2 = -2 - i. \]
Vérification : $(2 + i)^2 = 4 + 4i + i^2 = 3 + 4i = Z_1$. \checkmark

\textbf{Partie D : équations du réseau.}

\textbf{8.} Pour $z^2 - 2z + 10 = 0$ : $\Delta = (-2)^2 - 4 \times 10 = -36 < 0$. L'équation admet deux solutions conjuguées :
\[ z_1 = \dfrac{2 - 6i}{2} = 1 - 3i \qquad \text{et} \qquad z_2 = 1 + 3i. \]

\textbf{9. a)} $P(1) = 1 - 3 + 4 - 2 = 0$ : $1$ est bien racine.

\textbf{b)} On pose $P(z) = (z - 1)(az^2 + bz + c) = az^3 + (b - a)z^2 + (c - b)z - c$. Par identification avec $z^3 - 3z^2 + 4z - 2$ : $a = 1$, $-c = -2$ donc $c = 2$, $b - a = -3$ donc $b = -2$. Ainsi :
\[ z^3 - 3z^2 + 4z - 2 = (z - 1)(z^2 - 2z + 2). \]

\textbf{c)} Le discriminant de $z^2 - 2z + 2$ vaut $4 - 8 = -4 < 0$, d'où $z = \dfrac{2 \pm 2i}{2} = 1 \pm i$. L'ensemble des solutions est :
\[ \mathcal{S} = \{1 \ ;\ 1 - i \ ;\ 1 + i\}. \]

\textbf{10.} On place les points $A(1 - 3i)$, $B(1 + 3i)$, $C(1)$, $D(1 - i)$, $E(1 + i)$. Les points $A$ et $B$ sont symétriques par rapport à l'axe des réels, de même que $D$ et $E$ : les racines non réelles sont deux à deux \textbf{conjuguées}.
\end{exoresolu}

\begin{popfigure}
\begin{center}
\begin{tikzpicture}[scale=1.1]
\draw[->,popline] (-1.5,0) -- (3,0) node[below right] {Re};
\draw[->,popline] (0,-3.5) -- (0,3.5) node[above left] {Im};
\foreach \y in {-3,-2,-1,1,2,3} \draw[popline] (-0.06,\y) -- (0.06,\y);
\foreach \x in {1,2} \draw[popline] (\x,-0.06) -- (\x,0.06);
\fill[popink] (1,-3) circle (2.2pt) node[right] {$A(1-3i)$};
\fill[popink] (1,3) circle (2.2pt) node[right] {$B(1+3i)$};
\fill[poporange] (1,0) circle (2.2pt) node[above right] {$C(1)$};
\fill[popgreen] (1,-1) circle (2.2pt) node[right] {$D(1-i)$};
\fill[popgreen] (1,1) circle (2.2pt) node[right] {$E(1+i)$};
\draw[dashed,popblue] (1,-3) -- (1,3);
\draw[dashed,popblue] (1,-1) -- (1,1);
\end{tikzpicture}
\end{center}
\legende{Les solutions des deux équations : les racines non réelles sont symétriques par rapport à l'axe des réels (conjugaison).}
\end{popfigure}

\courssub{Bilan de l'activité}

Cette activité a permis de mettre en évidence les points essentiels suivants :

\begin{itemize}
\item Les quatre opérations dans $\mathbb{C}$ se ramènent toujours à des calculs dans $\mathbb{R}$, grâce à la règle $i^2 = -1$ et à la technique du \textbf{conjugué} pour les quotients : l'impédance en dérivation, pourtant donnée par une fraction de complexes, s'écrit simplement $Z_d = 2 - {1,5}i$.
\item Le produit $z\overline{z} = |z|^2$ est un réel positif : c'est ce qui permet au technicien de relier le formalisme complexe à une grandeur mesurable, la résistance apparente $|Z_1| = 5\ \SI{}{\ohm}$.
\item La recherche de racines carrées illustre la puissance de la méthode mixte : identification des parties réelle et imaginaire, complétée par la condition sur les modules $|w|^2 = |z|$.
\item Les équations à coefficients réels de discriminant négatif admettent toujours deux solutions \textbf{conjuguées l'une de l'autre}, ce qui se traduit géométriquement par une symétrie des points images par rapport à l'axe des réels. Cette propriété, observée ici sur les deux équations du réseau, est un fait général que tu retrouveras dans toute la suite du cours sur les nombres complexes.
\item Enfin, la factorisation d'un polynôme de degré 3 à partir d'une racine évidente transforme un problème de degré 3 en un problème de degré 2, déjà maîtrisé : c'est une stratégie de réduction qu'il faut retenir pour le Baccalauréat.
\end{itemize}

\begin{savaistu}[Pourquoi les électriciens aiment-ils les complexes ?]
La notation $a + ib$ utilisée en électricité a été introduite à la fin du XIX\textsuperscript{e} siècle par l'ingénieur Charles Steinmetz pour étudier les courants alternatifs : les tensions et courants sinusoïdaux deviennent de simples nombres complexes, et les lois d'association des circuits (série, dérivation) se calculent alors comme avec des résistances ordinaires. Les ingénieurs d'ENEO, comme ceux de tous les réseaux électriques du monde, utilisent quotidiennement ce formalisme. Les mathématiques de Terminale éclairent littéralement le quartier !
\end{savaistu}

\newpage

\coursec{Méthodes et exercices résolus}

\courssub{Mettre un nombre complexe sous forme algébrique}

\begin{methode}[Forme algébrique d'un nombre complexe]
Pour écrire un nombre complexe $z$ sous la forme $a + ib$ avec $a \in \mathbb{R}$ et $b \in \mathbb{R}$ :
\begin{enumerate}
\item Remplacer systématiquement $i^2$ par $-1$ dès qu'il apparaît (et $i^3 = -i$, $i^4 = 1$).
\item Développer les produits comme dans $\mathbb{R}$ (distributivité, identités remarquables), puis regrouper les termes réels et les termes en $i$.
\item Pour un \cle{quotient}, multiplier le numérateur et le dénominateur par le \cle{conjugué} du dénominateur : si $z = \dfrac{a+ib}{c+id}$ avec $(c+id) \neq 0$, alors
\[ z = \frac{(a+ib)(c-id)}{(c+id)(c-id)} = \frac{(a+ib)(c-id)}{c^2+d^2}. \]
\item Conclure en identifiant : $\operatorname{Re}(z) = a$ (partie réelle) et $\operatorname{Im}(z) = b$ (partie imaginaire).
\end{enumerate}
Réflexe : un dénominateur contenant $i$ doit \emph{toujours} disparaître avant de conclure.
\end{methode}

\begin{exoresolu}[Opérations et forme algébrique]
On donne $z_1 = 2 - 3i$ et $z_2 = 1 + 2i$.
\begin{enumerate}
\item Déterminer la forme algébrique de $z_1 + z_2$, $z_1 \times z_2$ et $z_1^2$.
\item Déterminer la forme algébrique de $Z = \dfrac{z_1}{z_2}$.
\end{enumerate}
\tcblower
\textbf{1. Somme.} On regroupe parties réelles et imaginaires :
\[ z_1 + z_2 = (2 - 3i) + (1 + 2i) = (2+1) + (-3+2)i = 3 - i. \]

\textbf{Produit.} On développe et on remplace $i^2$ par $-1$ :
\begin{align*}
z_1 \times z_2 &= (2-3i)(1+2i) = 2 + 4i - 3i - 6i^2 \\
&= 2 + i - 6\times(-1) = 2 + i + 6 = 8 + i.
\end{align*}

\textbf{Carré.} Avec l'identité remarquable $(a-b)^2 = a^2 - 2ab + b^2$ :
\[ z_1^2 = (2-3i)^2 = 4 - 12i + 9i^2 = 4 - 12i - 9 = -5 - 12i. \]

\textbf{2. Quotient.} Le conjugué du dénominateur $z_2 = 1+2i$ est $\overline{z_2} = 1 - 2i$. On multiplie numérateur et dénominateur par $1-2i$ :
\[ Z = \frac{2-3i}{1+2i} = \frac{(2-3i)(1-2i)}{(1+2i)(1-2i)}. \]
Le dénominateur vaut $1^2 + 2^2 = 1 + 4 = 5$ (identité $(c+id)(c-id) = c^2+d^2$). Le numérateur :
\[ (2-3i)(1-2i) = 2 - 4i - 3i + 6i^2 = 2 - 7i - 6 = -4 - 7i. \]
Donc
\[ Z = \frac{-4-7i}{5} = -\frac{4}{5} - \frac{7}{5}i. \]
\textbf{Conclusion :} $z_1 + z_2 = 3 - i$ ; $z_1 z_2 = 8 + i$ ; $z_1^2 = -5 - 12i$ ; $Z = -\dfrac{4}{5} - \dfrac{7}{5}i$.
\end{exoresolu}

\courssub{Calculer le conjugué et le module d'un complexe}

\begin{methode}[Conjugué et module : les réflexes]
Soit $z = a + ib$ avec $(a,b) \in \mathbb{R}^2$.
\begin{enumerate}
\item \cle{Conjugué} : $\bar{z} = a - ib$ (on change le signe de la partie imaginaire uniquement).
\item \cle{Module} : $|z| = \sqrt{a^2 + b^2}$ (toujours un réel positif ou nul ; $|z| = 0 \iff z = 0$).
\item Formule pivot : $z\bar{z} = |z|^2 = a^2 + b^2$. C'est elle qui justifie la technique du quotient.
\item Propriétés à utiliser pour aller vite (sans tout recalculer) :
\[ \overline{z + z'} = \bar{z} + \bar{z'}, \quad \overline{zz'} = \bar{z}\,\bar{z'}, \quad \overline{\left(\tfrac{z}{z'}\right)} = \frac{\bar{z}}{\bar{z'}} \ (z' \neq 0), \quad |zz'| = |z|\,|z'|, \quad \left|\tfrac{z}{z'}\right| = \frac{|z|}{|z'|} \ (z' \neq 0). \]
\item Caractérisations utiles : $z$ réel $\iff z = \bar{z}$ ; $z$ imaginaire pur $\iff z = -\bar{z}$.
\end{enumerate}
\end{methode}

\begin{exoresolu}[Conjugué, module et nature d'un complexe]
On pose $z = \dfrac{3+i}{1-i}$.
\begin{enumerate}
\item Déterminer la forme algébrique de $z$, puis $\bar{z}$.
\item Calculer $|z|$ de deux façons : directement, puis avec les propriétés du module.
\item Montrer que $w = z + \bar{z}$ est un nombre réel et le calculer.
\end{enumerate}
\tcblower
\textbf{1. Forme algébrique.} On multiplie par le conjugué du dénominateur, $1+i$ :
\[ z = \frac{(3+i)(1+i)}{(1-i)(1+i)} = \frac{3 + 3i + i + i^2}{1^2 + 1^2} = \frac{3 + 4i - 1}{2} = \frac{2 + 4i}{2} = 1 + 2i. \]
Donc $\bar{z} = 1 - 2i$.

\textbf{2. Module, méthode directe :}
\[ |z| = \sqrt{1^2 + 2^2} = \sqrt{5}. \]
\textbf{Avec les propriétés :} $\left|\dfrac{3+i}{1-i}\right| = \dfrac{|3+i|}{|1-i|} = \dfrac{\sqrt{9+1}}{\sqrt{1+1}} = \dfrac{\sqrt{10}}{\sqrt{2}} = \sqrt{\dfrac{10}{2}} = \sqrt{5}$. Les deux méthodes concordent.

\textbf{3. Nature de $w$.} Pour tout complexe, $z + \bar{z} = 2\operatorname{Re}(z)$, qui est réel. Ici :
\[ w = (1+2i) + (1-2i) = 2. \]
Vérification par la caractérisation : $\overline{w} = \bar{z} + \overline{\bar{z}} = \bar{z} + z = w$, donc $w$ est réel.
\textbf{Conclusion :} $z = 1+2i$, $\bar{z} = 1 - 2i$, $|z| = \sqrt{5}$ et $w = 2 \in \mathbb{R}$.
\end{exoresolu}

\courssub{Déterminer les racines carrées d'un nombre complexe}

\begin{methode}[Racines carrées sous forme algébrique]
Pour trouver les complexes $z = x + iy$ (avec $x, y$ réels) tels que $z^2 = a + ib$ :
\begin{enumerate}
\item Écrire $z^2 = (x^2 - y^2) + 2ixy$ et identifier parties réelle et imaginaire :
\[ x^2 - y^2 = a \qquad \text{et} \qquad 2xy = b. \]
\item Ajouter l'équation des modules : $|z|^2 = |a+ib|$, c'est-à-dire $x^2 + y^2 = \sqrt{a^2+b^2}$.
\item Résoudre le système en $x^2$ et $y^2$ par addition et soustraction :
\[ x^2 = \frac{a + \sqrt{a^2+b^2}}{2}, \qquad y^2 = \frac{-a + \sqrt{a^2+b^2}}{2}. \]
\item Le signe de $b$ fixe les signes relatifs de $x$ et $y$ grâce à $2xy = b$ : si $b > 0$, $x$ et $y$ ont le même signe ; si $b < 0$, ils sont de signes contraires.
\item Conclure : il y a exactement \cle{deux racines carrées opposées} $z_0$ et $-z_0$.
\end{enumerate}
\end{methode}

\begin{exoresolu}[Racines carrées de $3 - 4i$]
Déterminer les racines carrées du nombre complexe $Z = 3 - 4i$.
\tcblower
On cherche $z = x + iy$, avec $x \in \mathbb{R}$ et $y \in \mathbb{R}$, tel que $z^2 = 3 - 4i$.

\textbf{Étape 1 : identification.} $z^2 = x^2 - y^2 + 2ixy$, donc :
\[ x^2 - y^2 = 3 \quad (1) \qquad \text{et} \qquad 2xy = -4 \quad (2). \]

\textbf{Étape 2 : équation des modules.} $|z^2| = |z|^2 = x^2 + y^2$ et $|3-4i| = \sqrt{9+16} = \sqrt{25} = 5$, donc :
\[ x^2 + y^2 = 5 \quad (3). \]

\textbf{Étape 3 : résolution.} En additionnant (1) et (3) : $2x^2 = 8$, soit $x^2 = 4$, donc $x = 2$ ou $x = -2$. En soustrayant (1) de (3) : $2y^2 = 2$, soit $y^2 = 1$, donc $y = 1$ ou $y = -1$.

\textbf{Étape 4 : signes.} D'après (2), $2xy = -4 < 0$, donc $x$ et $y$ sont de \emph{signes contraires}. Les couples possibles sont $(x,y) = (2,-1)$ et $(x,y) = (-2,1)$.

\textbf{Vérification :} $(2-i)^2 = 4 - 4i + i^2 = 3 - 4i$. C'est correct.

\textbf{Conclusion :} les racines carrées de $3 - 4i$ sont $z_1 = 2 - i$ et $z_2 = -2 + i$.
\end{exoresolu}

\courssub{Résoudre une équation du second degré à coefficients réels dans $\mathbb{C}$}

\begin{methode}[Second degré à coefficients réels : discriminant négatif]
Pour résoudre $az^2 + bz + c = 0$ avec $a, b, c$ réels, $a \neq 0$ :
\begin{enumerate}
\item Calculer le discriminant $\Delta = b^2 - 4ac$.
\item Si $\Delta > 0$ : deux solutions réelles $z = \dfrac{-b \pm \sqrt{\Delta}}{2a}$.
\item Si $\Delta = 0$ : une solution réelle double $z_0 = \dfrac{-b}{2a}$.
\item Si $\Delta < 0$ : écrire $\Delta = -(-\Delta) = i^2 \times (-\Delta)$, d'où $\sqrt{\Delta} = i\sqrt{-\Delta}$ dans $\mathbb{C}$. Il y a \cle{deux solutions complexes conjuguées} :
\[ z_1 = \frac{-b - i\sqrt{-\Delta}}{2a} \qquad \text{et} \qquad z_2 = \frac{-b + i\sqrt{-\Delta}}{2a} = \overline{z_1}. \]
\item Conclure par l'ensemble des solutions $S = \{z_1 ; z_2\}$, et vérifier éventuellement que $z_1 + z_2 = -\dfrac{b}{a}$ et $z_1 z_2 = \dfrac{c}{a}$.
\end{enumerate}
\end{methode}

\begin{exoresolu}[Équation $z^2 - 2z + 5 = 0$]
Résoudre dans $\mathbb{C}$ l'équation $z^2 - 2z + 5 = 0$, puis calculer le module de chaque solution.
\tcblower
\textbf{Étape 1 : discriminant.} Ici $a = 1$, $b = -2$, $c = 5$ :
\[ \Delta = b^2 - 4ac = (-2)^2 - 4 \times 1 \times 5 = 4 - 20 = -16. \]

\textbf{Étape 2 : cas $\Delta < 0$.} Comme $\Delta = -16 < 0$, l'équation admet deux solutions complexes conjuguées. On écrit $\sqrt{-\Delta} = \sqrt{16} = 4$, donc $\sqrt{\Delta} = 4i$ dans $\mathbb{C}$.

\textbf{Étape 3 : solutions.}
\[ z_1 = \frac{-b - i\sqrt{-\Delta}}{2a} = \frac{2 - 4i}{2} = 1 - 2i, \qquad z_2 = \frac{2 + 4i}{2} = 1 + 2i. \]

\textbf{Vérification par somme et produit :} $z_1 + z_2 = 2 = -\dfrac{b}{a}$ et $z_1 z_2 = (1-2i)(1+2i) = 1 + 4 = 5 = \dfrac{c}{a}$. Correct.

\textbf{Modules :} $|z_1| = |z_2| = \sqrt{1^2 + 2^2} = \sqrt{5}$ (deux complexes conjugués ont le même module).

\textbf{Conclusion :} $S = \{1 - 2i \,;\, 1 + 2i\}$, et chaque solution a pour module $\sqrt{5}$.
\end{exoresolu}

\courssub{Factoriser un polynôme de degré 3 connaissant une racine}

\begin{methode}[Factorisation d'un polynôme de degré 3]
Soit $P(z) = az^3 + bz^2 + cz + d$ admettant une racine connue $z_0$ (c'est-à-dire $P(z_0) = 0$).
\begin{enumerate}
\item \cle{Vérifier} d'abord que $P(z_0) = 0$ par un calcul explicite (indispensable dans la rédaction).
\item Écrire la factorisation $P(z) = (z - z_0)(az^2 + b'z + c')$.
\item Déterminer $b'$ et $c'$ par \cle{identification des coefficients} : développer $(z-z_0)(az^2+b'z+c')$ et identifier avec $P(z)$, ou poser la division euclidienne de $P$ par $z - z_0$.
\item Résoudre l'équation du second degré $az^2 + b'z + c' = 0$ dans $\mathbb{C}$ (discriminant).
\item Conclure : $P(z) = a(z - z_0)(z - z_1)(z - z_2)$ et donner l'ensemble des solutions de $P(z) = 0$.
\end{enumerate}
\end{methode}

\begin{exoresolu}[Factorisation de $P(z) = z^3 - 3z^2 + 4z - 2$]
On considère le polynôme $P(z) = z^3 - 3z^2 + 4z - 2$.
\begin{enumerate}
\item Vérifier que $z_0 = 1$ est une racine de $P$.
\item Déterminer les réels $b'$ et $c'$ tels que $P(z) = (z-1)(z^2 + b'z + c')$.
\item Résoudre dans $\mathbb{C}$ l'équation $P(z) = 0$.
\end{enumerate}
\tcblower
\textbf{1. Vérification.} $P(1) = 1 - 3 + 4 - 2 = 0$. Donc $z_0 = 1$ est bien une racine de $P$, et $P$ est factorisable par $(z-1)$.

\textbf{2. Identification des coefficients.} Développons :
\begin{align*}
(z-1)(z^2 + b'z + c') &= z^3 + b'z^2 + c'z - z^2 - b'z - c' \\
&= z^3 + (b'-1)z^2 + (c'-b')z - c'.
\end{align*}
Par identification avec $P(z) = z^3 - 3z^2 + 4z - 2$ :
\[ b' - 1 = -3, \qquad c' - b' = 4, \qquad -c' = -2. \]
La première équation donne $b' = -2$ ; la troisième donne $c' = 2$. Vérification avec la deuxième : $c' - b' = 2 - (-2) = 4$. Cohérent. Donc :
\[ P(z) = (z-1)(z^2 - 2z + 2). \]

\textbf{3. Résolution de $P(z) = 0$.} Un produit est nul si l'un de ses facteurs est nul :
\begin{itemize}
\item $z - 1 = 0 \iff z = 1$ ;
\item $z^2 - 2z + 2 = 0$ : discriminant $\Delta = (-2)^2 - 4 \times 1 \times 2 = 4 - 8 = -4 < 0$. Deux solutions complexes conjuguées :
\[ z = \frac{2 \pm i\sqrt{4}}{2} = \frac{2 \pm 2i}{2} = 1 \pm i. \]
\end{itemize}
\textbf{Conclusion :} $P(z) = (z-1)(z - 1 - i)(z - 1 + i)$ et l'ensemble des solutions est $S = \{1 \,;\, 1 - i \,;\, 1 + i\}$.
\end{exoresolu}

\courssub{Exploiter les propriétés du conjugué dans un problème}

\begin{methode}[Montrer qu'un complexe est réel ou imaginaire pur]
Pour déterminer la \cle{nature} d'un nombre complexe $z$ sans forcément calculer sa forme algébrique :
\begin{enumerate}
\item Pour montrer que $z$ est \cle{réel} : prouver que $\bar{z} = z$, ou que $\operatorname{Im}(z) = 0$.
\item Pour montrer que $z$ est \cle{imaginaire pur} : prouver que $\bar{z} = -z$, ou que $\operatorname{Re}(z) = 0$.
\item Utiliser les propriétés algébriques du conjugué : $\overline{z+z'} = \bar{z}+\bar{z'}$, $\overline{zz'} = \bar{z}\bar{z'}$, $\overline{z^n} = \bar{z}^n$.
\item Formules utiles : $\operatorname{Re}(z) = \dfrac{z+\bar{z}}{2}$ et $\operatorname{Im}(z) = \dfrac{z-\bar{z}}{2i}$.
\end{enumerate}
\end{methode}

\begin{exoresolu}[Nature de $Z = \dfrac{z-i}{z+i}$ pour $|z| = 1$]
Soit $z$ un nombre complexe de module $1$, différent de $-i$. On pose $Z = \dfrac{z - i}{z + i}$.
\begin{enumerate}
\item Montrer que $\bar{z} = \dfrac{1}{z}$.
\item En déduire que $Z$ est un nombre imaginaire pur.
\end{enumerate}
\tcblower
\textbf{1.} On sait que $z\bar{z} = |z|^2$. Comme $|z| = 1$, on a $z\bar{z} = 1$. De plus $|z| = 1$ entraîne $z \neq 0$, donc on peut diviser par $z$ : $\bar{z} = \dfrac{1}{z}$.

\textbf{2.} Pour montrer que $Z$ est imaginaire pur, montrons que $\overline{Z} = -Z$. En utilisant les propriétés du conjugué (quotient et somme), et le fait que $\overline{i} = -i$ :
\[ \overline{Z} = \overline{\left(\frac{z-i}{z+i}\right)} = \frac{\bar{z} - \overline{i}}{\bar{z} + \overline{i}} = \frac{\bar{z} + i}{\bar{z} - i}. \]
Remplaçons $\bar{z}$ par $\dfrac{1}{z}$ (question 1) :
\[ \overline{Z} = \frac{\dfrac{1}{z} + i}{\dfrac{1}{z} - i} = \frac{\dfrac{1 + iz}{z}}{\dfrac{1 - iz}{z}} = \frac{1 + iz}{1 - iz}. \]
Factorisons $i$ au numérateur et $-i$ au dénominateur : $1 + iz = i(z - i)$ car $i(z-i) = iz - i^2 = iz + 1$ ; et $1 - iz = -i(z + i)$ car $-i(z+i) = -iz - i^2 = -iz + 1$. Donc :
\[ \overline{Z} = \frac{i(z-i)}{-i(z+i)} = \frac{i}{-i} \times \frac{z-i}{z+i} = -\,\frac{z-i}{z+i} = -Z. \]
Comme $\overline{Z} = -Z$, le nombre $Z$ est imaginaire pur.

\textbf{Vérification sur un exemple :} pour $z = 1$ (qui est bien de module $1$), $Z = \dfrac{1-i}{1+i} = \dfrac{(1-i)^2}{(1+i)(1-i)} = \dfrac{1 - 2i + i^2}{2} = \dfrac{-2i}{2} = -i$, qui est bien imaginaire pur.

\textbf{Conclusion :} pour tout complexe $z$ de module $1$ distinct de $-i$, le nombre $Z = \dfrac{z-i}{z+i}$ est un imaginaire pur.
\end{exoresolu}

\begin{attention}[Les 5 erreurs qui coûtent des points au Bac]
\begin{enumerate}
\item \textbf{Oublier de rendre réel le dénominateur.} Écrire $z = \dfrac{2-3i}{1+2i}$ et s'arrêter là n'est \emph{pas} une forme algébrique. Il faut multiplier par le conjugué du dénominateur et aboutir à $a + ib$ avec $a, b$ réels explicites.
\item \textbf{Confondre $i^2$ et $-i^2$.} On a $i^2 = -1$, donc $-6i^2 = +6$ et non $-6$. De même $(2i)^2 = 4i^2 = -4$, et non $4$. Cette erreur de signe fausse tout le calcul qui suit.
\item \textbf{Écrire $\sqrt{-16} = -4$ ou « pas de solution ».} Dans $\mathbb{C}$, si $\Delta < 0$, l'équation du second degré à coefficients réels admet \emph{toujours} deux solutions complexes conjuguées, obtenues avec $\sqrt{\Delta} = i\sqrt{-\Delta}$. La racine carrée d'un réel négatif n'existe pas dans $\mathbb{R}$ : on écrit $i\sqrt{-\Delta}$, jamais $\sqrt{\Delta}$ seul.
\item \textbf{Oublier la condition sur les signes pour les racines carrées.} Après avoir trouvé $x^2$ et $y^2$, il reste quatre couples $(\pm x, \pm y)$ possibles, mais seuls deux conviennent : l'équation $2xy = b$ impose le signe relatif de $x$ et $y$. Toujours vérifier en développant le carré trouvé.
\item \textbf{Factoriser sans vérifier la racine.} Avant d'écrire $P(z) = (z - z_0)Q(z)$, il faut \emph{montrer} que $P(z_0) = 0$ par un calcul explicite. Et après identification des coefficients, vérifier la cohérence avec \emph{toutes} les équations obtenues, pas seulement les deux premières.
\end{enumerate}
\end{attention}

\newpage

\coursec{Exercices}

\courssub{Je vérifie mes connaissances}

\begin{exercice}[QCM : forme algébrique]
Pour chaque question, une seule réponse est exacte. Recopie la lettre correspondante sur ta copie.
\begin{enumerate}
\item Le nombre complexe $z = (1+i)^2$ est égal à :
\begin{enumerate}
\item[a)] $2$ \quad b) $2i$ \quad c) $1+i$ \quad d) $2+2i$
\end{enumerate}
\item Le conjugué de $z = -3+5i$ est :
\begin{enumerate}
\item[a)] $3-5i$ \quad b) $-3-5i$ \quad c) $3+5i$ \quad d) $5-3i$
\end{enumerate}
\item Le module de $z = 1-i\sqrt{3}$ est :
\begin{enumerate}
\item[a)] $2$ \quad b) $\sqrt{2}$ \quad c) $4$ \quad d) $1+\sqrt{3}$
\end{enumerate}
\item L'inverse de $i$ est :
\begin{enumerate}
\item[a)] $i$ \quad b) $1$ \quad c) $-i$ \quad d) $-1$
\end{enumerate}
\end{enumerate}
\end{exercice}

\begin{exercice}[Vrai ou faux ? Justifier]
Pour chacune des affirmations suivantes, dire si elle est vraie ou fausse, puis \textbf{justifier} la réponse par une démonstration courte ou un contre-exemple.
\begin{enumerate}
\item « Si $|z| = |z'|$ alors $z = z'$ ou $z = \bar{z}'$. »
\item « Tout nombre complexe $z$ de module $1$ vérifie $z = \dfrac{1}{\bar{z}}$. »
\item « Le produit de deux imaginaires purs est toujours un nombre réel. »
\item « Pour tout complexe $z$, on a $z + \bar{z} = 2\,\mathrm{Re}(z)$. »
\end{enumerate}
\end{exercice}

\begin{exercice}[Définitions et propriétés]
\begin{enumerate}
\item Donner la définition du \cle{conjugue} d'un nombre complexe $z = a+ib$ (avec $a, b \in \mathbb{R}$).
\item Donner la définition du \cle{module} de $z = a+ib$.
\item Compléter les égalités suivantes, valables pour tous complexes $z$ et $z'$ :
\[ \overline{z+z'} = \ldots \qquad ; \qquad \overline{z \times z'} = \ldots \qquad ; \qquad |z \times z'| = \ldots \qquad ; \qquad z\,\bar{z} = \ldots \]
\item À quelle condition nécessaire et suffisante un complexe $z$ est-il réel ? Imaginaire pur ? (Répondre à l'aide du conjugué.)
\end{enumerate}
\end{exercice}

\begin{exercice}[Calculs directs]
Mettre chacun des nombres complexes suivants sous forme algébrique :
\[ z_1 = (2-3i) + (5+i) \quad ; \quad z_2 = (1+2i)(3-i) \quad ; \quad z_3 = \dfrac{1+i}{1-i} \quad ; \quad z_4 = i^{2025} \quad ; \quad z_5 = (2-i)^2. \]
\end{exercice}

\courssub{Je m'entraîne}

\begin{exercice}[Les quatre opérations]
On donne les complexes $z_1 = 1+i\sqrt{3}$ et $z_2 = \sqrt{3}-i$.
\begin{enumerate}
\item Écrire sous forme algébrique les nombres : $z_1 + z_2$ ; $z_1 z_2$ ; $\dfrac{z_1}{z_2}$ ; $z_1^3$.
\item Calculer les modules : $|z_1|$, $|z_2|$, $|z_1 z_2|$ et $\left|\dfrac{z_1}{z_2}\right|$. Vérifier sur cet exemple les propriétés du module d'un produit et d'un quotient.
\end{enumerate}
\end{exercice}

\begin{exercice}[Racines carrées et équation]
\begin{enumerate}
\item Déterminer, sous forme algébrique, les racines carrées du nombre complexe $z = -7 + 24i$.
\item En déduire la résolution dans $\mathbb{C}$ de l'équation :
\[ z^2 - (3+2i)z + 3 - 3i = 0. \]
(\emph{Indication :} calculer le discriminant $\Delta$ et utiliser la question 1.)
\end{enumerate}
\end{exercice}

\begin{exercice}[Équations du second degré]
\begin{enumerate}
\item Résoudre dans $\mathbb{C}$ les équations suivantes :
\[ \text{a) } z^2 + 4z + 13 = 0 \qquad ; \qquad \text{b) } 2z^2 - 6z + 5 = 0. \]
\item Le plan complexe est muni d'un repère orthonormé direct $(O\,;\vec{u},\vec{v})$. Placer les images $A$, $B$, $C$, $D$ des solutions trouvées.
\item Que remarque-t-on concernant les solutions de chaque équation ? Expliquer ce phénomène par une propriété du cours.
\end{enumerate}
\end{exercice}

\begin{exercice}[Factorisation d'un polynôme de degré 3]
On considère le polynôme $P$ défini sur $\mathbb{C}$ par :
\[ P(z) = z^3 - (4+i)z^2 + (5+4i)z - 5i. \]
\begin{enumerate}
\item Vérifier que $i$ est une racine de $P$.
\item Déterminer des complexes $a$, $b$, $c$ tels que, pour tout $z \in \mathbb{C}$ :
\[ P(z) = (z - i)(az^2 + bz + c). \]
\item Résoudre dans $\mathbb{C}$ l'équation $P(z) = 0$.
\end{enumerate}
\end{exercice}

\courssub{Je me prépare au Bac}

\begin{exercice}[Type Bac : étude d'un quotient]
On considère le nombre complexe $z = \dfrac{2-4i}{1+3i}$.
\begin{enumerate}
\item Écrire $z$ sous forme algébrique.
\item Montrer que $z + \bar{z}$ est un nombre réel et que $z - \bar{z}$ est un imaginaire pur. Donner leurs valeurs.
\item Calculer $|z|$ de deux manières différentes.
\item On s'intéresse aux puissances de $z$.
\begin{enumerate}
\item Calculer $z^2$ et montrer que $z^2$ est un imaginaire pur. Calculer également $z^4$.
\item En remarquant que $z^{n+4} = z^4 \times z^n$, déterminer tous les entiers naturels $n \geq 1$ tels que $z^n$ soit un imaginaire pur. (\emph{Indication :} discuter selon le reste de la division euclidienne de $n$ par $4$.)
\end{enumerate}
\end{enumerate}
\end{exercice}

\begin{exercice}[Type Bac : problème de synthèse avec interprétation géométrique]
Le plan complexe est rapporté à un repère orthonormé direct $(O\,;\vec{u},\vec{v})$ d'unité graphique $\SI{1}{cm}$. Un technicien d'Orange Cameroun modélise la position d'un relais téléphonique par l'affixe $z$ d'un point $M$ du plan.
\begin{enumerate}
\item On cherche les nombres complexes $z$ vérifiant simultanément :
\[ z\,\bar{z} = 25 \qquad \text{et} \qquad z + \bar{z} = 6. \]
\begin{enumerate}
\item On pose $z = x + iy$ avec $x, y \in \mathbb{R}$. Traduire les deux conditions par un système d'équations d'inconnues $x$ et $y$.
\item Résoudre ce système et donner les solutions $z_1$ et $z_2$ sous forme algébrique.
\item Retrouver ces solutions en montrant que $z$ et $\bar{z}$ sont les racines d'une équation du second degré à coefficients réels que l'on déterminera.
\end{enumerate}
\item Interprétation géométrique.
\begin{enumerate}
\item Montrer que l'ensemble des points $M$ d'affixe $z$ telle que $z\,\bar{z} = 25$ est un cercle $\mathcal{C}$ dont on précisera le centre et le rayon.
\item Montrer que l'ensemble des points $M$ d'affixe $z$ telle que $z + \bar{z} = 6$ est une droite $\mathcal{D}$ dont on donnera une équation.
\item Construire $\mathcal{C}$ et $\mathcal{D}$, placer les points $M_1$ et $M_2$ d'affixes $z_1$ et $z_2$, et vérifier graphiquement le résultat de la question 1.
\end{enumerate}
\end{enumerate}
\end{exercice}

\begin{exercice}[Type Bac : polynôme, géométrie et application]
Une entreprise d'électrification rurale étudie un circuit dont l'impédance complexe $z$ (en ohms) est solution de l'équation :
\[ (E) : \quad z^3 - 4iz^2 - 6z + 4i = 0. \]
\begin{enumerate}
\item 
\begin{enumerate}
\item Vérifier que $z_0 = 2i$ est une solution de $(E)$.
\item Déterminer les complexes $a$, $b$ et $c$ tels que, pour tout $z \in \mathbb{C}$ :
\[ z^3 - 4iz^2 - 6z + 4i = (z - 2i)(az^2 + bz + c). \]
\item Achever la résolution de $(E)$ dans $\mathbb{C}$. On note $z_1$ et $z_2$ les deux autres solutions, avec $\mathrm{Re}(z_1) > 0$.
\end{enumerate}
\item Le plan complexe est muni d'un repère orthonormé direct $(O\,;\vec{u},\vec{v})$. On note $M_0$, $M_1$, $M_2$ les points d'affixes respectives $z_0$, $z_1$, $z_2$.
\begin{enumerate}
\item Calculer les modules de $z_0$, $z_1$ et $z_2$. Que peut-on en déduire pour les points $M_1$ et $M_2$ ?
\item Placer les trois points dans le repère (unité : $\SI{2}{cm}$).
\item Calculer $\dfrac{z_2 - z_0}{z_1 - z_0}$ sous forme algébrique, puis en déduire la nature du triangle $M_0 M_1 M_2$.
\end{enumerate}
\end{enumerate}
\end{exercice}



\newpage

\coursec{Corrigés des exercices}

\courssub{Je vérifie mes connaissances}

\begin{corrige}[Exercice 1 --- QCM : forme algébrique]
\textbf{Énoncé :}
\begin{enumerate}
\item Calculer $(1+i)^2$.
\item Donner le conjugué de $z = -3+5i$.
\item Calculer le module de $z = 1 - i\sqrt{3}$.
\item Simplifier $\dfrac{1}{i}$.
\end{enumerate}
\tcblower
\textbf{Solution détaillée :}
\begin{enumerate}
\item $(1+i)^2 = 1 + 2i + i^2 = 1 + 2i - 1 = 2i$. Réponse \textbf{b)}.
\item Le conjugué de $a+ib$ est $a-ib$ : on ne change que le signe de la partie imaginaire. $\bar{z} = -3-5i$. Réponse \textbf{b)}.
\item $|z| = \sqrt{1^2 + (-\sqrt{3})^2} = \sqrt{1+3} = \sqrt{4} = 2$. Réponse \textbf{a)}.
\item $\dfrac{1}{i} = \dfrac{1 \times (-i)}{i \times (-i)} = \dfrac{-i}{-i^2} = \dfrac{-i}{1} = -i$. Réponse \textbf{c)}.
\end{enumerate}
\end{corrige}

\begin{corrige}[Exercice 2 --- Vrai ou faux ? Justifier]
\textbf{Énoncé :}
\begin{enumerate}
\item Si $|z| = |z'|$, alors $z = z'$ ou $z = \bar{z}'$.
\item Si $|z| = 1$, alors $z = \dfrac{1}{\bar{z}}$.
\item Le produit de deux imaginaires purs est un réel.
\item Pour tout complexe $z$, $z + \bar{z} = 2\,\mathrm{Re}(z)$.
\end{enumerate}
\tcblower
\textbf{Solution détaillée :}
\begin{enumerate}
\item \textbf{Faux.} Contre-exemple : $z = 1$ et $z' = i$. On a $|z| = 1 = |z'|$, pourtant $z \neq z'$ et $\bar{z}' = -i \neq z$. L'égalité des modules ne dit rien sur l'égalité des nombres.
\item \textbf{Vrai.} Si $|z| = 1$, alors $z\,\bar{z} = |z|^2 = 1$, donc $z = \dfrac{1}{\bar{z}}$ (et $z \neq 0$ puisque $|z|=1$).
\item \textbf{Vrai.} Soient $z = ib$ et $z' = ib'$ avec $b, b' \in \mathbb{R}$. Alors $zz' = i^2 bb' = -bb' \in \mathbb{R}$.
\item \textbf{Vrai.} Si $z = a+ib$ ($a, b \in \mathbb{R}$), alors $z + \bar{z} = (a+ib) + (a-ib) = 2a = 2\,\mathrm{Re}(z)$.
\end{enumerate}
\end{corrige}

\begin{corrige}[Exercice 3 --- Définitions et propriétés]
\textbf{Énoncé :} Rappeler les définitions et propriétés fondamentales suivantes :
\begin{enumerate}
\item Conjugué d'un nombre complexe $z = a+ib$.
\item Module d'un nombre complexe $z = a+ib$.
\item Propriétés du conjugué et du module (somme, produit, quotient, $z\bar{z}$).
\item Caractérisation des réels et des imaginaires purs par le conjugué.
\end{enumerate}
\tcblower
\textbf{Solution détaillée :}
\begin{enumerate}
\item Le conjugué de $z = a+ib$ est $\bar{z} = a - ib$.
\item Le module de $z = a+ib$ est le réel positif $|z| = \sqrt{a^2+b^2}$.
\item $\overline{z+z'} = \bar{z} + \bar{z}'$ ; $\overline{z \times z'} = \bar{z} \times \bar{z}'$ ; $\overline{\left(\dfrac{z}{z'}\right)} = \dfrac{\bar{z}}{\bar{z}'}$ (pour $z'\neq 0$) ; $|z \times z'| = |z| \times |z'|$ ; $\left|\dfrac{z}{z'}\right| = \dfrac{|z|}{|z'|}$ ; $z\,\bar{z} = |z|^2$.
\item $z$ est réel si et seulement si $\bar{z} = z$ ; $z$ est imaginaire pur si et seulement si $\bar{z} = -z$.
\end{enumerate}
\end{corrige}

\begin{corrige}[Exercice 4 --- Calculs directs]
\textbf{Énoncé :} Mettre sous forme algébrique :
\begin{enumerate}
\item $z_1 = (2-3i) + (5+i)$
\item $z_2 = (3-i)(1+2i)$
\item $z_3 = \dfrac{1+i}{1-i}$
\item $z_4 = i^{2025}$
\item $z_5 = (2-i)^2$
\end{enumerate}
\tcblower
\textbf{Solution détaillée :}
\begin{itemize}
\item $z_1 = (2+5) + (-3+1)i = 7 - 2i$.
\item $z_2 = 3 + 6i - i - 2i^2 = 3 + 5i + 2 = 5 + 5i$.
\item $z_3 = \dfrac{(1+i)(1+i)}{(1-i)(1+i)} = \dfrac{1 + 2i + i^2}{1 - i^2} = \dfrac{2i}{2} = i$.
\item Les puissances de $i$ sont périodiques de période $4$ : $i^4 = 1$. Or $2025 = 4 \times 506 + 1$, donc $z_4 = i^{2025} = (i^4)^{506} \times i = i$.
\item $z_5 = 4 - 4i + i^2 = 3 - 4i$.
\end{itemize}
\end{corrige}

\courssub{Je m'entraîne}

\begin{corrige}[Exercice 5 --- Les quatre opérations]
\textbf{Énoncé :} Soient $z_1 = 1+i\sqrt{3}$ et $z_2 = \sqrt{3}-i$.
\begin{enumerate}
\item Calculer la somme, le produit, le quotient $\dfrac{z_1}{z_2}$ et le cube $z_1^3$.
\item Calculer les modules $|z_1|$, $|z_2|$, $|z_1 z_2|$, $\left|\dfrac{z_1}{z_2}\right|$ et vérifier les propriétés du module.
\end{enumerate}
\tcblower
\textbf{Solution détaillée :}
\begin{enumerate}
\item \textbf{Somme :} $z_1 + z_2 = (1+\sqrt{3}) + (\sqrt{3}-1)i$.

\textbf{Produit :}
\[ z_1 z_2 = (1+i\sqrt{3})(\sqrt{3}-i) = \sqrt{3} - i + 3i - i^2\sqrt{3} = \sqrt{3} + 2i + \sqrt{3} = 2\sqrt{3} + 2i. \]

\textbf{Quotient :} on multiplie numérateur et dénominateur par le conjugué du dénominateur, $\overline{z_2} = \sqrt{3}+i$ :
\[ \dfrac{z_1}{z_2} = \dfrac{(1+i\sqrt{3})(\sqrt{3}+i)}{(\sqrt{3}-i)(\sqrt{3}+i)} = \dfrac{\sqrt{3} + i + 3i + i^2\sqrt{3}}{3+1} = \dfrac{4i}{4} = i. \]

\textbf{Cube :} $z_1^2 = 1 + 2i\sqrt{3} + 3i^2 = -2 + 2i\sqrt{3}$, puis
\[ z_1^3 = z_1^2 \times z_1 = (-2+2i\sqrt{3})(1+i\sqrt{3}) = -2 - 2i\sqrt{3} + 2i\sqrt{3} + 6i^2 = -8. \]
\item $|z_1| = \sqrt{1+3} = 2$ ; $|z_2| = \sqrt{3+1} = 2$ ; $|z_1 z_2| = \sqrt{(2\sqrt{3})^2 + 2^2} = \sqrt{12+4} = 4$ ; $\left|\dfrac{z_1}{z_2}\right| = |i| = 1$.

Vérification : $|z_1 z_2| = 4 = |z_1| \times |z_2|$ et $\left|\dfrac{z_1}{z_2}\right| = 1 = \dfrac{|z_1|}{|z_2|}$. Les propriétés du module sont bien confirmées.
\end{enumerate}
\end{corrige}

\begin{corrige}[Exercice 6 --- Racines carrées et équation]
\textbf{Énoncé :}
\begin{enumerate}
\item Déterminer les racines carrées de $-7+24i$ sous forme algébrique.
\item Résoudre dans $\mathbb{C}$ l'équation $z^2 - (3+2i)z + (3-3i) = 0$.
\end{enumerate}
\tcblower
\textbf{Solution détaillée :}
\begin{enumerate}
\item On cherche $\delta = a+ib$ ($a, b \in \mathbb{R}$) tel que $\delta^2 = -7+24i$. On obtient le système :
\[ \begin{cases} a^2 - b^2 = -7 \\ 2ab = 24 \\ a^2 + b^2 = |-7+24i| = \sqrt{49+576} = 25 \end{cases} \]
En additionnant et soustrayant la première et la troisième équations : $2a^2 = 18$ donc $a^2 = 9$, et $2b^2 = 32$ donc $b^2 = 16$. Comme $2ab = 24 > 0$, $a$ et $b$ sont de même signe. Les racines carrées de $-7+24i$ sont donc :
\[ \delta_1 = 3+4i \qquad \text{et} \qquad \delta_2 = -3-4i. \]
\item Le discriminant vaut :
\[ \Delta = (3+2i)^2 - 4(3-3i) = 9 + 12i + 4i^2 - 12 + 12i = -7 + 24i. \]
D'après la question 1, une racine carrée de $\Delta$ est $\delta = 3+4i$. Les solutions sont :
\[ z_1 = \dfrac{(3+2i) + (3+4i)}{2} = \dfrac{6+6i}{2} = 3+3i \qquad ; \qquad z_2 = \dfrac{(3+2i) - (3+4i)}{2} = \dfrac{-2i}{2} = -i. \]
\[ \mathcal{S} = \{3+3i \; ; \; -i\}. \]
\end{enumerate}
\end{corrige}

\begin{corrige}[Exercice 7 --- Équations du second degré]
\textbf{Énoncé :}
\begin{enumerate}
\item Résoudre dans $\mathbb{C}$ :
\begin{itemize}
\item[a)] $z^2 + 4z + 13 = 0$
\item[b)] $2z^2 - 6z + 5 = 0$
\end{itemize}
\item Représenter les solutions dans le plan complexe muni d'un repère orthonormé $(\vec{u}, \vec{v})$.
\item Quelle relation lie les deux solutions de chaque équation ? Généraliser.
\end{enumerate}
\tcblower
\textbf{Solution détaillée :}
\begin{enumerate}
\item \textbf{a)} $\Delta = 16 - 52 = -36 = (6i)^2$. Les solutions sont :
\[ z_1 = \dfrac{-4+6i}{2} = -2+3i \qquad ; \qquad z_2 = \dfrac{-4-6i}{2} = -2-3i. \]
\textbf{b)} $\Delta = 36 - 40 = -4 = (2i)^2$. Les solutions sont :
\[ z_3 = \dfrac{6+2i}{4} = \dfrac{3+i}{2} \qquad ; \qquad z_4 = \dfrac{6-2i}{4} = \dfrac{3-i}{2}. \]
\item 
\begin{popfigure}
\begin{center}
\begin{tikzpicture}[scale=1.1]
\draw[->,popline] (-3.2,0) -- (3.2,0) node[below left] {$\vec{u}$};
\draw[->,popline] (0,-3.5) -- (0,3.5) node[below left] {$\vec{v}$};
\draw[popline,dashed] (-2,-3) -- (-2,3);
\draw[popline,dashed] (1.5,-0.5) -- (1.5,0.5);
\fill[popblue] (-2,3) circle (2pt) node[above left] {$A$};
\fill[popblue] (-2,-3) circle (2pt) node[below left] {$B$};
\fill[poporange] (1.5,0.5) circle (2pt) node[above right] {$C$};
\fill[poporange] (1.5,-0.5) circle (2pt) node[below right] {$D$};
\fill[popdark] (0,0) circle (1.5pt) node[below left] {$O$};
\end{tikzpicture}
\end{center}
\legende{Images des solutions : $A$ et $B$ d'une part, $C$ et $D$ d'autre part, sont symétriques par rapport à l'axe des réels.}
\end{popfigure}
\item Pour chaque équation, les deux solutions sont \textbf{conjuguées l'une de l'autre} : $z_2 = \overline{z_1}$ et $z_4 = \overline{z_3}$. C'est une propriété générale : si une équation du second degré à \emph{coefficients réels} admet un discriminant strictement négatif, ses deux solutions dans $\mathbb{C}$ sont deux complexes conjugués, $\dfrac{-b + i\sqrt{-\Delta}}{2a}$ et $\dfrac{-b - i\sqrt{-\Delta}}{2a}$.
\end{enumerate}
\end{corrige}

\begin{corrige}[Exercice 8 --- Factorisation d'un polynôme de degré 3]
\textbf{Énoncé :} Soit le polynôme $P(z) = z^3 - (4+i)z^2 + (5+4i)z - 5i$.
\begin{enumerate}
\item Vérifier que $i$ est une racine de $P$.
\item Factoriser $P(z)$ sous la forme $(z-i)Q(z)$ où $Q$ est un trinôme du second degré.
\item Résoudre l'équation $P(z) = 0$ dans $\mathbb{C}$.
\end{enumerate}
\tcblower
\textbf{Solution détaillée :}
\begin{enumerate}
\item $P(i) = i^3 - (4+i)i^2 + (5+4i)i - 5i = -i + (4+i) + (5i - 4) - 5i = 0$. Donc $i$ est bien une racine de $P$.
\item On développe : $(z-i)(az^2+bz+c) = az^3 + (b-ia)z^2 + (c-ib)z - ic$. Par identification avec $P(z)$ :
\[ a = 1 \quad ; \quad b - ia = -(4+i) \quad ; \quad c - ib = 5+4i \quad ; \quad -ic = -5i. \]
On tire $a = 1$, $b = -(4+i) + i = -4$, $c = 5$. Vérification : $c - ib = 5 + 4i$ : correct. Donc :
\[ P(z) = (z-i)(z^2 - 4z + 5). \]
\item $P(z) = 0 \iff z = i$ ou $z^2 - 4z + 5 = 0$. Pour le trinôme : $\Delta = 16 - 20 = -4 = (2i)^2$, d'où $z = \dfrac{4 \pm 2i}{2} = 2 \pm i$.
\[ \mathcal{S} = \{i \; ; \; 2+i \; ; \; 2-i\}. \]
\end{enumerate}
\end{corrige}

\courssub{Je me prépare au Bac}

\begin{corrige}[Exercice 9 --- Type Bac : étude d'un quotient]
\textbf{Barème indicatif (sur 4 points) :} question 1 : 0,75 pt ; question 2 : 0,5 pt ; question 3 : 0,75 pt ; question 4.a : 0,75 pt ; question 4.b : 1,25 pt.

\textbf{Énoncé :} Soit $z = \dfrac{2-4i}{1+3i}$.
\begin{enumerate}
\item Mettre $z$ sous forme algébrique.
\item Vérifier que $z + \bar{z} \in \mathbb{R}$ et $z - \bar{z} \in i\mathbb{R}$.
\item Calculer $|z|$ de deux façons différentes.
\item \begin{enumerate}
\item Calculer $z^2$ et $z^4$. Quelle est leur nature ?
\item Déterminer les entiers naturels $n$ non nuls tels que $z^n$ soit un imaginaire pur.
\end{enumerate}
\end{enumerate}
\tcblower
\textbf{Solution détaillée :}
\begin{enumerate}
\item On multiplie par le conjugué du dénominateur :
\[ z = \dfrac{(2-4i)(1-3i)}{(1+3i)(1-3i)} = \dfrac{2 - 6i - 4i + 12i^2}{1+9} = \dfrac{-10 - 10i}{10} = -1 - i. \]
\item $z + \bar{z} = (-1-i) + (-1+i) = -2 \in \mathbb{R}$ ; $z - \bar{z} = (-1-i) - (-1+i) = -2i$, imaginaire pur. (C'est conforme au cours : $z+\bar{z} = 2\,\mathrm{Re}(z)$ et $z - \bar{z} = 2i\,\mathrm{Im}(z)$.)
\item \textbf{Première méthode} (forme algébrique) : $|z| = \sqrt{(-1)^2 + (-1)^2} = \sqrt{2}$.

\textbf{Deuxième méthode} (module d'un quotient) : $|z| = \dfrac{|2-4i|}{|1+3i|} = \dfrac{\sqrt{4+16}}{\sqrt{1+9}} = \dfrac{\sqrt{20}}{\sqrt{10}} = \sqrt{2}$.
\item 
\begin{enumerate}
\item $z^2 = (-1-i)^2 = 1 + 2i + i^2 = 2i$ : c'est bien un imaginaire pur. Puis $z^4 = (z^2)^2 = (2i)^2 = -4$, qui est réel.
\item Comme $z^4 = -4$, on a $z^{n+4} = -4\,z^n$ : multiplier par $-4$ (réel non nul) ne change ni le caractère réel ni le caractère imaginaire pur d'un nombre. Il suffit donc d'examiner les quatre premiers cas :
\begin{itemize}
\item $z^1 = -1-i$ : ni réel, ni imaginaire pur ;
\item $z^2 = 2i$ : imaginaire pur ;
\item $z^3 = z^2 \times z = 2i(-1-i) = 2 - 2i$ : ni réel, ni imaginaire pur ;
\item $z^4 = -4$ : réel.
\end{itemize}
En écrivant la division euclidienne $n = 4k + r$ avec $r \in \{1, 2, 3, 4\}$ (et $k \geq 0$), on a $z^n = (-4)^k z^r$. Le nombre $z^n$ est imaginaire pur si et seulement si $z^r$ l'est, c'est-à-dire si et seulement si $r = 2$.
\[ z^n \text{ est imaginaire pur} \iff n \equiv 2 \pmod{4} \iff n \in \{2, 6, 10, 14, \ldots\}. \]
\end{enumerate}
\end{enumerate}

\textbf{Points de vigilance :}
\begin{itemize}
\item Pour le quotient, il faut \emph{explicitement} écrire la multiplication par le conjugué du dénominateur : c'est la technique attendue et elle doit apparaître sur la copie.
\item Dans la question 4.b, il faut justifier que $(-4)^k$ est un réel non nul, donc que le caractère « imaginaire pur » de $z^n$ dépend uniquement de celui de $z^r$. Oublier cette justification fait perdre le point de raisonnement.
\item Attention : un imaginaire pur peut avoir une partie réelle nulle mais une partie imaginaire \emph{non nulle} ; $z^4 = -4$ est réel, pas imaginaire pur.
\end{itemize}
\end{corrige}

\begin{corrige}[Exercice 10 --- Type Bac : problème de synthèse avec interprétation géométrique]
\textbf{Barème indicatif (sur 5 points) :} question 1.a : 0,5 pt ; 1.b : 0,75 pt ; 1.c : 1 pt ; 2.a : 0,75 pt ; 2.b : 0,5 pt ; 2.c (figure + vérification) : 1,5 pt.

\textbf{Énoncé :} Soit $z \in \mathbb{C}$ tel que $\begin{cases} z\,\bar{z} = 25 \\ z + \bar{z} = 6 \end{cases}$.
\begin{enumerate}
\item \begin{enumerate}
\item Écrire ce système en fonction de $x = \mathrm{Re}(z)$ et $y = \mathrm{Im}(z)$.
\item En déduire les valeurs possibles de $z$.
\item Retrouver ces valeurs en utilisant la relation entre les racines d'un trinôme et ses coefficients (somme et produit).
\end{enumerate}
\item \begin{enumerate}
\item Déterminer l'ensemble des points $M$ d'affixe $z$ tels que $z\,\bar{z} = 25$.
\item Déterminer l'ensemble des points $M$ d'affixe $z$ tels que $z + \bar{z} = 6$.
\item Tracer ces deux ensembles dans un même repère (unité graphique $\SI{1}{cm}$) et vérifier graphiquement les solutions de la question 1.b.
\end{enumerate}
\end{enumerate}
\tcblower
\textbf{Solution détaillée :}
\begin{enumerate}
\item 
\begin{enumerate}
\item Avec $z = x+iy$ : $z\,\bar{z} = x^2 + y^2$ et $z + \bar{z} = 2x$. Le système s'écrit :
\[ \begin{cases} x^2 + y^2 = 25 \\ 2x = 6 \end{cases} \]
\item La deuxième équation donne $x = 3$ ; en reportant : $9 + y^2 = 25$, donc $y^2 = 16$, soit $y = 4$ ou $y = -4$. Les solutions sont :
\[ z_1 = 3 + 4i \qquad \text{et} \qquad z_2 = 3 - 4i. \]
\item Posons $S = z + \bar{z} = 6$ et $P = z\,\bar{z} = 25$. Deux nombres dont on connaît la somme $S$ et le produit $P$ sont les racines de l'équation $X^2 - SX + P = 0$, c'est-à-dire :
\[ X^2 - 6X + 25 = 0. \]
$\Delta = 36 - 100 = -64 = (8i)^2$, d'où $X = \dfrac{6 \pm 8i}{2} = 3 \pm 4i$. On retrouve bien $z_1 = 3+4i$ et $z_2 = 3-4i$.
\end{enumerate}
\item 
\begin{enumerate}
\item $z\,\bar{z} = 25 \iff |z|^2 = 25 \iff |z| = 5 \iff OM = 5$. L'ensemble cherché est le cercle $\mathcal{C}$ de centre $O$ et de rayon $5$.
\item $z + \bar{z} = 6 \iff 2x = 6 \iff x = 3$. L'ensemble cherché est la droite $\mathcal{D}$ d'équation $x = 3$ (parallèle à l'axe des imaginaires purs).
\item 
\begin{popfigure}
\begin{center}
\begin{tikzpicture}[scale=0.62]
\draw[->,popline] (-6,0) -- (6,0);
\draw[->,popline] (0,-6) -- (0,6);
\draw[popblue,thick] (0,0) circle (5);
\draw[poporange,thick] (3,-5.5) -- (3,5.5);
\fill[popdark] (3,4) circle (2.5pt) node[above right] {$M_1(3+4i)$};
\fill[popdark] (3,-4) circle (2.5pt) node[below right] {$M_2(3-4i)$};
\fill[popdark] (0,0) circle (2pt) node[below left] {$O$};
\node[popblue] at (-3.9,3.9) {$\mathcal{C}$};
\node[poporange] at (3,5.9) {$\mathcal{D}$};
\end{tikzpicture}
\end{center}
\legende{Le cercle $\mathcal{C}$ ($|z|=5$) et la droite $\mathcal{D}$ ($x=3$) se coupent exactement en $M_1$ et $M_2$.}
\end{popfigure}
Graphiquement, $\mathcal{C}$ et $\mathcal{D}$ se coupent en deux points d'abscisse $3$ et d'ordonnées $4$ et $-4$ : ce sont bien $M_1$ et $M_2$, ce qui confirme le calcul.
\end{enumerate}
\end{enumerate}

\textbf{Points de vigilance :}
\begin{itemize}
\item La chaîne d'équivalences $z\bar{z} = 25 \iff |z| = 5 \iff OM = 5$ doit être rédigée : c'est elle qui justifie que l'ensemble est un cercle, et non une simple affirmation.
\item Pour la question 1.c, citer explicitement la propriété « somme $S$ et produit $P$ : racines de $X^2 - SX + P = 0$ ».
\item Sur la figure, respecter l'unité graphique de $\SI{1}{cm}$ : le cercle a un rayon de $\SI{5}{cm}$, la droite coupe l'axe des réels en $x = 3$.
\end{itemize}
\end{corrige}

\begin{corrige}[Exercice 11 --- Type Bac : polynôme, géométrie et application]
\textbf{Barème indicatif (sur 5 points) :} question 1.a : 0,5 pt ; 1.b : 1 pt ; 1.c : 1 pt ; 2.a : 0,75 pt ; 2.b (figure) : 0,75 pt ; 2.c : 1 pt.

\textbf{Énoncé :} On considère l'équation $(E) : z^3 - 4iz^2 - 6z + 4i = 0$.
\begin{enumerate}
\item \begin{enumerate}
\item Vérifier que $z_0 = 2i$ est solution de $(E)$.
\item Factoriser le polynôme $P(z) = z^3 - 4iz^2 - 6z + 4i$ par $(z-2i)$.
\item Résoudre complètement $(E)$ dans $\mathbb{C}$. On note $z_1$ la solution de partie réelle strictement positive et $z_2$ l'autre.
\end{enumerate}
\item \begin{enumerate}
\item Calculer les modules $|z_0|$, $|z_1|$, $|z_2|$. Que peut-on dire des points $M_1$ et $M_2$ d'affixes respectives $z_1$ et $z_2$ par rapport au cercle de centre $O$ et de rayon $\sqrt{2}$ ?
\item Représenter dans un repère orthonormé (unité $\SI{2}{cm}$) les points $M_0$, $M_1$, $M_2$ et le cercle de centre $O$ passant par $M_1$ et $M_2$.
\item Étudier la nature du triangle $M_0M_1M_2$.
\end{enumerate}
\end{enumerate}
\tcblower
\textbf{Solution détaillée :}
\begin{enumerate}
\item 
\begin{enumerate}
\item Calculons : $(2i)^3 = 8i^3 = -8i$ ; $-4i(2i)^2 = -4i \times (-4) = 16i$ ; $-6 \times 2i = -12i$. Donc :
\[ -8i + 16i - 12i + 4i = 0. \]
$z_0 = 2i$ est bien solution de $(E)$.
\item On développe : $(z-2i)(az^2+bz+c) = az^3 + (b-2ia)z^2 + (c-2ib)z - 2ic$. Par identification :
\[ a = 1 \quad ; \quad b - 2ia = -4i \quad ; \quad c - 2ib = -6 \quad ; \quad -2ic = 4i. \]
On obtient $a = 1$, $b = -4i + 2i = -2i$, $c = -2$. Vérification : $c - 2ib = -2 - 2i(-2i) = -2 + 4i^2 = -6$ : correct. Donc :
\[ z^3 - 4iz^2 - 6z + 4i = (z - 2i)(z^2 - 2iz - 2). \]
\item Il reste à résoudre $z^2 - 2iz - 2 = 0$ : $\Delta = (-2i)^2 + 8 = -4 + 8 = 4 = 2^2$. D'où :
\[ z = \dfrac{2i \pm 2}{2} = \pm 1 + i. \]
Avec $\mathrm{Re}(z_1) > 0$ : $z_1 = 1 + i$ et $z_2 = -1 + i$.
\[ \mathcal{S} = \{2i \; ; \; 1+i \; ; \; -1+i\}. \]
\end{enumerate}
\item 
\begin{enumerate}
\item $|z_0| = |2i| = 2$ ; $|z_1| = \sqrt{1+1} = \sqrt{2}$ ; $|z_2| = \sqrt{1+1} = \sqrt{2}$. Comme $|z_1| = |z_2| = \sqrt{2}$, les points $M_1$ et $M_2$ appartiennent au cercle de centre $O$ et de rayon $\sqrt{2}$ (ils sont d'ailleurs symétriques par rapport à l'axe des imaginaires purs, car $z_2 = -\overline{z_1}$).
\item 
\begin{popfigure}
\begin{center}
\begin{tikzpicture}[scale=1.6]
\draw[->,popline] (-1.8,0) -- (1.8,0);
\draw[->,popline] (0,-0.6) -- (0,2.6);
\draw[popblue,thick,dashed] (0,0) circle (1.414);
\fill[poporange] (0,2) circle (1.8pt) node[above right] {$M_0(2i)$};
\fill[popblue] (1,1) circle (1.8pt) node[right] {$M_1(1+i)$};
\fill[popblue] (-1,1) circle (1.8pt) node[left] {$M_2(-1+i)$};
\draw[popdark] (0,2) -- (1,1) -- (-1,1) -- cycle;
\fill[popdark] (0,0) circle (1.2pt) node[below left] {$O$};
\end{tikzpicture}
\end{center}
\legende{Les points $M_0$, $M_1$, $M_2$ et le cercle de centre $O$ passant par $M_1$ et $M_2$.}
\end{popfigure}
\item Calculons :
\[ \dfrac{z_2 - z_0}{z_1 - z_0} = \dfrac{-1+i-2i}{1+i-2i} = \dfrac{-1-i}{1-i} = \dfrac{(-1-i)(1+i)}{(1-i)(1+i)} = \dfrac{-1 - i - i - i^2}{2} = \dfrac{-2i}{2} = -i. \]
On en tire $\left|\dfrac{z_2 - z_0}{z_1 - z_0}\right| = |-i| = 1$, c'est-à-dire $\dfrac{M_0M_2}{M_0M_1} = 1$ : le triangle est \textbf{isocèle en $M_0$}.

De plus, $M_0M_1 = |z_1 - z_0| = |1 - i| = \sqrt{2}$, donc $M_0M_2 = \sqrt{2}$ aussi, et $M_1M_2 = |z_1 - z_2| = |2| = 2$. Or :
\[ M_0M_1^2 + M_0M_2^2 = 2 + 2 = 4 = M_1M_2^2. \]
Par la réciproque du théorème de Pythagore, le triangle $M_0M_1M_2$ est \textbf{rectangle en $M_0$}. Conclusion : $M_0M_1M_2$ est un triangle \cle{rectangle isocèle} en $M_0$.
\end{enumerate}
\end{enumerate}

\textbf{Points de vigilance :}
\begin{itemize}
\item Pour l'identification des coefficients, la \emph{vérification} des équations non utilisées (ici $c - 2ib = -6$) est une exigence de rigueur appréciée du correcteur.
\item Attention au calcul du discriminant : $\Delta = b^2 - 4ac$ avec $b = -2i$, donc $b^2 = (-2i)^2 = -4$ et non $4$.
\item Pour conclure « rectangle », il faut citer la \emph{réciproque} du théorème de Pythagore ; le seul calcul du quotient $-i$ ne suffit pas à ce niveau du cours (l'interprétation angulaire sera vue avec la forme trigonométrique).
\item Sur la figure, l'unité est $\SI{2}{cm}$ : $M_0$ est donc à $\SI{4}{cm}$ de $O$ sur l'axe des imaginaires.
\end{itemize}
\end{corrige}

\newpage

\coursec{Fiche bilan}

\courssub{Carte mentale de la leçon}

\begin{popfigure}
\begin{tikzpicture}[mindmap, grow cyclic, every node/.style={concept, text=white, font=\small},
  root concept/.append style={concept color=popdark, font=\bfseries},
  level 1/.append style={level distance=3.6cm, sibling angle=60},
  level 2/.append style={level distance=2.2cm, sibling angle=45, font=\scriptsize}]
\node[root concept, align=center]{Nombres complexes\\ forme algébrique}
  child[concept color=popblue]{ node{L'ensemble $\mathbb{C}$}
    child{ node{$i^2=-1$} }
    child{ node{$z=a+ib$} }
  }
  child[concept color=poporange]{ node{Opérations}
    child{ node{Somme, produit} }
    child{ node{Quotient : multiplier par $\bar{z}$} }
  }
  child[concept color=popgreen]{ node{Conjugué}
    child{ node{$\bar{z}=a-ib$} }
    child{ node{$z\bar{z}=|z|^2$} }
  }
  child[concept color=popgold]{ node{Module}
    child{ node{$|z|=\sqrt{a^2+b^2}$} }
    child{ node{$|zz'|=|z||z'|$} }
  }
  child[concept color=poppurple]{ node{Racines carrées}
    child{ node{Système $x^2-y^2$, $2xy$, $x^2+y^2$} }
  }
  child[concept color=poppink]{ node{Équations}
    child{ node{$\Delta<0$ : 2 racines conjuguées} }
    child{ node{Factorisation degré 3} }
  };
\end{tikzpicture}
\legende{Carte mentale : les six blocs de compétences de la leçon sur les nombres complexes.}
\end{popfigure}

\courssub{L'essentiel à retenir}

\begin{aretenir}[Synthèse : Définitions et Formules fondamentales]
\textbf{Définitions clés}
\begin{itemize}
  \item \cle{Unité imaginaire} : nombre noté $i$ tel que $i^2=-1$.
  \item \cle{Forme algébrique} : tout complexe s'écrit de manière unique $z=a+ib$ avec $(a,b)\in\mathbb{R}^2$ ; $a=\mathrm{Re}(z)$ est la partie réelle, $b=\mathrm{Im}(z)$ la partie imaginaire.
  \item $z$ est \cle{réel} $\iff b=0$ ; $z$ est \cle{imaginaire pur} $\iff a=0$.
  \item Égalité : $a+ib=a'+ib' \iff a=a'$ et $b=b'$.
  \item \cle{Conjugué} : $\bar{z}=a-ib$.
  \item \cle{Module} : $|z|=\sqrt{a^2+b^2}$ (distance de l'image de $z$ à l'origine dans le plan complexe).
\end{itemize}

\textbf{Formules à connaître par cœur}
\begin{center}
\begin{tabularx}{\linewidth}{|l|X|}
\hline
\rowcolor{popblueL} \textbf{Notion} & \textbf{Formules} \\
\hline
Opérations & $(a+ib)+(a'+ib')=(a+a')+i(b+b')$ ; $(a+ib)(a'+ib')=(aa'-bb')+i(ab'+a'b)$ \\
\hline
Quotient ($z'\neq 0$) & $\dfrac{z}{z'}=\dfrac{z\,\bar{z'}}{|z'|^2}$ \quad (multiplier numérateur et dénominateur par le conjugué du dénominateur) \\
\hline
Conjugué & $\overline{z+z'}=\bar{z}+\bar{z'}$ ; $\overline{zz'}=\bar{z}\,\bar{z'}$ ; $\overline{\left(\dfrac{z}{z'}\right)}=\dfrac{\bar{z}}{\bar{z'}}$ ; $\overline{z^n}=\bar{z}^{\,n}$ \\
\hline
Réel / Imaginaire pur & $z$ réel $\iff z=\bar{z}$ ; $z$ imaginaire pur $\iff z=-\bar{z}$ ; $\mathrm{Re}(z)=\dfrac{z+\bar{z}}{2}$ ; $\mathrm{Im}(z)=\dfrac{z-\bar{z}}{2i}$ \\
\hline
Module & $z\bar{z}=|z|^2$ ; $|z|=|\bar{z}|=|-z|$ ; $|zz'|=|z||z'|$ ; $\left|\dfrac{z}{z'}\right|=\dfrac{|z|}{|z'|}$ ; $|z^n|=|z|^n$ \\
\hline
Inégalité triangulaire & $|z+z'|\le |z|+|z'|$ \\
\hline
Second degré ($\Delta<0$) & $z_1=\dfrac{-b-i\sqrt{-\Delta}}{2a}$ et $z_2=\dfrac{-b+i\sqrt{-\Delta}}{2a}=\bar{z_1}$ \\
\hline
\end{tabularx}
\end{center}
\end{aretenir}

\begin{methode}[Méthodes express]
\begin{enumerate}
  \item \cle{Forme algébrique d'un quotient} : multiplier le numérateur et le dénominateur par le conjugué du dénominateur, développer en utilisant $i^2=-1$, regrouper la partie réelle et la partie imaginaire pour obtenir la forme $a+ib$ avec $a,b\in\mathbb{R}$.
  \item \cle{Racines carrées de $z=a+ib$} : poser $\delta=x+iy$ ($x,y\in\mathbb{R}$) et résoudre le système
  \[
  \begin{cases}
  x^2-y^2=a \\
  2xy=b \\
  x^2+y^2=\sqrt{a^2+b^2}=|z|
  \end{cases}
  \]
  On en déduit $x^2=\dfrac{a+|z|}{2}$ et $y^2=\dfrac{|z|-a}{2}$. Les signes de $x$ et $y$ sont choisis pour que $2xy$ ait le signe de $b$. Les deux racines sont opposées : $\delta$ et $-\delta$.
  \item \cle{Équation $az^2+bz+c=0$ à coefficients réels} : calculer $\Delta=b^2-4ac$.
  \begin{itemize}
    \item Si $\Delta>0$ : deux solutions réelles distinctes.
    \item Si $\Delta=0$ : une solution réelle double.
    \item Si $\Delta<0$ : deux solutions complexes conjuguées $z=\dfrac{-b\pm i\sqrt{-\Delta}}{2a}$.
  \end{itemize}
  \item \cle{Factorisation d'un polynôme de degré 3} connaissant une racine $\alpha$ : écrire $P(z)=(z-\alpha)(z^2+bz+c)$, déterminer $b$ et $c$ par identification des coefficients (ou division euclidienne), puis résoudre l'équation du second degré $z^2+bz+c=0$ dans $\mathbb{C}$.
\end{enumerate}
\end{methode}

\begin{aretenir}[Checklist avant le Bac]
Avant l'épreuve, vérifie que tu sais :
\begin{itemize}
  \item[$\square$] écrire $i^2=-1$ et simplifier les puissances de $i$ ($i^3=-i$, $i^4=1$, $i^{4k+r}=i^r$) ;
  \item[$\square$] additionner, soustraire et multiplier deux complexes sous forme algébrique ;
  \item[$\square$] mettre un quotient sous forme algébrique $a+ib$ en multipliant par le conjugué du dénominateur ;
  \item[$\square$] donner le conjugué d'un complexe et utiliser la relation fondamentale $z\bar{z}=|z|^2$ ;
  \item[$\square$] calculer le module $|z|=\sqrt{a^2+b^2}$ et appliquer les propriétés $|zz'|=|z||z'|$, $\left|\frac{z}{z'}\right|=\frac{|z|}{|z'|}$ ;
  \item[$\square$] montrer qu'un complexe est réel ($z=\bar{z}$) ou imaginaire pur ($z=-\bar{z}$) ;
  \item[$\square$] déterminer les racines carrées d'un complexe par le système en $x$ et $y$ ;
  \item[$\square$] vérifier que les deux racines carrées trouvées sont bien opposées ;
  \item[$\square$] résoudre $az^2+bz+c=0$ avec $\Delta<0$ et écrire les deux solutions conjuguées sous la forme $\frac{-b}{2a}\pm i\frac{\sqrt{-\Delta}}{2a}$ ;
  \item[$\square$] factoriser un polynôme de degré 3 connaissant une racine, par identification ou division euclidienne ;
  \item[$\square$] résoudre complètement l'équation $P(z)=0$ associée et donner l'ensemble des solutions dans $\mathbb{C}$ ;
  \item[$\square$] rédiger proprement : forme finale $a+ib$ avec $a$ et $b$ réels explicites, sans $i$ au dénominateur, pas de racine au dénominateur si possible.
\end{itemize}
\end{aretenir}

\findecours

\end{document}
